%% marginfx --- supplementary material, formatted for JRSS-B.
%% Companion file: jrssb-main.tex (cross-referenced via xr-hyper).
\documentclass[12pt,a4paper]{article}

\usepackage[T1]{fontenc}
\usepackage{amsthm,amsmath,amsfonts,amssymb}
\usepackage{mathtools}
\usepackage[authoryear,round]{natbib}
\usepackage{graphicx}
\usepackage{booktabs,multirow,threeparttable,tabularx}
\usepackage[figuresright]{rotating}
\usepackage{enumerate}
\usepackage[a4paper,margin=1in]{geometry}
\usepackage{setspace}

\usepackage{xr-hyper}
\usepackage[colorlinks,citecolor=blue,urlcolor=blue,linkcolor=blue]{hyperref}
\externaldocument[M-]{jrssb-main}

\numberwithin{equation}{section}

%% ---- Theorem environments ----
%% JRSS-B asks that results be numbered sequentially and individually by type,
%% not by section, which is what these plain \newtheorem counters give.
\theoremstyle{plain}
\newtheorem{theorem}{Theorem}
\newtheorem{lemma}{Lemma}
\newtheorem{corollary}{Corollary}
\newtheorem{proposition}{Proposition}
\newtheorem{conjecture}{Conjecture}

\theoremstyle{definition}
\newtheorem{definition}{Definition}
\newtheorem{assumption}{Assumption}
\newtheorem{remark}{Remark}

%% ---- Author macros ----
\newcommand{\norm}[1]{\left\| #1 \right\|}
\newcommand{\R}{\mathbb{R}}
\newcommand{\E}{\mathbb{E}}
\newcommand{\Lp}{L^p(\Omega)}
\newcommand{\Lone}{L^1(\Omega)}
\newcommand{\Wonep}{W^{1,p}(\Omega)}
\newcommand{\Wtwop}{W^{2,p}(\Omega)}
\newcommand{\BV}{BV(\Omega)}
\newcommand{\cM}{\mathcal{M}(\Omega)^d}
\newcommand{\fn}{\hat{f}_n}
\newcommand{\xrightarrowp}{\xrightarrow{p}}
\newcommand{\wstar}{\xrightarrow{w*}}
\newcommand{\Prob}{\mathbb{P}}
\newcommand{\abs}[1]{\left| #1 \right|}
\newcommand{\Dh}{\Delta_{j,h}}
\newcommand{\Ptwo}[1]{\norm{#1}_{P,2}}
\newcommand{\Pq}[2]{\norm{#1}_{P,#2}}

%% Everything in this document carries an S prefix so that a reference
%% from the main paper is unambiguous.
\renewcommand{\thesection}{S\arabic{section}}
\renewcommand{\thetable}{S\arabic{table}}
\renewcommand{\thefigure}{S\arabic{figure}}

%% The stated results moved here from Section 3 restart their counters, so
%% without this the supplement's Theorem 1 and the main paper's Theorem 1 would
%% both be cited as "Theorem 1" from the main text.
\renewcommand{\thetheorem}{S\arabic{theorem}}
\renewcommand{\thelemma}{S\arabic{lemma}}
\renewcommand{\thecorollary}{S\arabic{corollary}}
\renewcommand{\theproposition}{S\arabic{proposition}}
\renewcommand{\thedefinition}{S\arabic{definition}}
\renewcommand{\theassumption}{S\arabic{assumption}}
\renewcommand{\theremark}{S\arabic{remark}}
\renewcommand{\theconjecture}{S\arabic{conjecture}}

\doublespacing

\begin{document}

\title{Supplement to ``Model-Agnostic Interpretation of Machine
Learning via Average Marginal Effects''}
\author{Jason Parker\thanks{Princeton, Texas, USA. Email:
\texttt{jparker588@gmail.com}.}}
\date{\today}
\maketitle

This supplement reports two of the three empirical applications, the
detailed results of the three simulation studies, and proofs of all
results stated in the main paper. References of the form Theorem~1,
Lemma~2, Section~2.4 and so on refer to the main paper; sections,
tables and equations carrying an \textsf{S} prefix are internal to this
supplement.

		% ===== moved from Section 3.2 of the main paper =====
	\section{Function Spaces and Well-Behaved Estimators}
	\label{app:spaces}

	The gradient results of the next two sections are stated in two
	function spaces, each the natural habitat of an estimator class.
	This section fixes them and the regularity conditions they carry.

	We work with two function spaces depending on the estimator class. The choice of
	function space is not arbitrary---each space is the natural habitat of the estimator
	class in question, in the sense that estimators of that class belong to the space
	with probability one under mild conditions.
	
	The \textbf{Sobolev space} $W^{k,p}(\Omega)$ for $k \geq 1$ and $p \geq 1$
	consists of functions whose weak derivatives up to order $k$ are in $\Lp$, with
	norm:
	\begin{equation}
		\norm{g}_{W^{k,p}} = \sum_{|\alpha| \leq k} \norm{\partial^\alpha g}_{L^p}
	\end{equation}
	The gradient $\nabla g \in L^p(\Omega)^d$ is a vector valued $L^p$ function for any
	$g \in \Wonep$. Sobolev spaces are the natural setting for smooth estimators because
	they quantify not just how well the function is approximated in $L^p$ but how well
	its derivatives are approximated simultaneously. This is exactly the property needed
	to justify gradient based interpretation---$L^p$ consistency of $\fn$ alone is
	insufficient because differentiation is an unbounded operator on $L^p$, and
	convergence of a function does not in general imply convergence of its derivatives
	without additional regularity \citep[Chapter~5]{evans2010pde}.
	
	The \textbf{space of functions of bounded variation} $\BV$ consists of functions
	$g \in \Lone$ whose distributional gradient $Dg$ is a vector valued Radon measure
	with finite total variation $|Dg|(\Omega) < \infty$. The $BV$ norm is:
	\begin{equation}
		\norm{g}_{BV} = \norm{g}_{L^1} + |Dg|(\Omega)
	\end{equation}
	The $BV$ framework extends the notion of gradient to functions that are not
	differentiable in the classical or even weak sense---precisely the situation with
	tree based estimators, which have jump discontinuities at leaf boundaries. For
	piecewise constant functions the distributional gradient is a sum of surface measures
	concentrated on the boundaries between constant regions, capturing the intuition that
	the ``derivative'' of a step function is a spike at the jump location. This makes
	$BV$ the natural space for random forests and boosting with tree base learners. The
	total variation $|Dg|(\Omega)$ measures the total size of all jumps, weighted by the
	boundary surface area---a quantity that is finite and well controlled for any
	piecewise constant function with finitely many pieces. The standard reference for
	the theory of $BV$ functions is \citet{ambrosio2000bv}.
	
	The relationship between the two spaces is important. If $f \in \Wonep$ for
	$p \geq 1$, then $f \in \BV$ with $Df = \nabla f \cdot \mathcal{L}^d$---that is,
	the distributional gradient of $f$ is absolutely continuous with respect to Lebesgue
	measure with density $\nabla f$. This embedding $W^{1,p} \hookrightarrow BV$ means
	that assuming $f \in \Wonep$ is sufficient for both the Sobolev and BV results below.
	We make this assumption throughout.
	
	\begin{definition}[Well-Behaved Estimator, Sobolev Sense]
		\label{def:wellbehaved_sobolev}
		We say that $\fn$ is \emph{well-behaved in the Sobolev sense} if
		\begin{equation}
			\norm{\fn}_{W^{2,p}} = O_p(1) \quad \text{and} \quad
			\norm{\fn - f}_{L^p} \xrightarrowp 0 ,
			\label{eq:wb_sobolev}
		\end{equation}
		where the first condition means that the sequence is tight in
		$W^{2,p}$: for every $\eta > 0$ there is an $M < \infty$ with
		\begin{equation}
			\sup_n \Prob\!\left(\norm{\fn}_{W^{2,p}} > M\right) \leq \eta .
			\label{eq:tight_sobolev}
		\end{equation}
	\end{definition}
	
	\begin{definition}[Well-Behaved Estimator, BV Sense]
		\label{def:wellbehaved_bv}
		We say that $\fn$ is \emph{well-behaved in the BV sense} if:
		\begin{equation}
			\sup_n \E\left[|D\fn|(\Omega)\right] < \infty \quad \text{and} \quad
			\norm{\fn - f}_{L^1} \xrightarrowp 0
		\end{equation}
	\end{definition}
	
	In both cases the well-behaved condition has the same structure: a
	boundedness condition in a strong space combined with consistency in a
	weaker space. The boundedness condition is the key regularity
	requirement---it says that the estimator sequence does not become
	increasingly irregular as $n$ grows, even as it fits the data more closely.
	This condition is verifiable from properties of the estimator class and is
	established for each class in the theorems below. The consistency condition
	is the standard requirement that the estimator converges to the truth and is
	available from existing results in each case.

	The two boundedness conditions are stated differently, and the difference is
	not cosmetic. Both $\fn$ and its norms are random, so a condition of the form
	$\sup_n \norm{\fn} < \infty$ with a single deterministic bound is stronger
	than any of the estimators below deliver: the boosting bound of
	Section~\ref{app:boosting} is proportional to $\norm{Y}_{L^2}$, which is
	$O_p(1)$ but not bounded by a constant. Definition~\ref{def:wellbehaved_sobolev}
	therefore asks only for tightness. Definition~\ref{def:wellbehaved_bv} asks
	for a bound in expectation, which is what the Mondrian argument of
	Section~\ref{app:forests} produces directly; by Markov's inequality it
	implies tightness of $|D\fn|(\Omega)$, so the two conditions are of the same
	type, and either is implied by an almost sure uniform bound. What both
	exclude is norms that diverge in probability.

		% ===== moved from Section 3 of the main paper =====
	\section{What Learner Gradients Estimate}
	\label{sec:interpretation}
	
	The inference theory of Section~\ref{M-sec:inference} used no property of the fitted
	learner beyond an $L^2$ rate. The question it leaves open is interpretive: the
	first component of the score averages the finite differences of the fitted model,
	and one wants to know what that object measures about the model --- and when the
	model's gradient information converges to the population's. This subsection answers
	both questions in the function spaces the learners naturally inhabit. Four lemmas
	establish that gradients exist and converge in the Sobolev space of smooth
	estimators and the BV space of piecewise constant ones; a pairing identity
	(Theorem~\ref{M-thm:pairing}) then expresses the window effect of any BV estimator
	exactly as its jump measure integrated against a locally averaged covariate
	density, which is, to our knowledge, the first characterization of the population
	quantity that the marginal effect of a tree-based model measures. All proofs are in
	the supplement.
	
	The parallel structure of Lemmas~\ref{lem:sobolev_existence}
	through~\ref{lem:bv_consistency} reflects the parallel structure of the Sobolev
	and BV frameworks and is the formal expression of the analogy between smooth and
	piecewise constant estimators. None of these results is needed for the validity of
	Theorem~\ref{M-thm:main}; they characterize the regime in which the \emph{uncorrected}
	plug-in average is already consistent, and hence in which the correction term of
	\eqref{M-eq:dml} is a refinement rather than a repair.
	
	\begin{lemma}[Derivative Existence in Sobolev Spaces]
		\label{lem:sobolev_existence}
		Let $\fn \in \Wtwop$ for some $p > 1$. Then $\nabla \fn \in W^{1,p}(\Omega)^d$
		exists as a well defined $L^p$ function and satisfies:
		\begin{equation}
			\norm{\nabla \fn}_{L^p} \leq \norm{\fn}_{W^{1,p}} \leq \norm{\fn}_{W^{2,p}}
		\end{equation}
	\end{lemma}
	
	Lemma~\ref{lem:sobolev_existence} is essentially definitional---membership in
	$\Wtwop$ implies by definition that the first and second weak derivatives exist and
	are in $L^p$. The content of the lemma is the norm bound, which establishes that
	the gradient operator $\nabla: W^{2,p} \to W^{1,p}$ is bounded and linear. This
	boundedness is what makes the subsequent consistency argument work---a bounded linear
	operator preserves convergence, so once we have convergence of $\fn$ in the right
	space, convergence of $\nabla \fn$ follows automatically. The reason we require
	$\fn \in \Wtwop$ rather than just $\Wonep$ is that we need one additional degree of
	differentiability to apply the compactness argument in
	Lemma~\ref{lem:sobolev_consistency}. The relevant theory of Sobolev spaces and their
	norms is developed in \citet[Chapter~5]{evans2010pde}.
	
	\begin{lemma}[Derivative Consistency in Sobolev Spaces]
		\label{lem:sobolev_consistency}
		Let $f \in \Wonep$ and suppose $\fn$ is well-behaved in the Sobolev sense. Then:
		\begin{equation}
			\norm{\nabla \fn - \nabla f}_{L^p} \xrightarrowp 0
		\end{equation}
	\end{lemma}
	
	Lemma~\ref{lem:sobolev_consistency} is the harder result and the one that does the
	real work in the Sobolev setting. The difficulty is that $L^p$ consistency of $\fn$
	alone does not imply $L^p$ consistency of $\nabla \fn$---differentiation is an
	unbounded operator on $L^p$ and limits and derivatives cannot be exchanged without
	additional structure. The uniform $\Wtwop$ boundedness condition provides that
	structure. The proof applies the Rellich--Kondrachov compact embedding theorem, which
	states that a uniformly bounded sequence in $\Wtwop$ has a subsequence converging
	strongly in $\Wonep$ \citep[Theorem~5,~Section~5.7]{evans2010pde}. Combined with
	the $L^p$ consistency condition identifying the limit as $f$, this gives strong
	convergence of $\fn$ to $f$ in $\Wonep$, from which $L^p$ convergence of the
	gradients follows immediately from the boundedness of $\nabla$ established in
	Lemma~\ref{lem:sobolev_existence}. The key insight is that the uniform boundedness
	condition ensures the estimator sequence lives in a compact subset of $\Wonep$,
	preventing the gradients from escaping to infinity even as the estimator fits the
	data more closely.
	
	\begin{lemma}[Derivative Existence in BV]
		\label{lem:bv_existence}
		Let $\fn \in \BV$. Then the distributional gradient $D\fn$ exists as a well defined
		vector valued Radon measure in $\cM$ with:
		\begin{equation}
			|D\fn|(\Omega) < \infty
		\end{equation}
		For piecewise constant estimators with leaf partition $\{R_j\}$ and leaf means
		$\{\bar{y}_j\}$, the distributional gradient takes the explicit form:
		\begin{equation}
			D\fn = \sum_{j \sim k} (\bar{y}_k - \bar{y}_j) \cdot \nu_{jk} \cdot
			\mathcal{H}^{d-1}\lfloor_{\partial R_j \cap \partial R_k}
		\end{equation}
		where $\nu_{jk}$ is the unit normal between adjacent leaves and $\mathcal{H}^{d-1}$
		is the $(d-1)$-dimensional Hausdorff measure on leaf boundaries.
	\end{lemma}
	
	Lemma~\ref{lem:bv_existence} is the BV analog of Lemma~\ref{lem:sobolev_existence}
	and is similarly structural in nature. The first part---that $BV$ membership implies
	the distributional gradient exists as a Radon measure---follows from the definition
	of $BV$ as developed in \citet[Chapter~3]{ambrosio2000bv}. The content of the lemma
	is the explicit formula for piecewise constant estimators, which makes concrete what
	the distributional gradient actually looks like for a tree based model. The formula
	says that the gradient is not a function at all but a measure concentrated on the
	boundaries between leaves, with mass proportional to the jump in response values
	across the boundary weighted by the boundary surface area. This is the precise sense
	in which a random forest or boosted tree has a gradient---not a pointwise derivative,
	which does not exist, but a distributional object that captures where and how much
	the estimator changes. This explicit representation is also what makes the total
	variation $|D\fn|(\Omega)$ computable in terms of observable quantities---the leaf
	means and the partition geometry---which is essential for verifying the well-behaved
	condition in the theorems below.
	
	\begin{lemma}[Derivative Consistency in BV]
		\label{lem:bv_consistency}
		Let $f \in \Wonep$ for some $p > 1$, so that $f \in \BV$ with
		$Df = \nabla f \cdot \mathcal{L}^d$. Suppose $\fn$ is well-behaved in the BV sense.
		Then:
		\begin{equation}
			D\fn \wstar Df \quad \text{in } \cM
		\end{equation}
		If additionally $\E[|D\fn|(\Omega)] \to |Df|(\Omega)$, then convergence is strict.
	\end{lemma}
	
	Lemma~\ref{lem:bv_consistency} is the hardest of the four function space lemmas and
	the BV analog of Lemma~\ref{lem:sobolev_consistency}. The convergence is necessarily
	weaker than in the Sobolev case---weak-$*$ convergence of measures rather than strong
	$L^p$ convergence---because the $BV$ space is not reflexive and the stronger norm
	convergence is not available in general. Weak-$*$ convergence means that for any
	continuous test function $\phi$:
	\begin{equation}
		\int_\Omega \phi \, d(D\fn) \to \int_\Omega \phi \, d(Df)
	\end{equation}
	which is the appropriate notion of convergence for the distributional gradients of
	piecewise constant estimators. In the context of average marginal effects, this means
	that integrated gradient information---which is exactly what the AME computes---
	converges to the correct population quantity. The proof applies the BV compactness
	theorem of \citet[Theorem~3.23]{ambrosio2000bv}, which is the analog of
	Rellich--Kondrachov in the BV setting: a uniformly bounded sequence in $BV$ has a
	subsequence converging in $L^1$, and the $L^1$ consistency condition identifies the
	limit as $f$. The strict convergence upgrade follows from the lower semicontinuity of
	the total variation functional and is the stronger result one would want for
	applications requiring convergence of the total size of gradient information rather
	than just its distribution.
	
	The parallel between Lemmas~\ref{lem:sobolev_consistency}
	and~\ref{lem:bv_consistency} is worth noting explicitly. In both cases the argument
	has the same structure: uniform boundedness in a strong space gives compactness,
	consistency in a weak space identifies the limit, and the combination gives
	convergence of the derivative in the appropriate topology. The difference is in the
	topology of the conclusion---strong $L^p$ convergence in the Sobolev case, weak-$*$
	convergence of measures in the BV case---and this difference is fundamental rather
	than technical. It reflects the fact that smooth estimators carry gradient
	information in a stronger sense than piecewise constant estimators, which is itself a
	reflection of the different inductive biases of the estimator classes.
	
	The pairing identity, Theorem~\ref{M-thm:pairing} of the main paper, is the
	capstone of the interpretation theory: by Lemma~\ref{lem:bv_existence} the
	distributional gradient of a piecewise constant estimator is a measure
	concentrated on leaf boundaries, and the identity integrates that measure
	against the smoothed density $K_h \star q$ of \eqref{M-eq:movingaverage}. One
	consequence of it is worth recording separately.

	\begin{corollary}
		\label{cor:pairing_limit}
		If in addition $q$ extends to a continuous function of compact support (for
		instance $w$ continuous and $p$ continuous on $\Omega$), then for fixed
		$g \in \BV$:
		\begin{equation}
			\lim_{h \downarrow 0} \int_\Omega (K_h \star q)\; dD_j g
			= \int_\Omega q\; dD_j g
		\end{equation}
		and when $g \in W^{1,1}(\Omega)$ the right side equals
		$\E[w(X)\, \partial_j g(X)]$, consistent with
		Proposition~\ref{M-prop:hlimit}.
	\end{corollary}
	
	Theorem~\ref{M-thm:pairing} settles a question that the interpretability literature
	has left informal: what does the average marginal effect of a random forest
	\emph{measure}? The answer given by \eqref{M-eq:pairing_tree} is exact and
	geometric: it is the total jump of the fitted surface across leaf boundaries in
	direction $j$, with each boundary segment weighted by its share of the covariate
	mass in an $h$-window --- the places where the model changes its prediction,
	weighted by how many observations live near those places. The identity reconciles
	the two readings of a marginal effect that Section~2 could only juxtapose. Read
	with $g = \hat{f}$, \eqref{M-eq:pairing} is a descriptive statement about the fitted
	model, available in finite samples with no assumptions at all. Read through
	Theorem~\ref{M-thm:main}, the same plug-in quantity, corrected by the orthogonal
	score, estimates the population functional $\theta_h(f)$ at the parametric rate.
	The descriptive and the inferential interpretations are the two sides of
	\eqref{M-eq:dml}, and the pairing identity is what makes the first of them precise.
	
	\section{Gradient Convergence of Specific Learners}
	\label{sec:learners}
	
	The theorems of this subsection verify the well-behaved conditions for three
	estimator classes and conclude, via the lemma pairs above, that the learner's own
	gradient information converges to the population's. This is the strongest
	interpretability property a learner can have --- the fitted model's local behaviour,
	not merely its window averages, tracks the truth --- and it implies consistency of
	the uncorrected plug-in average for the small-$h$ estimand. We emphasize the
	division of labor: these results calibrate how much interpretive weight the fitted
	gradients themselves can bear, while the inference of Section~\ref{M-sec:inference}
	requires none of them. In each case the argument has two components: verifying that
	the estimator satisfies the well-behaved condition, and citing the appropriate
	lemma pair. The first component is the estimator-specific contribution; the second
	is the functional analytic machinery established above.
	
	\begin{theorem}[Gradient Consistency for Neural Networks]
		\label{thm:neural_nets}
		Let $\fn$ be a neural network estimator with smooth activations $\sigma \in C^2$,
		bounded weight norms $\prod_\ell \norm{W_\ell^{(n)}} = O_p(1)$, and
		$\norm{\fn - f}_{L^p} \xrightarrowp 0$ for $p > d$. Then $\fn$ is well-behaved in
		the Sobolev sense and:
		\begin{equation}
			\norm{\nabla \fn - \nabla f}_{L^p} \xrightarrowp 0
		\end{equation}
	\end{theorem}
	
	The conditions of Theorem~\ref{thm:neural_nets} are natural and verifiable. Smooth
	activations --- sigmoid, tanh, GELU --- are standard in practice and ensure the network
	function is twice differentiable. The bounded weight norm condition formalizes the
	effect of $L^2$ regularization during training, which prevents the network weights from
	growing without bound. The $L^p$ consistency condition is the standard requirement that
	the network is learning the true function. \citet{hornik1990universal} establish that
	multilayer feedforward networks with smooth activation functions are capable of
	approximating an arbitrary function and its derivatives to arbitrary accuracy, providing
	the foundation for the consistency condition and establishing that the function class is
	rich enough to contain the estimators we consider. The condition $p > d$ is required by
	the pointwise differentiability theorem of \citet[Theorem~6.5]{evans2015measure} and
	is satisfied for all $p$ larger than the covariate dimension --- a mild regularity
	condition that is implicit in the classical Sobolev embedding theory.
	
	\begin{proof}[Proof sketch]
		The chain rule applied through the network layers gives $\fn \in \Wtwop$ with:
		\begin{equation}
			\norm{\fn}_{W^{2,p}} \lesssim \prod_\ell \norm{W_\ell^{(n)}}^2 \cdot
			\norm{\sigma''}_\infty
		\end{equation}
		which is $O_p(1)$ by the weight norm assumption, the fitted weights being
		random. The $L^p$ consistency condition and
		Lemma~\ref{lem:sobolev_consistency} then give the result. Full proof
		in Section~\ref{app:neural_nets} of the supplement.
	\end{proof}
	
	\begin{remark}
		Theorem~\ref{thm:neural_nets} is stated for smooth activations. ReLU networks
		require a modified argument because ReLU is not $C^2$ and the chain rule calculation
		does not apply directly. Extension to ReLU networks is possible through a
		distributional argument but complicates the presentation and is omitted here.
	\end{remark}
	
	\begin{theorem}[Gradient Consistency for Gradient Boosted Models]
		\label{thm:boosting}
		Let $\fn = \hat{F}_m = (I - (I-S_n)^{m+1})Y$ be an $L_2$Boost estimator
		with a \emph{fixed} symmetric linear learner $S_n$ whose eigenvalues
		satisfy $0 < \lambda_k^{(n)} \leq 1$ for all $k$, and with base learners
		$h_m = S_n u_m$ lying in $W^{2,p}$. Suppose in addition that the
		spectral gap and the smoothing bound are \emph{uniform in $n$}:
		\begin{equation}
			\bar{\rho} := \sup_n \max_k \left(1 - \lambda_k^{(n)}\right) < 1,
			\qquad
			\sup_n \norm{S_n}_{L^2 \to W^{2,p}} < \infty .
			\label{eq:boost_uniform}
		\end{equation}
		Suppose $\norm{\fn - f}_{L^p} \xrightarrowp 0$. Then $\fn$ is
		well-behaved in the Sobolev sense and
		$\norm{\nabla \fn - \nabla f}_{L^p} \xrightarrowp 0$.
	\end{theorem}

	The consistency hypothesis is in $L^p$, the same exponent as the
	conclusion and as the second condition of
	Definition~\ref{def:wellbehaved_sobolev}; an $L^1$ hypothesis would not
	suffice, since on a bounded domain $L^p$ consistency implies $L^1$
	consistency and not the reverse. For $p \leq 2$ it is implied by $L^2$
	consistency, again by H\"{o}lder on the bounded domain $\Omega$, and $L^2$
	is the exponent in which the risk results of \citet{buhlmann2003boosting}
	for this estimator are stated.
	
	Theorem~\ref{thm:boosting} is deliberately narrow, and it is worth being
	explicit about what it does not reach. It applies to $L_2$Boost driven by a
	single fixed symmetric linear smoother --- a smoothing spline at a fixed
	smoothing parameter is the canonical example --- for which
	\citet{buhlmann2003boosting} supply the representation and the spectral
	decay. It does \emph{not} apply to XGBoost or LightGBM, the implementations
	used in Section~\ref{M-sec:empirical}. Those grow each tree greedily from the
	current residual, so there is no fixed operator $S$, the representation
	$\hat{F}_m = (I - (I-S)^{m+1})Y$ fails, and the geometric decay that the
	proof turns on is unavailable. Remark~\ref{rem:boosting_trees} sets out why
	no choice of tree base learner can be brought inside the hypotheses, and what
	is and is not known about those implementations in consequence.
	
	The eigenvalue condition $0 < \lambda_k \leq 1$ is the key regularity condition on
	the base learner. \citet{buhlmann2003boosting} impose this condition explicitly in
	their Theorem~1 and establish that under it the bias decays exponentially fast while
	the variance increases only by exponentially diminishing amounts --- the exponential
	bias--variance trade-off that makes $L_2$Boost analytically tractable. The condition
	holds for smoothing spline learners and, generically, for tree stump learners whose
	splits have positive but incomplete explanatory power, which is the typical case in
	practice. Learners with any eigenvalue equal to zero are degenerate in the sense of
	Corollary~1 of \citet{buhlmann2003boosting}: $\mathcal{B}_m \equiv S$ for all $m$ and
	boosting does not improve on the base learner. The functional gradient descent
	perspective of \citet{mason2000boosting} provides the conceptual framework for
	understanding the boosting path as gradient descent in function space.
	
	\begin{proof}[Proof sketch]
		By Proposition~1 of \citet{buhlmann2003boosting} the residual at step $m$
		is $u_m = (I - S_n)^m Y$, and by Proposition~2 and Theorem~1, under the
		eigenvalue condition, it decays geometrically:
		\begin{equation}
			\norm{u_m}_{L^2} \leq C\rho^m,
			\qquad \rho = \max_k\left(1-\lambda_k^{(n)}\right) \leq \bar{\rho} < 1,
			\qquad C = \norm{Y}_{L^2},
		\end{equation}
		in the normalized empirical $L^2$ norm, for which $C = O_p(1)$ under
		Assumption~\ref{M-ass:sampling}. Each base learner fit satisfies
		$\norm{h_m}_{W^{2,p}} \leq \norm{S_n}_{L^2 \to W^{2,p}} \norm{u_m}_{L^2}$,
		so by linearity and the triangle inequality
		\begin{equation}
			\norm{\fn}_{W^{2,p}}
			\leq \sum_{m \geq 1} \norm{h_m}_{W^{2,p}}
			\leq \frac{C\,\bar{\rho}}{1 - \bar{\rho}} \sup_n \norm{S_n}_{L^2 \to W^{2,p}}
			= O_p(1),
		\end{equation}
		the two suprema in \eqref{eq:boost_uniform} making the bound free of $n$
		apart from the factor $C$. This is the first condition of
		Definition~\ref{def:wellbehaved_sobolev}; the second is assumed, and
		Lemma~\ref{lem:sobolev_consistency} gives the conclusion. Full proof in
		Section~\ref{app:boosting} of the supplement.
	\end{proof}
	
	\begin{theorem}[Gradient Consistency for Mondrian Forests]
		\label{thm:forests}
		Let $\fn$ be a Mondrian random forest estimator with budget $\lambda_n \asymp
		n^{1/(d+2)}$ and any fixed number $B \geq 1$ of trees, with $f$ Lipschitz
		on $[0,1]^d$. Then $\fn$ is well-behaved in the BV sense and:
		\begin{equation}
			D\fn \wstar Df \quad \text{in } \cM
		\end{equation}
	\end{theorem}
	
	Theorem~\ref{thm:forests} is the most technically demanding result in the paper.
	Random forests are genuinely different from neural networks and boosted models in a
	way that makes gradient consistency harder to establish --- the partition structure is
	random and the response averages within leaves depend on the data in a complex way.
	The restriction to Mondrian forests, where splits are made independently of the data,
	is the key simplification that makes the proof tractable. For Mondrian forests the
	partition geometry and the response variation can be controlled separately, whereas for
	Breiman forests with data-dependent splits the two are coupled in a way that resists
	clean analysis. The $L^2$ consistency result under the budget $\lambda_n \asymp
	n^{1/(d+2)}$ follows from \citet{mourtada2017mondrian}, who establish the minimax rate
	$O(n^{-2/(d+2)})$ for Lipschitz regression functions, and $L^1$ consistency follows
	by Cauchy--Schwarz on the bounded domain $\Omega$.
	
	The budget condition $\lambda_n \asymp n^{1/(d+2)}$ is the only condition the
	theorem imposes, and it is the natural one: it balances the bias from coarse
	partitions against the variance from fine ones at the minimax optimal rate.
	What is perhaps surprising is that the same rate balances the two
	contributions to the \emph{total variation} of the fitted tree. The jumps
	across leaf boundaries carry a signal part, which reflects the variation of
	$f$ itself, and a noise part, which reflects the sampling error in the leaf
	means; at this budget both are $O(1)$, and the partition may refine without
	limit without the total variation growing. No condition on the number of
	trees is needed.
	
	\begin{proof}[Proof sketch]
		Each tree is piecewise constant and therefore in $BV$ by
		Lemma~\ref{lem:bv_existence}, and the forest average is a convex
		combination, so $|D\fn|(\Omega) \leq B^{-1}\sum_b |DT_b|(\Omega)$. It
		therefore suffices to bound a single tree, and no cancellation between
		trees is required.

		For one tree, the distributional gradient is carried by the split faces
		(Lemma~\ref{lem:nodeface}), and the jump across a face splits into a
		signal part and a noise part. Reducing both to one-dimensional sections
		turns the signal part into the expected aspect ratio
		$\E[\mathrm{diam}(l(X))/s_1(l(X))]$ of the cell containing a uniform
		point, which Lemma~\ref{lem:aspect} bounds by $2d-1+O(1/\lambda)$ using
		the Mondrian cell law of \citet{mourtada2017mondrian} --- a bound with no
		rate in it, because the variation of a Lipschitz function over a refining
		partition depends on the partition's shape and not on how fine it is. The
		noise part is bounded by Lemma~\ref{lem:noise} at
		$O(\lambda^{1+d/2}n^{-1/2})$, which is $O(1)$ precisely when
		$\lambda \asymp n^{1/(d+2)}$, since $(1+d/2)/(d+2) = 1/2$.
		Lemma~\ref{lem:bv_consistency} then gives $D\fn \xrightarrow{w*} Df$ in
		$\cM$. Full proof in Section~\ref{app:forests}.
	\end{proof}
	
	\begin{remark}
		Extension of Theorem~\ref{thm:forests} to Breiman forests with data-dependent splits
		requires additional control over the interaction between partition geometry and
		response variation. The data-dependent splitting rule couples these two quantities in
		a way that makes the cross-term cancellation argument unavailable --- splits are
		chosen precisely to maximize response differences across boundaries, which is the
		worst case for BV control. This extension is an important direction for future work
		and is discussed further in Section~\ref{M-sec:conclusion}.
	\end{remark}
	
	\section{Nuisance Estimation and Rates}
	\label{sec:rates}
	
	Assumption~\ref{M-ass:nuisance} asks for two regression-type $L^2$ rates whose
	product is $o_p(n^{-1/2})$; the benchmark sufficient condition is that each
	nuisance converge faster than $n^{-1/4}$. This subsection records when the
	learners of Section~2 meet the benchmark for $\hat{f}$, and how the representer
	$\alpha_h$ is estimated without ever estimating a density.
	
	For the regression function, the required rates are exactly the $L^2$ risk bounds
	of the modern nonparametric literature, and a single smoothness threshold
	organizes all three learner classes. For deep neural networks with ReLU
	activations, \citet{schmidthieber2020} establishes the rate
	$n^{-s/(2s + d)}$ up to logarithmic factors for $s$-H\"{o}lder regression
	functions (with the effective $d$ reduced under compositional structure), so the
	$n^{-1/4}$ benchmark is met precisely when:
	\begin{equation}
		\frac{s}{2s + d} > \frac{1}{4}
		\quad \Longleftrightarrow \quad
		s > \frac{d}{2}
		\label{eq:threshold}
	\end{equation}
	For Mondrian forests, the minimax-optimal rates of \citet{mourtada2017mondrian}
	take the same form $n^{-s/(2s + d)}$ for H\"{o}lder smoothness $s \leq 2$, so the
	benchmark is again \eqref{eq:threshold}: at Lipschitz smoothness $s = 1$ this
	requires $d < 2$, and exploiting forest smoothing up to $s = 2$ accommodates
	$d \leq 3$. For $L_2$Boost, \citet{buhlmann2003boosting} establish minimax rates
	for smoothing-spline base learners with the boosting iteration acting as the
	regularization parameter, and the benchmark corresponds to smoothness exceeding
	one half in their univariate and additive settings --- with additive structure,
	as usual, breaking the dependence on the ambient dimension. The reading we intend
	is not that these conditions are innocuous --- \eqref{eq:threshold} is a real
	restriction, and in high dimensions it binds --- but that the currency in which
	the conditions are denominated has changed: what inference requires of the learner
	is an $L^2$ error rate of the kind the learning-theory literature actually
	proves, rather than root-$n$ convergence in a Sobolev norm, which no
	growing-complexity learner satisfies. For Breiman forests with data-dependent
	splits, $L^2$ consistency is available under structural assumptions
	\citep{scornet2015forests}, and polynomial $L^2$ rates for the original
	algorithm --- sample CART splits and column subsampling included --- have been
	established under sufficient-impurity-decrease conditions by
	\citet{chi2022forests}; the resulting exponents can fall short of the
	$n^{-1/4}$ benchmark, but the product form of
	Assumption~\ref{M-ass:nuisance}(b) permits pairing a slow forest rate with a fast
	representer estimate, and Proposition~\ref{M-prop:robust} covers the
	known-representer case with no forest rate at all. The open question relevant to
	this paper is therefore an $L^2$ rate question, a substantially more tractable
	target than the gradient-consistency question that
	Section~\ref{sec:learners} leaves open for that class.
	
	The representer is estimated by Riesz regression
	\citep{chernozhukov2022riesz}, which exploits the same linearity that produced
	Lemma~\ref{M-lem:orth}. By Lemma~\ref{M-lem:riesz}, for any candidate
	$\alpha \in L^2(P)$:
	\begin{equation}
		\E\left[\alpha(X)^2\right] - 2\,\E\left[w(X)\, \Dh \alpha(X)\right]
		= \Ptwo{\alpha - \alpha_h}^2 - \Ptwo{\alpha_h}^2
		\label{eq:rieszloss}
	\end{equation}
	so the population loss on the left is minimized uniquely at $\alpha_h$, and its
	sample analog:
	\begin{equation}
		\hat{L}(\alpha) = \frac{1}{n} \sum_{i=1}^{n} \left[
		\alpha(\mathbf{x}_i)^2
		- 2\, w(\mathbf{x}_i)\,
		\frac{\alpha(\mathbf{x}_i + h e_j) - \alpha(\mathbf{x}_i - h e_j)}{2h}
		\right]
	\end{equation}
	is computable directly: minimizing $\hat{L}$ over a class of functions ---
	penalized linear sieves, neural networks, or forests --- requires only the ability
	to evaluate candidates at shifted points, and at no stage is the density $p$, let
	alone a derivative of it, estimated. \citet{chernozhukov2022riesz} provide
	$L^2$ convergence guarantees for such estimators over exactly these classes, with
	rates governed by the complexity of the class and the approximability of
	$\alpha_h$. Since $\alpha_h$ obeys the explicit bound of Lemma~\ref{M-lem:riesz},
	the minimizer is truncated at that bound, which enforces
	Assumption~\ref{M-ass:nuisance}(c) by construction without affecting the $L^2$
	error. Two further remarks. First, the product form of
	Assumption~\ref{M-ass:nuisance}(b) permits asymmetry: a learner with a fast rate
	for $f$ tolerates a slow Riesz estimate, and conversely. Second, in designed or
	simulated data where the covariate density is known --- as in the Monte Carlo
	study of Section~5 --- $\alpha_h$ is available in closed form and the representer
	error is identically zero; Proposition~\ref{M-prop:robust} then applies, and valid
	inference requires no rate --- indeed no consistency --- from the learner at all.
	
	\section{Verifying a New Learner}
	\label{sec:newlearners}
	
	The architecture above is deliberately modular, and it is worth stating plainly
	what a reader must check to bring a learner not treated here under the results.
	There are two obligations, and they are separate. Valid inference requires only
	an $L^2(P)$ convergence rate for $\hat{f}$ fast enough that its product with the
	representer rate is $o_p(n^{-1/2})$, the benchmark sufficient condition being
	$n^{-1/4}$ for each. Nothing else about the learner enters
	Theorem~\ref{M-thm:main}: not differentiability, not the existence of a gradient,
	not membership in any particular function space. Consistency of the learner's
	own gradient information --- the object of Section~\ref{sec:learners} --- is a
	separate question, required only if one wishes to interpret the uncorrected
	plug-in, and it is answered by identifying the natural function space of the
	estimator class and applying the appropriate compactness theorem.
	
	The practical consequence is that adding a learner is usually a matter of
	locating an existing rate rather than proving a new theorem.
	Table~\ref{tab:learner_checklist} records the position for the classes treated
	here and for several that are not. Penalized linear estimators are the easiest
	case: the lasso attains $L^2$ rates of order $\sqrt{s \log d / n}$ under
	standard sparsity and restricted-eigenvalue conditions, comfortably inside the
	benchmark whenever $s \log d = o(\sqrt{n})$, and it is a linear estimator in a
	Sobolev space, so both obligations are discharged by citation. Smoothing splines
	and series estimators are similar on both counts, attaining the classical
	$n^{-s/(2s+1)}$ rate in the univariate case and living in $W^{1,p}$ by
	construction. Kernel ridge regression over a reproducing kernel Hilbert space --- the function
	space generated by the chosen kernel, in which evaluation at a point is a
	continuous linear functional --- has minimax $L^2$ rates of the same form. For a
	sufficiently smooth kernel that space embeds continuously in $W^{1,p}$, so the
	interpretation argument of Section~\ref{sec:learners} applies unchanged; for a
	kernel that is not smooth it need not, and the embedding is the thing to check. Generalized additive models
	inherit the univariate rate componentwise, which is what makes additive
	structure so useful here: it breaks the dependence on ambient dimension in
	\eqref{eq:threshold}. For support vector regression with a smooth kernel and for
	Bayesian additive regression trees, $L^2$ rates are available under conditions
	in the respective literatures, though in the latter case posterior contraction
	rather than a frequentist risk bound is the natural object and the translation
	requires care. Any learner of bounded variation --- which includes every
	tree-based method, however its splits are chosen --- falls under the pairing
	identity of Theorem~\ref{M-thm:pairing} without further argument, so what its
	window effect measures is characterized even when its rate is unknown.
	
	Two cases in the table deserve separate mention because they are where the
	framework is currently incomplete. Breiman forests have the inference obligation
	open and the interpretation obligation conjectured, for the reasons given above.
	Learners whose output is not a fixed function of the data --- online or
	streaming estimators, for instance --- fall outside the sampling framework of
	Assumption~\ref{M-ass:sampling} altogether and would require a different
	treatment.

\begin{table}[htbp]
	\centering
	\caption{What must be verified to apply the framework to a learner. The
		inference column is what Theorem~\ref{M-thm:main} consumes; the interpretation
		column is what Section~\ref{sec:learners} requires in order to read the
		uncorrected plug-in as a gradient.}
	\label{tab:learner_checklist}
	\begin{tabular}{lll}
		\toprule
		Learner & Inference: $L^2$ rate & Interpretation: function space \\
		\midrule
		Linear and GLM        & parametric        & $W^{1,p}$, exact \\
		Lasso, ridge          & established       & $W^{1,p}$ \\
		Splines, series       & established       & $W^{1,p}$ \\
		Kernel ridge          & established       & $W^{1,p}$ via smooth kernel \\
		GAM                   & componentwise     & $W^{1,p}$ \\
		Neural networks       & established       & $W^{1,p}$, proved here \\
		$L_2$Boost, fixed smoother & established  & $W^{1,p}$/BV, proved here \\
		Tree boosting (XGBoost) & established     & BV; gradient \emph{open} \\
		Mondrian forests      & minimax optimal   & BV, proved here \\
		Breiman forests       & \emph{open}       & BV; gradient \emph{conjectured} \\
		SVM (smooth kernel)   & available         & $W^{1,p}$ via kernel; not treated \\
		BART                  & contraction rate  & BV; not treated \\
		\bottomrule
	\end{tabular}
\end{table}


		% ===== moved from Section 4 of the main paper =====
	\section{Implementation Details for the Empirical Applications}
	\label{app:empirical_details}

	Across the three empirical applications we fit four model classes: logistic or linear regression as the parametric baseline, a random forest, XGBoost, and a feedforward neural network. The random forest uses 200 trees with a maximum depth of 8. XGBoost uses 200 boosting rounds with a maximum depth of 4 and a learning rate of 0.05. The neural network architecture and training schedule depend on the outcome type. For binary classification, the network is a feedforward architecture with two hidden layers of 64 and 32 units trained with the Adam optimizer at a learning rate of $10^{-3}$, with early stopping on a 10\% validation split with patience of 20 epochs. For continuous regression, the network is deeper --- four hidden layers of 256, 128, 64, and 32 units --- trained at a learning rate of $3 \times 10^{-4}$ with patience of 50 epochs, reflecting the larger capacity required to fit the continuous outcome surface and the longer training horizon needed for the optimization to converge. For regression outcomes, the network includes the input normalization layer and target standardization wrapper described in Section~2 to ensure scale-invariant gradient updates and predictions on the natural outcome scale. Estimation follows the debiased cross-fitted procedure of Section~2.4 with $K=5$ folds: within each fold the learner and the Riesz representer are fit on the other four folds and evaluated on the held-out fold, and standard errors come from the sample variance of the orthogonal score. No learner is ever differenced at an observation it was trained on. Every model class is refit from scratch on each fold, including the neural networks, whose early stopping schedule is applied within the fold. The fit statistics reported at the foot of each table are computed from a single full-sample fit and are descriptive of the learner rather than of the estimator; they are not used in the estimation. Where applicable, predicted probabilities for classification models are extracted from \texttt{predict\_proba(X)[:,1]} as discussed in Section~2.

	\begin{table}[htbp]
		\centering
		\small
		\begin{threeparttable}
		\caption{Window settings for the empirical applications. Step size $h_j$ and the fraction of observations receiving trimming weight zero.}
		\label{tab:window_settings}
		\begin{tabular}{llrr}
			\toprule
			Dataset & Feature & $h_j$ & Trimmed \\
			\midrule
			UCI Adult Income & age & 0.6855 & 1.3\% \\
			& female & --- & 0.0\% \\
			& education\_num & 0.5 & 1.4\% \\
			& hours\_per\_week & 0.6196 & 0.3\% \\
			& government & --- & 0.0\% \\
			& self\_employed & --- & 0.0\% \\
			& white\_collar & --- & 0.0\% \\
			& blue\_collar & --- & 0.0\% \\
			& service & --- & 0.0\% \\
			& capital\_gain & 373 & 92.3\% \\
			& capital\_loss & 20.15 & 95.3\% \\
			& married & --- & 0.0\% \\
			\midrule
			UCI Credit Card Default & age & 0.5 & 0.2\% \\
			& female & --- & 0.0\% \\
			& education & 0.5 & 36.8\% \\
			& married & --- & 0.0\% \\
			& avg\_pay\_status & 0.04911 & 7.0\% \\
			& months\_delayed & 0.5 & 70.9\% \\
			& avg\_bill\_amt & 3,163 & 0.0\% \\
			& avg\_pay\_amt & 507 & 11.7\% \\
			& pay\_ratio & 0.01839 & 18.4\% \\
			\midrule
			Ames Housing & GrLivArea & 26.27 & 0.1\% \\
			& BedroomAbvGr & 0.5 & 0.5\% \\
			& FullBath & 0.5 & 2.9\% \\
			& HalfBath & 0.5 & 63.4\% \\
			& OverallQual & 0.5 & 1.4\% \\
			& OverallCond & 0.5 & 1.6\% \\
			& YearBuilt & 1.51 & 1.4\% \\
			& YearRemodAdd & 1.032 & 14.5\% \\
			& LotArea & 499 & 1.2\% \\
			& GarageArea & 10.69 & 5.6\% \\
			& nbhd\_high & --- & 0.0\% \\
			& nbhd\_mid & --- & 0.0\% \\
			\bottomrule
			\end{tabular}
			\begin{tablenotes}
			\footnotesize
			\item \textit{Notes:} $h_j$ is the default adaptive step size, $\max(10^{-4},\, 0.05\hat{\sigma}_j)$, floored at $0.5$ for integer-valued features so that a count is contrasted over a whole unit. Features declared binary use the level contrast $f(x \mid x_j=1) - f(x \mid x_j=0)$, for which $h_j$ is undefined and no trimming applies. The trimmed fraction is the share of observations lying within $h_j$ of the boundary of the observed support of $x_j$, which receive weight zero. The estimand is the \emph{weighted} effect $\E[w(X)\Dh f(X)]$ of Definition~\ref{M-def:window}, not the effect among retained observations: by \eqref{M-eq:trimscale} the two differ by the factor $\E[w(X)]$, which is one minus the trimmed fraction. For a feature trimmed at $92.3\%$ the reported number is about one thirteenth of its conditional counterpart, so entries in this table are not comparable across features without that rescaling.
			\end{tablenotes}
		\end{threeparttable}
		\end{table}

	\section{Specification Search for the Credit Default Application}
	\label{app:credit_spec}

	\begin{table}[htbp]
		\centering
		\small
		\resizebox{\textwidth}{!}{%
			\begin{threeparttable}
				\caption{UCI Credit Card Default: Specification Search. Average marginal effects on P(Default).}
				\label{tab:credit_default_spec_search}
				\begin{tabular}{lrrrrrrrr}
					\toprule
					& \multicolumn{4}{c}{Logistic} & \multicolumn{4}{c}{Neural Net} \\
					\cmidrule(lr){2-5} \cmidrule(lr){6-9}
					& Spec A & A+B & A+C & A+B+C & Spec A & A+B & A+C & A+B+C \\
					\midrule
					\multicolumn{9}{l}{\textit{Panel A: Demographics}} \\
					\quad age & -0.0015$^{***}$ & -0.0002 & -0.0005 & 0.0001 & -0.0026$^{***}$ & -0.0006$^{*}$ & -0.0010$^{***}$ & -0.0000 \\
					& (0.0004) & (0.0003) & (0.0004) & (0.0003) & (0.0004) & (0.0003) & (0.0004) & (0.0003) \\
					\quad female & -0.0358$^{***}$ & -0.0206$^{***}$ & -0.0291$^{***}$ & -0.0198$^{***}$ & -0.0364$^{***}$ & -0.0205$^{***}$ & -0.0265$^{***}$ & -0.0176$^{***}$ \\
					& (0.0050) & (0.0046) & (0.0049) & (0.0046) & (0.0050) & (0.0046) & (0.0049) & (0.0045) \\
					\quad education & -0.0001 & -0.0027 & -0.0054$^{**}$ & -0.0040$^{*}$ & -0.0008 & -0.0035 & -0.0048$^{**}$ & -0.0040$^{*}$ \\
					& (0.0022) & (0.0022) & (0.0023) & (0.0022) & (0.0021) & (0.0021) & (0.0022) & (0.0022) \\
					\quad married & 0.0253$^{***}$ & 0.0189$^{***}$ & 0.0303$^{***}$ & 0.0206$^{***}$ & 0.0260$^{***}$ & 0.0152$^{**}$ & 0.0269$^{***}$ & 0.0186$^{***}$ \\
					& (0.0068) & (0.0064) & (0.0066) & (0.0063) & (0.0068) & (0.0064) & (0.0067) & (0.0063) \\
					\midrule
					\multicolumn{9}{l}{\textit{Panel B: Payment History}} \\
					\quad avg\_pay\_status & -- & -0.0140$^{***}$ & -- & -0.0218$^{***}$ & -- & 0.0073$^{*}$ & -- & 0.0040 \\
					&  & (0.0041) &  & (0.0075) &  & (0.0041) &  & (0.0072) \\
					\quad months\_delayed & -- & 0.0347$^{***}$ & -- & 0.0301$^{***}$ & -- & 0.0298$^{***}$ & -- & 0.0254$^{***}$ \\
					&  & (0.0011) &  & (0.0012) &  & (0.0011) &  & (0.0012) \\
					\midrule
					\multicolumn{9}{l}{\textit{Panel C: Bill and Payment Amounts}} \\
					\quad avg\_bill\_amt & -- & -- & -0.0000$^{***}$ & -0.0000$^{**}$ & -- & -- & -0.0000 & -0.0000$^{*}$ \\
					&  &  & (0.0000) & (0.0000) &  &  & (0.0000) & (0.0000) \\
					\quad avg\_pay\_amt & -- & -- & -0.0000$^{***}$ & 0.0000 & -- & -- & -0.0000$^{***}$ & -0.0000$^{***}$ \\
					&  &  & (0.0000) & (0.0000) &  &  & (0.0000) & (0.0000) \\
					\quad pay\_ratio & -- & -- & -0.3521$^{***}$ & -0.2651$^{***}$ & -- & -- & -0.2572$^{***}$ & -0.1943$^{**}$ \\
					&  &  & (0.0731) & (0.0730) &  &  & (0.0792) & (0.0773) \\
					\midrule
					McFadden $R^2$ & 0.003 & 0.130 & 0.031 & 0.137 & 0.010 & 0.155 & 0.055 & 0.166 \\
					Accuracy & 0.779 & 0.803 & 0.779 & 0.803 & 0.779 & 0.808 & 0.779 & 0.808 \\
					\bottomrule
					\end{tabular}
					\begin{tablenotes}
					\footnotesize
					\item \textit{Notes:} $^{***}$p$<$0.01, $^{**}$p$<$0.05, $^{*}$p$<$0.10. Estimates are the debiased cross-fitted window AME with $K=5$ folds; standard errors from the orthogonal score in parentheses. Step sizes $h_j$ and trimmed fractions are reported in Table~\ref{tab:window_settings}.
					\end{tablenotes}
			\end{threeparttable}
		}%
	\end{table}
	
	Table~\ref{tab:credit_default_spec_search} presents the specification search. The dominant finding is the importance of payment history and of the payment ratio. Months delayed enters at 0.035 for logistic regression and 0.030 for the neural network in specification A+B and remains at 0.030 and 0.025 in the full specification, while the payment ratio, once it is available in specifications A+C and A+B+C, is much the largest effect in the table: $-$0.265 for logistic regression and $-$0.194 for the neural network in the full specification. The demographic panel tells a different story from the Adult Income dataset of Section~\ref{app:adult}: age is statistically indistinguishable from zero in nearly every specification for both models, and the female indicator, while consistently negative and significant, shrinks substantially once payment history is controlled for, falling from $-$0.036 to $-$0.020 for logistic regression and from $-$0.036 to $-$0.018 for the neural network. This pattern of attenuation is cleanly legible in the AME framework because the estimates are on a common probability scale across specifications.

	What is most striking in this table is how little separates the two columns. In specification A the female AME is $-$0.0358 for logistic regression and $-$0.0364 for the neural network; married is 0.0253 and 0.0260; in the full specification the pair is $-$0.0198 and $-$0.0176, and 0.0206 and 0.0186. A parametric model and a four-layer network, differing in McFadden $R^2$ by three percentage points, produce estimates that agree to the third decimal place across sixteen model-specification cells. This is what the orthogonal score is supposed to deliver and what the plug-in did not: since each column estimates the same $\theta_{j,h}$ and the correction absorbs the first-order effect of the fitted regression, the remaining gap between the two columns reflects only the second-order term and sampling noise. The standard errors are also now nearly identical between the two models --- 0.0046 and 0.0045 on the female indicator, 0.0012 and 0.0012 on months delayed --- rather than differing by the order of magnitude the refitting bootstrap reported, where the network's standard error for months delayed was 0.032 against the logistic model's 0.002. Both figures were describing the stability of a particular fitting procedure under resampling rather than the sampling variability of an estimate of a fixed target.

		\section{Comparison with SHAP and Partial Dependence for the
	Credit Default Application}
	\label{app:credit_methods}

	\begin{sidewaystable}[htbp]
		\centering
		\small
		\begin{threeparttable}
			\caption{UCI Credit Card Default: Method Comparison. AME, SHAP, and PDP estimates for P(Default), full specification (A+B+C).}
			\label{tab:credit_default_method_comparison}
			\begin{tabular}{lrrrrrrrrrrrr}
				\toprule
				& \multicolumn{3}{c}{Logistic} & \multicolumn{3}{c}{Random Forest} & \multicolumn{3}{c}{XGBoost} & \multicolumn{3}{c}{Neural Net} \\
				\cmidrule(lr){2-4} \cmidrule(lr){5-7} \cmidrule(lr){8-10} \cmidrule(lr){11-13}
				& AME & SHAP & PDP & AME & SHAP & PDP & AME & SHAP & PDP & AME & SHAP & PDP \\
				\midrule
				age & 0.0001 & -0.0003 & 0.0002 & -0.0002 & 0.0011 & -0.0000 & -0.0005 & 0.0004 & -0.0002 & -0.0000 & -0.0010 & 0.0012 \\
				female & -0.0198 & -0.0005 & -0.0198 & -0.0161 & 0.0001 & -0.0062 & -0.0175 & 0.0002 & -0.0146 & -0.0176 & 0.0009 & -0.0200 \\
				education & -0.0040 & 0.0002 & 0.0002 & -0.0018 & 0.0006 & -0.0057 & 0.0034 & -0.0005 & -0.0356 & -0.0040 & -0.0016 & 0.0082 \\
				married & 0.0206 & -0.0038 & 0.0220 & 0.0222 & 0.0002 & 0.0054 & 0.0205 & 0.0002 & 0.0142 & 0.0186 & 0.0022 & 0.0217 \\
				\midrule
				avg\_pay\_status & -0.0218 & -0.0144 & -0.0246 & 0.0061 & 0.0029 & 0.0610 & 0.0907 & 0.0027 & 0.0404 & 0.0040 & -0.0116 & 0.0097 \\
				months\_delayed & 0.0301 & -0.0148 & 0.1107 & 0.0134 & -0.0008 & 0.0456 & 0.0191 & 0.0001 & 0.0404 & 0.0254 & -0.0154 & 0.0653 \\
				\midrule
				avg\_bill\_amt & -0.0000 & -0.0001 & 0.0000 & -0.0000 & -0.0018 & -0.0000 & -0.0000 & -0.0046 & -0.0000 & -0.0000 & 0.0036 & 0.0000 \\
				avg\_pay\_amt & 0.0000 & 0.0026 & -0.0000 & -0.0000 & -0.0006 & -0.0000 & -0.0000 & -0.0005 & -0.0000 & -0.0000 & -0.0025 & -0.0000 \\
				pay\_ratio & -0.2651 & 0.0100 & -0.0444 & -0.4890 & 0.0012 & -0.0998 & -0.5836 & 0.0005 & -0.0953 & -0.1943 & 0.0015 & -0.0170 \\
				\bottomrule
				\end{tabular}
				\begin{tablenotes}
				\footnotesize
				\item \textit{Notes:} AME = Average Marginal Effect (marginfx). SHAP = mean signed SHAP value (GradientExplainer). PDP = slope of partial dependence plot. Full specification (A+B+C) only. No standard errors reported for SHAP and PDP.
				\end{tablenotes}
		\end{threeparttable}
	\end{sidewaystable}
	
	Table~\ref{tab:credit_default_method_comparison} again shows the two quantities parting company. For the payment history variables the separation is large: the AME for months delayed is 0.030 for logistic regression while the mean signed SHAP value is $-$0.015, opposite in sign. The payment ratio illustrates the most consequential divergence: the random forest and XGBoost AMEs are $-$0.489 and $-$0.584, indicating that a higher payment ratio is associated with a large reduction in default probability, while the corresponding SHAP values are 0.0012 and 0.0005, near zero and of the opposite sign. This is again a difference of estimand rather than a defect: a feature can carry a large average derivative while contributing little on average to the displacement of predictions from their baseline. The practical concern is that a researcher reading the SHAP column as a marginal effect would conclude that the payment ratio is essentially irrelevant to default risk, and, with no standard errors attached to either the SHAP or the PDP column, would have no statistical basis on which to check that conclusion.

	The partial dependence column separates from the AME column here in the same pattern as in the Ames application of Section~\ref{app:ames}, and for the same reason: a partial dependence slope is a plug-in average over a single full-sample fit, so it carries the regularization bias that the correction removes. The separation tracks the quality of the underlying fit. On the female indicator, where the logistic model is close to correctly specified, the AME of $-$0.0198 and the PDP slope of $-$0.0198 agree exactly; for the random forest on the same variable they are $-$0.0161 and $-$0.0062. On the payment ratio, the hardest nuisance problem in the design, the gap is largest and always towards zero: $-$0.265 against $-$0.044 for logistic regression, $-$0.489 against $-$0.100 for the random forest, $-$0.584 against $-$0.095 for XGBoost, $-$0.194 against $-$0.017 for the neural network. Months delayed runs the other way, with PDP slopes of 0.111, 0.046, 0.040 and 0.065 against debiased AMEs of 0.030, 0.013, 0.019 and 0.025, which is what one expects when an uncorrected estimator distributes a shared effect between two correlated features differently than the corrected one does. Taken across all three applications, the attenuation is the dominant pattern rather than a universal one: restricting attention to the effects the debiased estimator resolves away from zero, the partial dependence slope is smaller in absolute value than the AME in 62 of 84 cases, with a median ratio of 0.47 on Ames Housing, 0.60 here and 0.95 on Adult Income. It is most pronounced for the random forest, the learner whose regularization bias the simulations isolate.

	\section{UCI Adult Income}
	\label{app:adult}
	
	The first dataset is the UCI Adult Income dataset \citep{UCI, kohavi1996}, a widely used benchmark for binary classification in the social sciences. We combine the training and test sets for a total of 48,842 observations. The outcome is whether an individual earns more than \$50,000 per year. Features are organized into three panels: demographics (age, sex, years of education), work characteristics (hours worked per week, occupation group, work class), and financial variables (capital gains, capital losses, marital status). The specification search considers logistic regression and XGBoost across four specifications, allowing the method to be evaluated both where the parametric model is approximately correctly specified and where the black-box learner is substantially better fitted. Years of education takes the floored step size $h_j = 0.5$, so its effect is the one-year contrast; age and hours per week have enough spread that the adaptive rule already exceeds the floor, at $h_j = 0.686$ and $0.620$. Trimming is negligible for every feature except capital gains and capital losses, which lose 92.3\% and 95.3\% of the sample for the reason discussed below.
	
	\begin{table}[htbp]
		\centering
		\small
		
		\resizebox{\textwidth}{!}{%
		\begin{threeparttable}
			\caption{UCI Adult Income: Specification Search. Average marginal effects on P(Income $>$ \$50k).}
			\label{tab:adult_spec_search}
			
			\begin{tabular}{lrrrrrrrr}
				\toprule
				& \multicolumn{4}{c}{Logistic} & \multicolumn{4}{c}{XGBoost} \\
				\cmidrule(lr){2-5} \cmidrule(lr){6-9}
				& Spec A & A+B & A+C & A+B+C & Spec A & A+B & A+C & A+B+C \\
				\midrule
				\multicolumn{9}{l}{\textit{Panel A: Demographics}} \\
				\quad age & 0.0084$^{***}$ & 0.0081$^{***}$ & 0.0041$^{***}$ & 0.0043$^{***}$ & 0.0059$^{***}$ & 0.0059$^{***}$ & 0.0031$^{***}$ & 0.0033$^{***}$ \\
				& (0.0001) & (0.0001) & (0.0001) & (0.0001) & (0.0001) & (0.0001) & (0.0001) & (0.0001) \\
				\quad female & -0.1749$^{***}$ & -0.1670$^{***}$ & -0.0122$^{***}$ & -0.0153$^{**}$ & -0.1604$^{***}$ & -0.1507$^{***}$ & -0.0089$^{**}$ & -0.0128$^{**}$ \\
				& (0.0034) & (0.0045) & (0.0046) & (0.0061) & (0.0033) & (0.0041) & (0.0042) & (0.0056) \\
				\quad education\_num & 0.0464$^{***}$ & 0.0351$^{***}$ & 0.0371$^{***}$ & 0.0284$^{***}$ & 0.0606$^{***}$ & 0.0458$^{***}$ & 0.0417$^{***}$ & 0.0325$^{***}$ \\
				& (0.0007) & (0.0008) & (0.0006) & (0.0007) & (0.0007) & (0.0008) & (0.0006) & (0.0007) \\
				\midrule
				\multicolumn{9}{l}{\textit{Panel B: Work Characteristics}} \\
				\quad hours\_per\_week & -- & 0.0035$^{***}$ & -- & 0.0027$^{***}$ & -- & 0.0084$^{***}$ & -- & 0.0038$^{***}$ \\
				&  & (0.0002) &  & (0.0001) &  & (0.0002) &  & (0.0001) \\
				\quad government & -- & 0.0059 & -- & -0.0004 & -- & -0.0072 & -- & -0.0060 \\
				&  & (0.0067) &  & (0.0061) &  & (0.0066) &  & (0.0055) \\
				\quad self\_employed & -- & 0.0218$^{***}$ & -- & -0.0085$^{**}$ & -- & 0.0192$^{***}$ & -- & -0.0080$^{**}$ \\
				&  & (0.0058) &  & (0.0043) &  & (0.0056) &  & (0.0040) \\
				\quad white\_collar & -- & 0.0649$^{***}$ & -- & 0.0561$^{***}$ & -- & 0.0480$^{***}$ & -- & 0.0478$^{***}$ \\
				&  & (0.0053) &  & (0.0049) &  & (0.0051) &  & (0.0045) \\
				\quad blue\_collar & -- & -0.0246$^{***}$ & -- & -0.0248$^{***}$ & -- & -0.0482$^{***}$ & -- & -0.0300$^{***}$ \\
				&  & (0.0036) &  & (0.0031) &  & (0.0034) &  & (0.0028) \\
				\quad service & -- & -0.0744$^{***}$ & -- & -0.0456$^{***}$ & -- & -0.0795$^{***}$ & -- & -0.0342$^{***}$ \\
				&  & (0.0024) &  & (0.0023) &  & (0.0024) &  & (0.0020) \\
				\midrule
				\multicolumn{9}{l}{\textit{Panel C: Financial}} \\
				\quad capital\_gain & -- & -- & 0.0000$^{***}$ & 0.0000$^{***}$ & -- & -- & 0.0000$^{***}$ & 0.0000$^{***}$ \\
				&  &  & (0.0000) & (0.0000) &  &  & (0.0000) & (0.0000) \\
				\quad capital\_loss & -- & -- & 0.0000$^{***}$ & 0.0000$^{***}$ & -- & -- & -0.0001$^{***}$ & -0.0000$^{***}$ \\
				&  &  & (0.0000) & (0.0000) &  &  & (0.0000) & (0.0000) \\
				\quad married & -- & -- & 0.2842$^{***}$ & 0.2785$^{***}$ & -- & -- & 0.2442$^{***}$ & 0.2423$^{***}$ \\
				&  &  & (0.0048) & (0.0054) &  &  & (0.0045) & (0.0047) \\
				\midrule
				McFadden $R^2$ & 0.199 & 0.233 & 0.376 & 0.400 & 0.257 & 0.296 & 0.467 & 0.492 \\
				Accuracy & 0.797 & 0.806 & 0.841 & 0.847 & 0.805 & 0.818 & 0.864 & 0.872 \\
				\bottomrule
				\end{tabular}
				\begin{tablenotes}
				\footnotesize
				\item \textit{Notes:} $^{***}$p$<$0.01, $^{**}$p$<$0.05, $^{*}$p$<$0.10. Estimates are the debiased cross-fitted window AME with $K=5$ folds; standard errors from the orthogonal score in parentheses. Step sizes $h_j$ and trimmed fractions are reported in Table~\ref{tab:window_settings}.
				\end{tablenotes}
		\end{threeparttable}
		}%
	\end{table}
	
	Table~\ref{tab:adult_spec_search} presents the specification search results for the UCI Adult Income dataset. This is the largest of the three applications at 48,842 observations, and the estimates are correspondingly the most precisely determined in the paper. The demographic panel tells a consistent story across all specifications and both model classes. Age and education have positive effects on the probability of earning above \$50,000, and the female indicator is large and negative in the specifications that omit marital status, at $-$0.175 for logistic regression and $-$0.160 for XGBoost, collapsing to $-$0.012 and $-$0.009 once marital status enters in specification A+C. That collapse is the single most dramatic omitted variable pattern in the paper, and it is legible only because the AME puts every specification on the same probability scale. Education runs from 0.046 to 0.028 for logistic regression and 0.061 to 0.033 for XGBoost as controls accumulate. Marital status itself carries much the largest effect in the table, 0.278 and 0.242 in the full specification.

	The two model classes track each other closely throughout, and their standard errors are now of the same order everywhere --- 0.0061 and 0.0056 on the female indicator in the full specification, 0.0007 and 0.0007 on education. Under the refitting bootstrap the black-box standard errors were sometimes an order of magnitude smaller than the parametric ones, which the earlier draft read as an efficiency gain from the better-fitting model. That reading was mistaken. The two models estimate the same $\theta_{j,h}$ at the same sample size, so their standard errors ought to be similar, and the earlier gap measured how stable each fitting procedure was under resampling rather than how precisely either pinned down the target. The influence-function standard errors do behave as sampling variability should: they shrink with $n$, they are largest for the features carrying the least independent variation, and they do not reward a model for being reproducibly wrong.

	The financial panel illustrates the limits of what this design can say about capital gains and losses. Both variables are measured in raw dollars, so the reported effect is the response to a one-dollar change, and it displays as zero to four decimals in every cell but one --- as it did under the plug-in. The estimates are nonetheless significant at the one percent level, because the standard errors are smaller still; with nearly fifty thousand observations a substantively negligible effect can be sharply distinguished from zero, and this panel is a reminder that the stars answer a different question from the one the reader usually cares about. They also carry by far the heaviest trimming in the paper. Both are zero for the overwhelming majority of the sample and zero is the minimum of their observed support, so 92.3\% and 95.3\% of observations receive weight zero and the reported effect is identified only among filers who have nonzero gains or losses. The sign reversal between logistic regression and XGBoost that the earlier draft flagged in specification A+C does not survive the correction; both are zero to the displayed precision. Nothing here should be read as evidence about capital income, and a study interested in that question would want the variable re-expressed on a scale where a unit change is meaningful.
	
	\begin{table}[htbp]
		\centering
		\small
		\begin{threeparttable}
			\caption{UCI Adult Income: Model Comparison. Average marginal effects on P(Income $>$ \$50k), full specification (A+B+C).}
			\label{tab:adult_model_comparison}
			\begin{tabular}{lrrrr}
				\toprule
				Feature & Logistic & Random Forest & XGBoost & Neural Net \\
				\midrule
				age & 0.0043$^{***}$ & 0.0056$^{***}$ & 0.0033$^{***}$ & 0.0041$^{***}$ \\
				& (0.0001) & (0.0001) & (0.0001) & (0.0001) \\
				female & -0.0153$^{**}$ & -0.0210$^{***}$ & -0.0128$^{**}$ & -0.0166$^{***}$ \\
				& (0.0061) & (0.0058) & (0.0056) & (0.0059) \\
				education\_num & 0.0284$^{***}$ & 0.0336$^{***}$ & 0.0325$^{***}$ & 0.0284$^{***}$ \\
				& (0.0007) & (0.0007) & (0.0007) & (0.0007) \\
				\midrule
				hours\_per\_week & 0.0027$^{***}$ & 0.0077$^{***}$ & 0.0038$^{***}$ & 0.0049$^{***}$ \\
				& (0.0001) & (0.0001) & (0.0001) & (0.0001) \\
				government & -0.0004 & -0.0025 & -0.0060 & -0.0060 \\
				& (0.0061) & (0.0056) & (0.0055) & (0.0060) \\
				self\_employed & -0.0085$^{**}$ & -0.0103$^{**}$ & -0.0080$^{**}$ & -0.0071$^{*}$ \\
				& (0.0043) & (0.0041) & (0.0040) & (0.0042) \\
				white\_collar & 0.0561$^{***}$ & 0.0606$^{***}$ & 0.0478$^{***}$ & 0.0383$^{***}$ \\
				& (0.0049) & (0.0047) & (0.0045) & (0.0047) \\
				blue\_collar & -0.0248$^{***}$ & -0.0316$^{***}$ & -0.0300$^{***}$ & -0.0533$^{***}$ \\
				& (0.0031) & (0.0030) & (0.0028) & (0.0031) \\
				service & -0.0456$^{***}$ & -0.0262$^{***}$ & -0.0342$^{***}$ & -0.0758$^{***}$ \\
				& (0.0023) & (0.0021) & (0.0020) & (0.0023) \\
				\midrule
				capital\_gain & 0.0000$^{***}$ & 0.0000$^{***}$ & 0.0000$^{***}$ & 0.0000$^{***}$ \\
				& (0.0000) & (0.0000) & (0.0000) & (0.0000) \\
				capital\_loss & 0.0000$^{***}$ & -0.0000$^{***}$ & -0.0000$^{***}$ & 0.0000$^{***}$ \\
				& (0.0000) & (0.0000) & (0.0000) & (0.0000) \\
				married & 0.2785$^{***}$ & 0.2504$^{***}$ & 0.2423$^{***}$ & 0.2438$^{***}$ \\
				& (0.0054) & (0.0048) & (0.0047) & (0.0048) \\
				\midrule
				McFadden $R^2$ & 0.400 & 0.450 & 0.492 & 0.449 \\
				Accuracy & 0.847 & 0.863 & 0.872 & 0.858 \\
				\bottomrule
				\end{tabular}
				\begin{tablenotes}
				\footnotesize
				\item \textit{Notes:} $^{***}$p$<$0.01, $^{**}$p$<$0.05, $^{*}$p$<$0.10. Estimates are the debiased cross-fitted window AME with $K=5$ folds; standard errors from the orthogonal score in parentheses. Step sizes $h_j$ and trimmed fractions are reported in Table~\ref{tab:window_settings}.
				\end{tablenotes}
		\end{threeparttable}
	\end{table}
	
	Table~\ref{tab:adult_model_comparison} compares all four model classes on the full specification, and with nearly fifty thousand observations it is the cleanest test in the paper of whether the correction delivers a common estimand. It does. The education AME is 0.028, 0.034, 0.033 and 0.028 for logistic regression, random forest, XGBoost and the neural network, and marital status is 0.279, 0.250, 0.242 and 0.244 --- a spread of 0.005 and 0.036 respectively across four architectures with McFadden $R^2$ ranging from 0.400 to 0.492. The female indicator, which the plug-in spread from $-$0.027 down to $-$0.009 and left insignificant for the neural network, is $-$0.015, $-$0.021, $-$0.013 and $-$0.017 after correction, all four significant and all four within a standard error or two of one another. The variable the earlier draft singled out as sensitive to model specification turns out not to have been.

	The standard errors tell the same story from the other side. For a common target at a common sample size, four estimators using different nuisances should report standard errors of similar magnitude, and here they do: 0.0056 to 0.0061 on the female indicator, 0.0047 to 0.0054 on marital status, 0.0020 to 0.0023 on the service indicator. Under the plug-in they did not. The random forest reported 0.0001 on the female indicator and 0.0006 on marital status, one to two orders of magnitude below every other model, which the earlier draft attributed to ensemble averaging producing a smoother prediction surface and therefore more stable finite differences across bootstrap replicates. That description was accurate as a description of the bootstrap and wrong as a statement about precision: the replicates agreed with each other because every one of them was regularized the same way, not because the quantity was well determined. The convergence of the four standard-error columns to a common magnitude is what the orthogonal score is supposed to produce, and it is the more convincing for happening at the sample size where the asymptotics have the best claim to apply.
	
	\begin{sidewaystable}[htbp]
		\centering
		\small
		\begin{threeparttable}
			\caption{UCI Adult Income: Method Comparison. AME, SHAP, and PDP estimates for P(Income $>$ \$50k), full specification (A+B+C).}
			\label{tab:adult_method_comparison}
			\begin{tabular}{lrrrrrrrrrrrr}
				\toprule
				& \multicolumn{3}{c}{Logistic} & \multicolumn{3}{c}{Random Forest} & \multicolumn{3}{c}{XGBoost} & \multicolumn{3}{c}{Neural Net} \\
				\cmidrule(lr){2-4} \cmidrule(lr){5-7} \cmidrule(lr){8-10} \cmidrule(lr){11-13}
				& AME & SHAP & PDP & AME & SHAP & PDP & AME & SHAP & PDP & AME & SHAP & PDP \\
				\midrule
				age & 0.0043 & -0.0270 & 0.0028 & 0.0056 & -0.0068 & 0.0024 & 0.0033 & -0.0097 & 0.0036 & 0.0041 & 0.0044 & 0.0033 \\
				female & -0.0153 & 0.0022 & -0.0270 & -0.0210 & 0.0011 & -0.0218 & -0.0128 & 0.0015 & -0.0122 & -0.0166 & -0.0001 & -0.0222 \\
				education\_num & 0.0284 & -0.1186 & 0.0269 & 0.0336 & -0.0077 & 0.0137 & 0.0325 & -0.0093 & 0.0170 & 0.0284 & -0.0051 & 0.0207 \\
				\midrule
				hours\_per\_week & 0.0027 & -0.0691 & 0.0033 & 0.0077 & -0.0054 & 0.0020 & 0.0038 & -0.0080 & 0.0028 & 0.0049 & -0.0032 & 0.0028 \\
				government & -0.0004 & -0.0009 & 0.0026 & -0.0025 & -0.0003 & 0.0051 & -0.0060 & -0.0008 & -0.0003 & -0.0060 & 0.0002 & 0.0082 \\
				self\_employed & -0.0085 & 0.0016 & -0.0265 & -0.0103 & -0.0003 & 0.0032 & -0.0080 & -0.0002 & -0.0055 & -0.0071 & -0.0008 & -0.0045 \\
				white\_collar & 0.0561 & 0.0127 & 0.0538 & 0.0606 & 0.0031 & 0.0600 & 0.0478 & 0.0025 & 0.0474 & 0.0383 & 0.0007 & 0.0312 \\
				blue\_collar & -0.0248 & 0.0136 & -0.0273 & -0.0316 & 0.0010 & -0.0240 & -0.0300 & 0.0004 & -0.0251 & -0.0533 & -0.0008 & -0.0531 \\
				service & -0.0456 & -0.0027 & -0.0427 & -0.0262 & 0.0003 & -0.0176 & -0.0342 & 0.0005 & -0.0343 & -0.0758 & 0.0011 & -0.0558 \\
				\midrule
				capital\_gain & 0.0000 & -0.2627 & 0.0000 & 0.0000 & -0.0224 & 0.0000 & 0.0000 & -0.0264 & -0.0001 & 0.0000 & 0.0042 & 0.0000 \\
				capital\_loss & 0.0000 & 0.0027 & -- & -0.0000 & -0.0034 & -- & -0.0000 & -0.0062 & -- & 0.0000 & -0.0115 & -- \\
				married & 0.2785 & -0.0768 & 0.2749 & 0.2504 & -0.0064 & 0.2446 & 0.2423 & -0.0049 & 0.2357 & 0.2438 & 0.0149 & 0.1958 \\
				\bottomrule
				\end{tabular}
				\begin{tablenotes}
				\footnotesize
				\item \textit{Notes:} AME = Average Marginal Effect (marginfx). SHAP = mean signed SHAP value (GradientExplainer). PDP = slope of partial dependence plot. Full specification (A+B+C) only. No standard errors reported for SHAP and PDP.
				\end{tablenotes}
		\end{threeparttable}
	\end{sidewaystable}
	
	Table~\ref{tab:adult_method_comparison} compares AMEs against SHAP values and partial dependence plot slopes across all four models. The SHAP column measures something else and is reported here to make the difference concrete rather than to adjudicate between the two. Mean signed SHAP values are on a different scale entirely and in several cases carry the opposite sign: for logistic regression the SHAP value for education is $-$0.119 against an AME of 0.028, and for marital status it is $-$0.077 against an AME of 0.279 --- opposite in sign on the two largest effects in the application. The SHAP value for age is negative for three of the four models while every AME is positive. The most striking entry is capital gains, whose SHAP value of $-$0.263 for logistic regression is the largest SHAP magnitude anywhere in the table, attached to a variable whose AME is zero to four decimals and 92\% of whose observations are trimmed. None of this indicates an error in SHAP or in its implementation. A Shapley value is an additive attribution of a prediction relative to a reference distribution, and its average over the sample is not an estimate of the average rate of change of the prediction in a feature. The two quantities answer different questions and there is no reason to expect them to agree in magnitude or in sign. The point of the comparison is interpretive rather than evaluative: SHAP summaries are frequently read as though larger values meant larger effects on the outcome, and that reading is not supported by what the quantity measures.

	The partial dependence column behaves as it did in the other two applications: it agrees with the debiased AME where the underlying plug-in is nearly unbiased and falls short of it elsewhere. For logistic regression, the best-specified model here, agreement is close on the large effects --- marital status 0.279 against 0.275, the white collar indicator 0.056 against 0.054, service $-$0.046 against $-$0.043, education 0.028 against 0.027. For the two tree ensembles the same features separate substantially: the random forest education AME of 0.034 sits against a PDP slope of 0.014, and its hours-per-week AME of 0.008 against a slope of 0.002. A piecewise constant fit has little local slope to report between splits, and the partial dependence average inherits that flatness while the orthogonal correction does not. Partial dependence slopes for capital losses are unavailable for every model, a limitation of the grid construction when almost all mass sits at a single value. As with SHAP, no standard errors accompany either column, so a researcher working from them has point estimates of unknown precision and no way to test whether any variable has a statistically distinguishable effect.
	
	\section{Ames Housing}
	\label{app:ames}
	
	The third dataset is the Ames Housing dataset \citep{decock2011}, comprising approximately 1,460 observations of residential home sales in Ames, Iowa between 2006 and 2010. The outcome is sale price in dollars, so average marginal effects are interpreted in dollar units --- a one-unit change in a continuous feature is associated with a $\widehat{\theta}_{j,h}$ change in predicted sale price, on average across the untrimmed sample. Features are organized into physical characteristics (above-ground living area, number of bedrooms, number of bathrooms), quality and condition (overall quality rating, condition rating, year built, year remodeled), and lot and location (lot area, garage area, neighborhood tier indicators). The specification search considers linear regression and a random forest, covering the regression case and completing the sweep across outcome types. Six of the ten continuous features here are integer counts or ratings and so take the floored step size $h_j = 0.5$, which makes their reported effects one-unit contrasts: one more bedroom, one more point of overall quality. One feature requires a narrower reading than the rest. Half bathrooms take only the values zero, one and two, and both outer levels lie on the boundary of the observed support, so 63.4\% of the sample is trimmed and the half-bath effect is identified only among homes that have exactly one. The remaining features in this application are trimmed at 15\% or below.
	
	\begin{table}[htbp]
		\centering
		\small
		\resizebox{\textwidth}{!}{%
			\begin{threeparttable}
				\caption{Ames Housing: Specification Search. Average marginal effects on Sale Price (USD).}
				\label{tab:ames_housing_spec_search}
				\begin{tabular}{lrrrrrrrr}
					\toprule
					& \multicolumn{4}{c}{Linear} & \multicolumn{4}{c}{Random Forest} \\
					\cmidrule(lr){2-5} \cmidrule(lr){6-9}
					& Spec A & A+B & A+C & A+B+C & Spec A & A+B & A+C & A+B+C \\
					\midrule
					\multicolumn{9}{l}{\textit{Panel A: Physical Characteristics}} \\
					\quad GrLivArea & 138.96$^{***}$ & 106.16$^{***}$ & 84.82$^{***}$ & 68.05$^{***}$ & 128.77$^{***}$ & 103.42$^{***}$ & 90.19$^{***}$ & 81.50$^{***}$ \\
					& (20.51) & (15.58) & (14.25) & (12.14) & (13.06) & (8.98) & (8.97) & (8.77) \\
					\quad BedroomAbvGr & -33209.66$^{***}$ & -14046.40$^{*}$ & -24282.55$^{***}$ & -26454.60$^{*}$ & -28632.01$^{***}$ & -11415.18$^{***}$ & -11292.79$^{**}$ & -9212.02 \\
					& (5594.51) & (7435.20) & (9027.62) & (14881.54) & (3382.73) & (3920.43) & (5576.82) & (8925.07) \\
					\quad FullBath & 18589.00$^{***}$ & -10533.55$^{**}$ & 7261.14 & -10578.42 & 16236.01$^{***}$ & -6616.40$^{*}$ & 5652.46 & 7858.42 \\
					& (5990.54) & (4912.27) & (4987.91) & (9969.12) & (3779.18) & (3670.56) & (3754.82) & (8988.79) \\
					\quad HalfBath & -11228.94$^{**}$ & -8180.81$^{**}$ & -9092.88$^{**}$ & 1028.31 & -7189.23$^{*}$ & -6638.05$^{***}$ & -11578.88$^{**}$ & -10317.96$^{**}$ \\
					& (4624.11) & (3242.49) & (4451.09) & (5697.41) & (4264.04) & (2525.05) & (4601.22) & (4713.92) \\
					\midrule
					\multicolumn{9}{l}{\textit{Panel B: Quality and Condition}} \\
					\quad OverallQual & -- & 15819.18$^{***}$ & -- & 5436.28 & -- & 11911.79$^{***}$ & -- & 16255.82$^{***}$ \\
					&  & (1589.02) &  & (4964.23) &  & (1342.87) &  & (4829.24) \\
					\quad OverallCond & -- & 6477.48$^{*}$ & -- & 2890.60 & -- & 8281.97$^{***}$ & -- & 9190.11$^{***}$ \\
					&  & (3687.03) &  & (5496.46) &  & (2434.71) &  & (3492.01) \\
					\quad YearBuilt & -- & 443.39 & -- & -141.08 & -- & 1033.46$^{***}$ & -- & 1323.39$^{***}$ \\
					&  & (277.71) &  & (404.21) &  & (190.55) &  & (419.89) \\
					\quad YearRemodAdd & -- & 192.69 & -- & 1743.45$^{**}$ & -- & 556.20$^{***}$ & -- & -311.44 \\
					&  & (161.82) &  & (740.09) &  & (173.98) &  & (753.53) \\
					\midrule
					\multicolumn{9}{l}{\textit{Panel C: Lot and Location}} \\
					\quad LotArea & -- & -- & 5.51 & 23.02$^{**}$ & -- & -- & -7.53 & -12.51 \\
					&  &  & (3.93) & (11.37) &  &  & (8.06) & (11.14) \\
					\quad GarageArea & -- & -- & 40.08$^{*}$ & -102.78 & -- & -- & 103.31$^{***}$ & 90.33 \\
					&  &  & (20.94) & (71.57) &  &  & (34.64) & (59.74) \\
					\quad nbhd\_high & -- & -- & 128801.76$^{***}$ & 89812.31$^{**}$ & -- & -- & 66978.55$^{***}$ & 43512.42$^{*}$ \\
					&  &  & (34103.92) & (38393.43) &  &  & (16110.93) & (24635.87) \\
					\quad nbhd\_mid & -- & -- & 77339.38$^{**}$ & 74336.31$^{*}$ & -- & -- & 42305.81$^{***}$ & 33767.43 \\
					&  &  & (33663.73) & (39182.81) &  &  & (15387.40) & (26676.67) \\
					\midrule
					$R^2$ & 0.584 & 0.758 & 0.739 & 0.802 & 0.829 & 0.937 & 0.926 & 0.953 \\
					RMSE & 51,243.9 & 39,080.1 & 40,596.4 & 35,355.5 & 32,830.9 & 19,881.6 & 21,534.2 & 17,244.4 \\
					\bottomrule
					\end{tabular}
					\begin{tablenotes}
					\footnotesize
					\item \textit{Notes:} $^{***}$p$<$0.01, $^{**}$p$<$0.05, $^{*}$p$<$0.10. Estimates are the debiased cross-fitted window AME with $K=5$ folds; standard errors from the orthogonal score in parentheses. Step sizes $h_j$ and trimmed fractions are reported in Table~\ref{tab:window_settings}.
					\end{tablenotes}
			\end{threeparttable}
		}%
	\end{table}
	
	Table~\ref{tab:ames_housing_spec_search} presents the specification search. The regression setting allows AMEs to be interpreted directly in dollar units, which makes the substantive findings immediately legible. Above-ground living area carries an AME of \$139 per square foot in the physical-characteristics specification for linear regression, falling to \$106, \$85 and finally \$68 as quality, condition, and location controls are added --- a pattern of positive omitted variable bias that makes intuitive sense, since larger homes tend to be in better neighborhoods and of higher quality. The random forest traces the same path from \$129 to \$81. That the two model classes track each other this closely across the sequence is the first thing to note about the debiased estimates: because the orthogonal correction removes the first-order dependence on the fitted regression, the two columns are estimating the same window effect rather than reporting the slope each model happens to believe, and the linear column is no longer simply the ordinary least squares coefficient. The standard errors are now of comparable size across the two model classes --- \$12 for linear regression and \$9 for the random forest on GrLivArea in the full specification --- rather than differing by the order of magnitude that the refitting bootstrap produced, where the ensemble's smooth prediction surface made every replicate agree with every other and drove the reported standard error to \$1.93. Those intervals were narrow around a biased centre; the influence-function intervals are wider and honest.

	The cost of that honesty is visible in the quality panel. The overall quality rating is estimated at \$15,819 per point for linear regression and \$11,912 for the random forest in specification A+B, both precisely determined, but in the full specification the linear estimate falls to \$5,436 with a standard error of \$4,964 and is no longer distinguishable from zero, while the random forest gives \$16,256 with a standard error of \$4,829. With $n = 1{,}460$ observations, twelve features and a second-order sieve for the representer, the full specification is close to the limit of what this sample can support: the estimator reports that limit rather than concealing it. The neighbourhood tier indicators show the same attenuation as before --- the high neighbourhood effect falls from \$128,802 to \$89,812 for linear regression and from \$66,979 to \$43,512 for the random forest as quality controls are added --- but with standard errors near \$25{,}000 to \$38{,}000, the attenuation is a pattern in the point estimates rather than a difference the data can resolve.
	
	\begin{table}[htbp]
		\centering
		\small
		\begin{threeparttable}
			\caption{Ames Housing: Model Comparison. Average marginal effects on Sale Price (USD), full specification (A+B+C).}
			\label{tab:ames_housing_model_comparison}
			\begin{tabular}{lrrrr}
				\toprule
				Feature & Linear & Random Forest & XGBoost & Neural Net \\
				\midrule
				GrLivArea & 68.05$^{***}$ & 81.50$^{***}$ & 70.27$^{***}$ & 66.20$^{***}$ \\
				& (12.14) & (8.77) & (6.60) & (6.24) \\
				BedroomAbvGr & -26454.60$^{*}$ & -9212.02 & -3464.31 & -12638.60$^{***}$ \\
				& (14881.54) & (8925.07) & (4508.19) & (3830.14) \\
				FullBath & -10578.42 & 7858.42 & 4761.27 & -5055.91 \\
				& (9969.12) & (8988.79) & (7661.72) & (5080.88) \\
				HalfBath & 1028.31 & -10317.96$^{**}$ & -10906.00$^{***}$ & -2175.44 \\
				& (5697.41) & (4713.92) & (3940.81) & (3286.03) \\
				\midrule
				OverallQual & 5436.28 & 16255.82$^{***}$ & 14913.45$^{***}$ & 8524.39$^{***}$ \\
				& (4964.23) & (4829.24) & (4271.98) & (2394.41) \\
				OverallCond & 2890.60 & 9190.11$^{***}$ & 12102.59$^{***}$ & 7234.42$^{**}$ \\
				& (5496.46) & (3492.01) & (2519.57) & (2920.49) \\
				YearBuilt & -141.08 & 1323.39$^{***}$ & 1352.73$^{***}$ & 230.04 \\
				& (404.21) & (419.89) & (366.04) & (267.61) \\
				YearRemodAdd & 1743.45$^{**}$ & -311.44 & -289.89 & 942.55$^{***}$ \\
				& (740.09) & (753.53) & (673.39) & (351.83) \\
				\midrule
				LotArea & 23.02$^{**}$ & -12.51 & -10.83 & 10.04$^{**}$ \\
				& (11.37) & (11.14) & (9.91) & (4.81) \\
				GarageArea & -102.78 & 90.33 & 92.29$^{*}$ & -6.11 \\
				& (71.57) & (59.74) & (49.36) & (27.32) \\
				nbhd\_high & 89812.31$^{**}$ & 43512.42$^{*}$ & 33476.93$^{**}$ & 35458.39$^{**}$ \\
				& (38393.43) & (24635.87) & (16103.49) & (17705.63) \\
				nbhd\_mid & 74336.31$^{*}$ & 33767.43 & 17468.12 & 36951.86$^{**}$ \\
				& (39182.81) & (26676.67) & (18286.73) & (17998.56) \\
				\midrule
				$R^2$ & 0.802 & 0.953 & 0.954 & 0.932 \\
				RMSE & 35,355.5 & 17,244.4 & 17,036.6 & 20,699.9 \\
				\bottomrule
				\end{tabular}
				\begin{tablenotes}
				\footnotesize
				\item \textit{Notes:} $^{***}$p$<$0.01, $^{**}$p$<$0.05, $^{*}$p$<$0.10. Estimates are the debiased cross-fitted window AME with $K=5$ folds; standard errors from the orthogonal score in parentheses. Step sizes $h_j$ and trimmed fractions are reported in Table~\ref{tab:window_settings}.
				\end{tablenotes}
		\end{threeparttable}
	\end{table}
	
	Table~\ref{tab:ames_housing_model_comparison} compares all four models on the full specification. All four are well fitted, with $R^2$ ranging from 0.802 for linear regression to 0.954 for XGBoost, the random forest at 0.953 and the neural network at 0.932 with normalization. The GrLivArea AME is \$68, \$81, \$70, and \$66 for linear regression, random forest, XGBoost, and neural network respectively --- a range of \$15 across four models that differ fundamentally in their functional form, and all four intervals overlap comfortably. Under the debiased estimator this agreement means something different from what it meant under the plug-in. Every column is now an estimate of the same $\theta_{j,h}$, so agreement is evidence that the four nuisance estimates are all good enough for the correction to do its work, not evidence that four different models happen to have similar slopes. Disagreement, correspondingly, is now diagnostic of nuisance quality rather than of functional form.

	The most consequential change from the plug-in results is in the standard errors, and it runs in the opposite direction to the one the refitting bootstrap suggested. Under that bootstrap the neural network's standard errors were implausibly tight --- on the order of single dollars for several variables, \$7 for bedrooms and \$8 for overall quality --- because resampling propagated uncertainty through the model fitting but the smooth differentiable surface produced nearly identical finite differences across replicates. The influence-function standard errors for the same quantities are \$3,830 and \$2,394. The random forest shows the same correction: the high-neighborhood standard error rises from \$256 to \$24,636. Those earlier figures were not precision but the bootstrap's inability to see a bias that every replicate shared, which is the failure mode of Remark~\ref{M-rem:refit}, and their disappearance is the clearest single indication in the empirical results that the orthogonal score is doing what it was constructed to do.

	Read with those intervals attached, the bedroom variable is more equivocal than the plug-in made it look. The AME is $-$\$26,455 for linear regression, $-$\$9,212 for the random forest, $-$\$3,464 for XGBoost, and $-$\$12,639 for the neural network. All four are negative, consistent with the well-known hedonic finding that holding square footage fixed, more bedrooms means smaller rooms, but only the neural network estimate is significant at conventional levels and the linear estimate only marginally so. The plug-in reported all four as significant with standard errors between \$7 and \$3,042. The honest reading is that the sign is stable across architectures while this sample cannot pin down the magnitude, which is a weaker conclusion than the paper could previously state and a better supported one.
	
	\begin{sidewaystable}[htbp]
		\centering
		\small
		\begin{threeparttable}
			\caption{Ames Housing: Method Comparison. AME, SHAP, and PDP estimates for Sale Price (USD), full specification (A+B+C).}
			\label{tab:ames_housing_method_comparison}
			\begin{tabular}{lrrrrrrrrrrrr}
				\toprule
				& \multicolumn{3}{c}{Linear} & \multicolumn{3}{c}{Random Forest} & \multicolumn{3}{c}{XGBoost} & \multicolumn{3}{c}{Neural Net} \\
				\cmidrule(lr){2-4} \cmidrule(lr){5-7} \cmidrule(lr){8-10} \cmidrule(lr){11-13}
				& AME & SHAP & PDP & AME & SHAP & PDP & AME & SHAP & PDP & AME & SHAP & PDP \\
				\midrule
				GrLivArea & 68.05 & 3423.49 & 65.57 & 81.50 & 61.69 & 38.64 & 70.27 & -971.88 & 39.06 & 66.20 & -- & 58.62 \\
				BedroomAbvGr & -26454.60 & 145.82 & -6189.02 & -9212.02 & 90.87 & -1272.27 & -3464.31 & 158.44 & -4456.61 & -12638.60 & -- & -9294.43 \\
				FullBath & -10578.42 & 31.01 & -6287.42 & 7858.42 & 34.06 & 543.63 & 4761.27 & 54.86 & 193.04 & -5055.91 & -- & 2427.62 \\
				HalfBath & 1028.31 & -298.76 & -6967.90 & -10317.96 & 50.46 & -125.75 & -10906.00 & 30.53 & -1495.20 & -2175.44 & -- & -2723.69 \\
				\midrule
				OverallQual & 5436.28 & 1635.24 & 16465.22 & 16255.82 & 596.70 & 20776.14 & 14913.45 & 713.13 & 20130.92 & 8524.39 & -- & 10802.73 \\
				OverallCond & 2890.60 & 980.44 & 6311.48 & 9190.11 & 161.26 & 2030.44 & 12102.59 & -481.60 & 5071.47 & 7234.42 & -- & 6365.50 \\
				YearBuilt & -141.08 & -415.01 & 440.47 & 1323.39 & 288.29 & 209.05 & 1352.73 & 569.00 & 397.33 & 230.04 & -- & 421.27 \\
				YearRemodAdd & 1743.45 & 86.86 & 101.51 & -311.44 & 373.04 & 118.06 & -289.89 & 499.84 & 152.00 & 942.55 & -- & 561.16 \\
				\midrule
				LotArea & 23.02 & 267.44 & 0.61 & -12.51 & -183.50 & 1.33 & -10.83 & 216.28 & 2.38 & 10.04 & -- & 2.35 \\
				GarageArea & -102.78 & 93.22 & 38.68 & 90.33 & -884.31 & 23.77 & 92.29 & -704.59 & 28.98 & -6.11 & -- & 41.48 \\
				nbhd\_high & 89812.31 & 77.50 & 51433.02 & 43512.42 & -506.36 & 5058.64 & 33476.93 & 50.70 & 20307.34 & 35458.39 & -- & 12399.98 \\
				nbhd\_mid & 74336.31 & 266.38 & 7508.07 & 33767.43 & 40.68 & 5949.97 & 17468.12 & -134.71 & 9744.97 & 36951.86 & -- & 11392.44 \\
				\bottomrule
				\end{tabular}
				\begin{tablenotes}
				\footnotesize
				\item \textit{Notes:} AME = Average Marginal Effect (marginfx). SHAP = mean signed SHAP value (GradientExplainer). PDP = slope of partial dependence plot. Full specification (A+B+C) only. No standard errors reported for SHAP and PDP.
				\end{tablenotes}
		\end{threeparttable}
	\end{sidewaystable}
	
	Table~\ref{tab:ames_housing_method_comparison} makes the difference between the two quantities easiest to see, since the units are dollars and the magnitudes are large enough to be concrete. SHAP values are unavailable for the neural network in this table, reflecting a limitation of the GradientExplainer implementation for this architecture. For the other three models the two columns are far apart. The SHAP value for GrLivArea is \$3,423 for the linear model against an AME of \$68, roughly fifty times larger --- which is what one should expect, since the SHAP figure is the average absolute displacement a typical house's living area contributes to its predicted price, not the price change associated with one additional square foot. The SHAP values for garage area ($-$\$884), lot area ($-$\$183), and the high neighborhood indicator ($-$\$506) are negative for the random forest while the AME estimates are positive for all three. And as before, no standard errors attach to the SHAP column, so a reader who did want to treat those figures as effects would have no way to gauge their precision.

	The partial dependence column requires a different reading than it did under the plug-in, and the change is itself informative. A partial dependence slope is computed by perturbing a single full-sample fit and averaging, which makes it a plug-in average of exactly the kind the orthogonal correction is designed to repair; it carries no correction term and no standard error. Where the plug-in is nearly unbiased the two therefore continue to agree closely: on above-ground living area, the best supported and smoothest feature in the data, the AME is \$68 and the PDP slope \$66 for linear regression, \$66 and \$59 for the neural network. Where the regularization bias bites, they separate, and they separate in the direction the theory predicts --- towards zero for the uncorrected quantity. The bedroom AME of $-$\$26,455 for linear regression sits against a PDP slope of $-$\$6,189; the high-neighbourhood AME of \$43,512 for the random forest sits against a PDP slope of \$5,059, and the mid-neighbourhood AME of \$33,767 against \$5,950. These are the count and binary contrasts, where the fitted surface is flattest between levels and where the plug-in has the most to lose. The pattern that once read as mutual corroboration between two estimators now reads as a measurement of the bias separating them, which is the same quantity Simulation~2 isolates in the coverage of the refitting bootstrap. The Ames Housing results complete the sweep across outcome types: the framework produces dollar-denominated marginal effects with valid standard errors for a continuous regression outcome, using exactly the same estimation procedure as the binary classification examples, with no modification required.
	
	\section{Detailed simulation results}
	\label{app:sim_detail}

	This section reports the three simulation studies in full: every
	sample size, every data generating process, every feature, and the
	sieve-representer and bounded-support arms of Simulation 3. The
	design, the learners and their hyperparameters are described in
	Section~\ref{M-sec:simulation} of the main paper, and
	Table~\ref{M-tab:sim_summary} there summarizes the results that
	follow.

	\subsection{Simulation 1: Recovering the Window Effect}
	
	\begin{table}[htbp]
		\centering
		\small
		\begin{threeparttable}
			\caption{Simulation 1: AME Recovery --- Linear DGP (Classification). Mean bias and RMSE (in parentheses) across 500 Monte Carlo iterations.}
			\label{tab:sim1_classification_linear}
			\begin{tabular}{lrrrrrr}
				\toprule
				& & \multicolumn{5}{c}{Bias (RMSE)} \\
				\cmidrule(lr){3-7}
				Variable & True AME & $n=250$ & $n=500$ & $n=1,000$ & $n=2,500$ & $n=5,000$ \\
				\midrule
				\multicolumn{7}{l}{\textit{Panel: Logistic}} \\
				\quad x1 & 0.199 & -0.0058 & -0.0034 & -0.0014 & -0.0008 & -0.0006 \\
				&  & (0.0201) & (0.0146) & (0.0106) & (0.0064) & (0.0048) \\
				\quad x2 & 0.298 & -0.0109 & -0.0057 & -0.0026 & -0.0010 & -0.0005 \\
				&  & (0.0229) & (0.0153) & (0.0105) & (0.0066) & (0.0046) \\
				\quad x3 & 0.000 & 0.0012 & 0.0006 & 0.0006 & 0.0003 & 0.0002 \\
				&  & (0.0198) & (0.0136) & (0.0099) & (0.0063) & (0.0045) \\
				\quad x4 & 0.000 & 0.0008 & -0.0001 & -0.0008 & -0.0001 & 0.0003 \\
				&  & (0.0197) & (0.0136) & (0.0100) & (0.0062) & (0.0045) \\
				\midrule
				\multicolumn{7}{l}{\textit{Panel: Random Forest}} \\
				\quad x1 & 0.199 & 0.0126 & 0.0044 & -0.0018 & -0.0111 & -0.0157 \\
				&  & (0.0412) & (0.0239) & (0.0160) & (0.0142) & (0.0168) \\
				\quad x2 & 0.298 & 0.0199 & 0.0113 & -0.0008 & -0.0102 & -0.0158 \\
				&  & (0.0500) & (0.0309) & (0.0174) & (0.0143) & (0.0172) \\
				\quad x3 & 0.000 & 0.0007 & 0.0010 & 0.0005 & 0.0002 & 0.0002 \\
				&  & (0.0284) & (0.0171) & (0.0104) & (0.0056) & (0.0034) \\
				\quad x4 & 0.000 & -0.0008 & -0.0013 & -0.0012 & -0.0001 & 0.0002 \\
				&  & (0.0267) & (0.0177) & (0.0106) & (0.0055) & (0.0034) \\
				\midrule
				\multicolumn{7}{l}{\textit{Panel: XGBoost}} \\
				\quad x1 & 0.199 & 0.0590 & 0.0265 & 0.0107 & -0.0016 & -0.0049 \\
				&  & (0.0782) & (0.0401) & (0.0223) & (0.0108) & (0.0083) \\
				\quad x2 & 0.298 & 0.0643 & 0.0315 & 0.0085 & -0.0015 & -0.0046 \\
				&  & (0.0893) & (0.0490) & (0.0218) & (0.0116) & (0.0083) \\
				\quad x3 & 0.000 & 0.0008 & -0.0002 & 0.0006 & 0.0002 & 0.0002 \\
				&  & (0.0332) & (0.0193) & (0.0110) & (0.0056) & (0.0033) \\
				\quad x4 & 0.000 & 0.0008 & -0.0009 & -0.0006 & 0.0000 & 0.0002 \\
				&  & (0.0344) & (0.0198) & (0.0116) & (0.0054) & (0.0032) \\
				\midrule
				\multicolumn{7}{l}{\textit{Panel: Neural Net}} \\
				\quad x1 & 0.199 & -0.0078 & -0.0007 & 0.0005 & 0.0003 & -0.0002 \\
				&  & (0.0222) & (0.0156) & (0.0115) & (0.0074) & (0.0058) \\
				\quad x2 & 0.298 & -0.0162 & -0.0017 & 0.0004 & 0.0004 & 0.0002 \\
				&  & (0.0276) & (0.0174) & (0.0116) & (0.0073) & (0.0052) \\
				\quad x3 & 0.000 & 0.0014 & 0.0010 & 0.0007 & 0.0003 & 0.0004 \\
				&  & (0.0204) & (0.0144) & (0.0103) & (0.0069) & (0.0057) \\
				\quad x4 & 0.000 & 0.0009 & -0.0003 & -0.0007 & -0.0001 & 0.0004 \\
				&  & (0.0203) & (0.0143) & (0.0105) & (0.0070) & (0.0055) \\
				\bottomrule
				\end{tabular}
				\begin{tablenotes}
				\footnotesize
				\item \textit{Notes:} Mean bias and RMSE (in parentheses) computed over 500 Monte Carlo iterations. True AMEs computed via Monte Carlo integration with $n=1{,}000{,}000$ observations. Noise features x3 and x4 have true AME of zero.
				\end{tablenotes}
		\end{threeparttable}
	\end{table}
	
	Table~\ref{tab:sim1_classification_linear} reports bias and RMSE for the linear classification DGP, where the true AMEs of $x_1$ and $x_2$ are 0.199 and 0.298 --- these are not 2.0 and 3.0 because the sigmoid compresses the scale. Logistic regression recovers the true AMEs with negligible bias at all sample sizes, and RMSE declines from around 0.020 at $n=250$ to 0.005 at $n=5{,}000$ at the expected $O(n^{-1/2})$ rate. The neural network behaves similarly, with small bias at $n=250$ that disappears by $n=500$. XGBoost shows larger bias at small samples --- around 0.06 for both $x_1$ and $x_2$ at $n=250$ --- reflecting the well-known tendency of gradient boosted models to underfit at small sample sizes, but bias falls sharply and RMSE converges by $n=2{,}500$. Random forests show a different pattern: bias is small at $n=250$ and $n=500$ but turns negative and persistent at $n=2{,}500$ and $n=5{,}000$, reaching around $-0.016$ for both signal features at the largest sample size. This reflects the regularization bias inherent in random forests --- the tree depth and subsampling constraints impose shrinkage that does not vanish with sample size at the default tuning settings. The noise features $x_3$ and $x_4$ have bias and RMSE near zero throughout for all model classes, confirming that the estimator does not spuriously attribute effects to irrelevant features.
	
	\begin{table}[htbp]
		\centering
		\small
		\begin{threeparttable}
			\caption{Simulation 1: AME Recovery --- Nonlinear DGP (Classification). Mean bias and RMSE (in parentheses) across 500 Monte Carlo iterations.}
			\label{tab:sim1_classification_nonlinear}
			\begin{tabular}{lrrrrrr}
				\toprule
				& & \multicolumn{5}{c}{Bias (RMSE)} \\
				\cmidrule(lr){3-7}
				Variable & True AME & $n=250$ & $n=500$ & $n=1,000$ & $n=2,500$ & $n=5,000$ \\
				\midrule
				\multicolumn{7}{l}{\textit{Panel: Logistic}} \\
				\quad x1 & 0.000 & 0.0002 & 0.0015 & 0.0009 & 0.0002 & -0.0001 \\
				&  & (0.0230) & (0.0166) & (0.0123) & (0.0078) & (0.0055) \\
				\quad x2 & 0.298 & -0.0077 & -0.0042 & -0.0017 & -0.0007 & -0.0005 \\
				&  & (0.0215) & (0.0152) & (0.0109) & (0.0065) & (0.0048) \\
				\quad x3 & 0.000 & 0.0003 & 0.0001 & -0.0002 & -0.0001 & 0.0003 \\
				&  & (0.0229) & (0.0158) & (0.0112) & (0.0074) & (0.0052) \\
				\quad x4 & 0.000 & -0.0012 & 0.0005 & 0.0006 & -0.0001 & -0.0000 \\
				&  & (0.0252) & (0.0159) & (0.0108) & (0.0075) & (0.0051) \\
				\midrule
				\multicolumn{7}{l}{\textit{Panel: Random Forest}} \\
				\quad x1 & 0.000 & 0.0014 & 0.0001 & 0.0000 & -0.0003 & -0.0000 \\
				&  & (0.0407) & (0.0256) & (0.0173) & (0.0114) & (0.0082) \\
				\quad x2 & 0.298 & 0.0177 & 0.0060 & -0.0040 & -0.0120 & -0.0155 \\
				&  & (0.0521) & (0.0278) & (0.0182) & (0.0157) & (0.0171) \\
				\quad x3 & 0.000 & 0.0019 & -0.0000 & 0.0003 & 0.0001 & 0.0003 \\
				&  & (0.0285) & (0.0166) & (0.0100) & (0.0052) & (0.0031) \\
				\quad x4 & 0.000 & -0.0026 & -0.0007 & 0.0005 & 0.0000 & -0.0001 \\
				&  & (0.0295) & (0.0165) & (0.0099) & (0.0054) & (0.0031) \\
				\midrule
				\multicolumn{7}{l}{\textit{Panel: XGBoost}} \\
				\quad x1 & 0.000 & 0.0039 & 0.0003 & 0.0004 & -0.0000 & -0.0003 \\
				&  & (0.0549) & (0.0324) & (0.0209) & (0.0123) & (0.0079) \\
				\quad x2 & 0.298 & 0.0471 & 0.0221 & 0.0066 & -0.0028 & -0.0054 \\
				&  & (0.0787) & (0.0396) & (0.0208) & (0.0114) & (0.0087) \\
				\quad x3 & 0.000 & 0.0038 & 0.0007 & 0.0007 & 0.0002 & 0.0003 \\
				&  & (0.0307) & (0.0172) & (0.0108) & (0.0051) & (0.0028) \\
				\quad x4 & 0.000 & -0.0016 & 0.0010 & 0.0006 & 0.0001 & 0.0001 \\
				&  & (0.0329) & (0.0184) & (0.0105) & (0.0052) & (0.0027) \\
				\midrule
				\multicolumn{7}{l}{\textit{Panel: Neural Net}} \\
				\quad x1 & 0.000 & 0.0006 & 0.0015 & 0.0002 & 0.0001 & 0.0000 \\
				&  & (0.0234) & (0.0189) & (0.0140) & (0.0094) & (0.0069) \\
				\quad x2 & 0.298 & -0.0132 & -0.0026 & 0.0002 & 0.0005 & 0.0001 \\
				&  & (0.0249) & (0.0155) & (0.0112) & (0.0077) & (0.0064) \\
				\quad x3 & 0.000 & 0.0001 & 0.0002 & 0.0007 & 0.0003 & 0.0003 \\
				&  & (0.0215) & (0.0138) & (0.0101) & (0.0068) & (0.0053) \\
				\quad x4 & 0.000 & -0.0007 & 0.0004 & 0.0005 & -0.0001 & 0.0002 \\
				&  & (0.0230) & (0.0144) & (0.0100) & (0.0070) & (0.0056) \\
				\bottomrule
				\end{tabular}
				\begin{tablenotes}
				\footnotesize
				\item \textit{Notes:} Mean bias and RMSE (in parentheses) computed over 500 Monte Carlo iterations. True AMEs computed via Monte Carlo integration with $n=1{,}000{,}000$ observations. Noise features x3 and x4 have true AME of zero.
				\end{tablenotes}
		\end{threeparttable}
	\end{table}
	
	Table~\ref{tab:sim1_classification_nonlinear} reports results for the nonlinear classification DGP, where $P(y=1) = \sigma(2x_1^2 + 3x_2)$. The true AME of $x_1$ is zero despite the quadratic term, because the sigmoid and quadratic effects cancel in expectation at these parameter values as confirmed by the Monte Carlo ground truth. All four model classes recover this near-zero AME accurately, with bias well below 0.005 across all sample sizes. The more important result is $x_2$, which has true AME of 0.298. Logistic regression misspecifies the DGP --- the true relationship is nonlinear in $x_1$ --- but still recovers the $x_2$ AME accurately because $x_2$ enters linearly and the misspecification in $x_1$ does not propagate to the $x_2$ marginal effect. This is an important practical result: AMEs from a misspecified parametric model can still be reliable for variables that enter the true DGP in a manner consistent with the model's functional form. The black-box models show the same pattern as the linear DGP: XGBoost has elevated RMSE at small samples, random forests develop negative bias at large samples, and neural networks perform comparably to logistic regression throughout.
	
	\begin{table}[htbp]
		\centering
		\small
		\begin{threeparttable}
			\caption{Simulation 1: AME Recovery --- Interaction DGP (Classification). Mean bias and RMSE (in parentheses) across 500 Monte Carlo iterations.}
			\label{tab:sim1_classification_interaction}
			\begin{tabular}{lrrrrrr}
				\toprule
				& & \multicolumn{5}{c}{Bias (RMSE)} \\
				\cmidrule(lr){3-7}
				Variable & True AME & $n=250$ & $n=500$ & $n=1,000$ & $n=2,500$ & $n=5,000$ \\
				\midrule
				\multicolumn{7}{l}{\textit{Panel: Logistic}} \\
				\quad x1 & 0.199 & -0.0065 & -0.0042 & -0.0021 & -0.0008 & -0.0007 \\
				&  & (0.0219) & (0.0153) & (0.0112) & (0.0071) & (0.0051) \\
				\quad x2 & 0.298 & -0.0083 & -0.0035 & -0.0018 & -0.0008 & -0.0007 \\
				&  & (0.0222) & (0.0154) & (0.0102) & (0.0071) & (0.0050) \\
				\quad x3 & 0.000 & 0.0009 & -0.0002 & 0.0005 & 0.0003 & -0.0001 \\
				&  & (0.0206) & (0.0155) & (0.0107) & (0.0070) & (0.0046) \\
				\quad x4 & 0.000 & 0.0007 & -0.0000 & -0.0001 & 0.0001 & 0.0002 \\
				&  & (0.0212) & (0.0159) & (0.0106) & (0.0066) & (0.0047) \\
				\midrule
				\multicolumn{7}{l}{\textit{Panel: Random Forest}} \\
				\quad x1 & 0.199 & 0.0114 & 0.0051 & -0.0031 & -0.0110 & -0.0154 \\
				&  & (0.0423) & (0.0264) & (0.0169) & (0.0144) & (0.0167) \\
				\quad x2 & 0.298 & 0.0178 & 0.0104 & -0.0023 & -0.0101 & -0.0140 \\
				&  & (0.0498) & (0.0320) & (0.0200) & (0.0148) & (0.0158) \\
				\quad x3 & 0.000 & 0.0019 & 0.0004 & 0.0002 & 0.0002 & -0.0000 \\
				&  & (0.0276) & (0.0170) & (0.0104) & (0.0059) & (0.0034) \\
				\quad x4 & 0.000 & -0.0005 & -0.0011 & -0.0001 & 0.0001 & -0.0000 \\
				&  & (0.0267) & (0.0177) & (0.0111) & (0.0057) & (0.0033) \\
				\midrule
				\multicolumn{7}{l}{\textit{Panel: XGBoost}} \\
				\quad x1 & 0.199 & 0.0511 & 0.0225 & 0.0079 & -0.0014 & -0.0046 \\
				&  & (0.0737) & (0.0392) & (0.0219) & (0.0105) & (0.0082) \\
				\quad x2 & 0.298 & 0.0496 & 0.0228 & 0.0042 & -0.0021 & -0.0047 \\
				&  & (0.0796) & (0.0439) & (0.0227) & (0.0125) & (0.0091) \\
				\quad x3 & 0.000 & 0.0004 & 0.0005 & 0.0004 & 0.0000 & 0.0000 \\
				&  & (0.0310) & (0.0205) & (0.0111) & (0.0064) & (0.0035) \\
				\quad x4 & 0.000 & -0.0012 & -0.0015 & -0.0004 & 0.0001 & 0.0001 \\
				&  & (0.0338) & (0.0205) & (0.0118) & (0.0060) & (0.0037) \\
				\midrule
				\multicolumn{7}{l}{\textit{Panel: Neural Net}} \\
				\quad x1 & 0.199 & -0.0089 & -0.0013 & 0.0001 & 0.0000 & -0.0004 \\
				&  & (0.0233) & (0.0158) & (0.0117) & (0.0079) & (0.0061) \\
				\quad x2 & 0.298 & -0.0114 & -0.0002 & 0.0007 & 0.0005 & -0.0002 \\
				&  & (0.0257) & (0.0170) & (0.0131) & (0.0085) & (0.0067) \\
				\quad x3 & 0.000 & 0.0008 & 0.0004 & 0.0008 & 0.0001 & 0.0000 \\
				&  & (0.0193) & (0.0148) & (0.0099) & (0.0070) & (0.0053) \\
				\quad x4 & 0.000 & 0.0005 & -0.0002 & 0.0000 & -0.0001 & 0.0006 \\
				&  & (0.0206) & (0.0151) & (0.0103) & (0.0069) & (0.0057) \\
				\bottomrule
				\end{tabular}
				\begin{tablenotes}
				\footnotesize
				\item \textit{Notes:} Mean bias and RMSE (in parentheses) computed over 500 Monte Carlo iterations. True AMEs computed via Monte Carlo integration with $n=1{,}000{,}000$ observations. Noise features x3 and x4 have true AME of zero.
				\end{tablenotes}
		\end{threeparttable}
	\end{table}
	
	Table~\ref{tab:sim1_classification_interaction} reports results for the interaction DGP, where $P(y=1) = \sigma(2x_1 + 3x_2 + 2x_1 x_2)$. The true AMEs of $x_1$ and $x_2$ remain 0.199 and 0.298 --- the same as the linear DGP --- because the interaction term $2x_1 x_2$ has zero expectation under $N(0,1)$ and the Monte Carlo ground truth confirms near-identical values. The key finding is that all four model classes recover these AMEs with bias and RMSE patterns essentially identical to the linear DGP, despite the true DGP including a multiplicative interaction. This is by design of the AME estimand: the average marginal effect averages over the joint distribution of the features, so the interaction term is absorbed into the average derivative rather than requiring explicit modelling. A logistic regression that omits the interaction term recovers the correct AME because the interaction's contribution to the marginal effect of $x_1$ is $\mathbb{E}[2x_2] = 0$ and similarly for $x_2$. This is an important robustness result --- AMEs are not sensitive to interaction misspecification when the interactions involve zero-mean features, which is the common case after centring.
	
	\begin{table}[htbp]
		\centering
		\small
		\begin{threeparttable}
			\caption{Simulation 1: AME Recovery --- Linear DGP (Regression). Mean bias and RMSE (in parentheses) across 1,000 Monte Carlo iterations.}
			\label{tab:sim1_regression_linear}
			\begin{tabular}{lrrrrrr}
				\toprule
				& & \multicolumn{5}{c}{Bias (RMSE)} \\
				\cmidrule(lr){3-7}
				Variable & True AME & $n=250$ & $n=500$ & $n=1,000$ & $n=2,500$ & $n=5,000$ \\
				\midrule
				\multicolumn{7}{l}{\textit{Panel: Linear}} \\
				\quad x1 & 2.000 & -0.0017 & -0.0001 & -0.0000 & -0.0006 & -0.0003 \\
				&  & (0.0659) & (0.0445) & (0.0322) & (0.0207) & (0.0137) \\
				\quad x2 & 3.000 & -0.0043 & -0.0017 & 0.0007 & -0.0004 & -0.0003 \\
				&  & (0.0636) & (0.0453) & (0.0316) & (0.0197) & (0.0140) \\
				\quad x3 & 0.000 & -0.0001 & -0.0006 & -0.0005 & 0.0012 & -0.0003 \\
				&  & (0.0650) & (0.0457) & (0.0312) & (0.0199) & (0.0140) \\
				\quad x4 & 0.000 & -0.0034 & 0.0006 & -0.0002 & -0.0003 & -0.0001 \\
				&  & (0.0642) & (0.0452) & (0.0315) & (0.0196) & (0.0139) \\
				\midrule
				\multicolumn{7}{l}{\textit{Panel: Random Forest}} \\
				\quad x1 & 2.000 & -0.0973 & -0.0332 & -0.0387 & -0.0703 & -0.0911 \\
				&  & (0.2064) & (0.1189) & (0.0850) & (0.0824) & (0.0956) \\
				\quad x2 & 3.000 & -0.0215 & -0.0105 & -0.0433 & -0.0797 & -0.0953 \\
				&  & (0.2338) & (0.1456) & (0.0997) & (0.0958) & (0.1025) \\
				\quad x3 & 0.000 & -0.0015 & -0.0005 & -0.0005 & 0.0001 & 0.0000 \\
				&  & (0.0562) & (0.0290) & (0.0137) & (0.0036) & (0.0011) \\
				\quad x4 & 0.000 & -0.0021 & -0.0010 & -0.0005 & -0.0002 & 0.0000 \\
				&  & (0.0574) & (0.0297) & (0.0138) & (0.0037) & (0.0011) \\
				\midrule
				\multicolumn{7}{l}{\textit{Panel: XGBoost}} \\
				\quad x1 & 2.000 & 0.4114 & 0.1780 & 0.0366 & -0.0337 & -0.0467 \\
				&  & (0.4644) & (0.2227) & (0.0917) & (0.0554) & (0.0547) \\
				\quad x2 & 3.000 & 0.3122 & 0.0948 & -0.0172 & -0.0558 & -0.0601 \\
				&  & (0.4103) & (0.1884) & (0.0947) & (0.0754) & (0.0677) \\
				\quad x3 & 0.000 & 0.0007 & 0.0016 & -0.0003 & 0.0011 & -0.0001 \\
				&  & (0.1090) & (0.0550) & (0.0313) & (0.0153) & (0.0083) \\
				\quad x4 & 0.000 & -0.0080 & 0.0002 & 0.0011 & -0.0006 & 0.0001 \\
				&  & (0.1077) & (0.0597) & (0.0314) & (0.0149) & (0.0081) \\
				\midrule
				\multicolumn{7}{l}{\textit{Panel: Neural Net}} \\
				\quad x1 & 2.000 & -0.0659 & -0.0208 & -0.0063 & -0.0021 & -0.0008 \\
				&  & (0.1161) & (0.0544) & (0.0344) & (0.0293) & (0.0329) \\
				\quad x2 & 3.000 & -0.1231 & -0.0349 & -0.0102 & -0.0031 & -0.0006 \\
				&  & (0.1765) & (0.0638) & (0.0361) & (0.0324) & (0.0372) \\
				\quad x3 & 0.000 & 0.0014 & -0.0006 & -0.0003 & 0.0009 & -0.0006 \\
				&  & (0.0788) & (0.0484) & (0.0330) & (0.0250) & (0.0280) \\
				\quad x4 & 0.000 & -0.0002 & 0.0008 & -0.0003 & -0.0004 & -0.0017 \\
				&  & (0.0752) & (0.0496) & (0.0336) & (0.0247) & (0.0287) \\
				\bottomrule
				\end{tabular}
				\begin{tablenotes}
				\footnotesize
				\item \textit{Notes:} Mean bias and RMSE (in parentheses) computed over 1,000 Monte Carlo iterations. True AMEs computed via Monte Carlo integration with $n=1{,}000{,}000$ observations. Noise features x3 and x4 have true AME of zero.
				\end{tablenotes}
		\end{threeparttable}
	\end{table}
	
	Table~\ref{tab:sim1_regression_linear} turns to the regression setting with the linear DGP, $y = 2x_1 + 3x_2 + \varepsilon$. Linear regression recovers the true AMEs of 2.0 and 3.0 with negligible bias across all sample sizes, and RMSE declines from around 0.066 at $n=250$ to 0.014 at $n=5{,}000$, consistent with $O(n^{-1/2})$ convergence. The regression setting reveals the regularization bias of random forests more starkly than the classification setting: bias reaches $-0.091$ for $x_1$ and $-0.095$ for $x_2$ at $n=5{,}000$, and RMSE is not declining at the largest sample sizes, indicating that the bias is not vanishing with $n$ at default tuning. This is a known property of random forests in the regression case --- the ensemble averaging and subsampling impose shrinkage towards zero that does not vanish at these sample sizes. XGBoost shows large positive bias at small samples --- 0.41 for $x_1$ at $n=250$ --- but converges quickly, with bias falling below 0.05 by $n=2{,}500$. The neural network shows moderate negative bias at small samples that falls towards zero by $n=2{,}500$, though RMSE does not decline monotonically at the largest sample sizes, suggesting residual variance from the optimization landscape. The noise features have RMSE near zero for linear regression and declining rapidly for the black-box models, confirming that the estimator correctly attributes near-zero effects to irrelevant features in the regression setting.
	
	\begin{table}[htbp]
		\centering
		\small
		\begin{threeparttable}
			\caption{Simulation 1: AME Recovery --- Nonlinear DGP (Regression). Mean bias and RMSE (in parentheses) across 1,000 Monte Carlo iterations.}
			\label{tab:sim1_regression_nonlinear}
			\begin{tabular}{lrrrrrr}
				\toprule
				& & \multicolumn{5}{c}{Bias (RMSE)} \\
				\cmidrule(lr){3-7}
				Variable & True AME & $n=250$ & $n=500$ & $n=1,000$ & $n=2,500$ & $n=5,000$ \\
				\midrule
				\multicolumn{7}{l}{\textit{Panel: Linear}} \\
				\quad x1 & 0.000 & 0.0241 & 0.0261 & 0.0148 & 0.0048 & 0.0015 \\
				&  & (0.3868) & (0.2820) & (0.1966) & (0.1272) & (0.0904) \\
				\quad x2 & 3.000 & -0.0145 & -0.0041 & -0.0012 & -0.0002 & 0.0001 \\
				&  & (0.1921) & (0.1276) & (0.0930) & (0.0572) & (0.0417) \\
				\quad x3 & 0.000 & -0.0025 & 0.0001 & -0.0018 & -0.0020 & -0.0009 \\
				&  & (0.1889) & (0.1343) & (0.0941) & (0.0576) & (0.0414) \\
				\quad x4 & 0.000 & -0.0114 & -0.0003 & 0.0023 & 0.0006 & -0.0018 \\
				&  & (0.1857) & (0.1328) & (0.0927) & (0.0594) & (0.0430) \\
				\midrule
				\multicolumn{7}{l}{\textit{Panel: Random Forest}} \\
				\quad x1 & 0.000 & 0.0071 & 0.0097 & 0.0053 & 0.0006 & 0.0007 \\
				&  & (0.3433) & (0.2547) & (0.1836) & (0.1185) & (0.0903) \\
				\quad x2 & 3.000 & 0.0142 & -0.1152 & -0.1701 & -0.1867 & -0.1849 \\
				&  & (0.2704) & (0.2008) & (0.2004) & (0.1974) & (0.1916) \\
				\quad x3 & 0.000 & 0.0001 & 0.0008 & 0.0003 & -0.0002 & 0.0000 \\
				&  & (0.0549) & (0.0274) & (0.0138) & (0.0055) & (0.0027) \\
				\quad x4 & 0.000 & -0.0009 & -0.0007 & -0.0001 & -0.0000 & -0.0000 \\
				&  & (0.0573) & (0.0278) & (0.0138) & (0.0054) & (0.0025) \\
				\midrule
				\multicolumn{7}{l}{\textit{Panel: XGBoost}} \\
				\quad x1 & 0.000 & 0.0552 & 0.0038 & -0.0117 & -0.0137 & -0.0057 \\
				&  & (0.4662) & (0.2744) & (0.1829) & (0.1040) & (0.0735) \\
				\quad x2 & 3.000 & 0.2725 & 0.0816 & -0.0183 & -0.0702 & -0.0794 \\
				&  & (0.3937) & (0.1981) & (0.1058) & (0.0885) & (0.0862) \\
				\quad x3 & 0.000 & -0.0040 & -0.0022 & -0.0020 & -0.0011 & -0.0006 \\
				&  & (0.1124) & (0.0704) & (0.0446) & (0.0191) & (0.0097) \\
				\quad x4 & 0.000 & -0.0052 & -0.0028 & -0.0021 & 0.0003 & 0.0006 \\
				&  & (0.1121) & (0.0636) & (0.0381) & (0.0189) & (0.0096) \\
				\midrule
				\multicolumn{7}{l}{\textit{Panel: Neural Net}} \\
				\quad x1 & 0.000 & 0.0201 & 0.0178 & 0.0065 & -0.0003 & -0.0013 \\
				&  & (0.3086) & (0.1925) & (0.1390) & (0.0911) & (0.0760) \\
				\quad x2 & 3.000 & -0.1545 & -0.0202 & -0.0150 & -0.0098 & -0.0022 \\
				&  & (0.2601) & (0.0710) & (0.0461) & (0.0343) & (0.0395) \\
				\quad x3 & 0.000 & 0.0017 & 0.0004 & 0.0007 & 0.0009 & -0.0004 \\
				&  & (0.1419) & (0.0634) & (0.0400) & (0.0293) & (0.0321) \\
				\quad x4 & 0.000 & -0.0028 & -0.0001 & -0.0013 & -0.0004 & -0.0016 \\
				&  & (0.1410) & (0.0626) & (0.0406) & (0.0281) & (0.0323) \\
				\bottomrule
				\end{tabular}
				\begin{tablenotes}
				\footnotesize
				\item \textit{Notes:} Mean bias and RMSE (in parentheses) computed over 1,000 Monte Carlo iterations. True AMEs computed via Monte Carlo integration with $n=1{,}000{,}000$ observations. Noise features x3 and x4 have true AME of zero.
				\end{tablenotes}
		\end{threeparttable}
	\end{table}
	
	Table~\ref{tab:sim1_regression_nonlinear} reports results for the nonlinear regression DGP, $y = 2x_1^2 + 3x_2 + \varepsilon$, where the true AME of $x_1$ is zero by symmetry of $N(0,1)$ and the true AME of $x_2$ is 3.0. Linear regression recovers the $x_2$ AME accurately with RMSE declining at the standard rate, and correctly attributes a near-zero AME to $x_1$ despite misspecifying its functional form --- the same robustness result observed in the classification case. However, the RMSE for $x_1$ under linear regression is substantially larger than in the linear DGP --- 0.387 at $n=250$ versus 0.066 --- reflecting increased estimation variance from the misspecified model. The most striking result is random forests: the bias for $x_2$ reaches $-0.185$ at $n=5{,}000$ and is not declining, indicating that the regularization bias is amplified in the nonlinear setting where the forest must capture the quadratic relationship in $x_1$ while simultaneously estimating the linear effect of $x_2$. XGBoost and the neural network both show elevated RMSE relative to the linear DGP but declining bias, suggesting that these models are better able to adapt to the nonlinear DGP with sufficient data. The $x_1$ AME of zero is recovered accurately by all model classes at moderate sample sizes, which is a reassuring result for practitioners who want to test whether nonlinear features have nonzero average effects.
	
	\begin{table}[htbp]
		\centering
		\small
		\begin{threeparttable}
			\caption{Simulation 1: AME Recovery --- Interaction DGP (Regression). Mean bias and RMSE (in parentheses) across 1,000 Monte Carlo iterations.}
			\label{tab:sim1_regression_interaction}
			\begin{tabular}{lrrrrrr}
				\toprule
				& & \multicolumn{5}{c}{Bias (RMSE)} \\
				\cmidrule(lr){3-7}
				Variable & True AME & $n=250$ & $n=500$ & $n=1,000$ & $n=2,500$ & $n=5,000$ \\
				\midrule
				\multicolumn{7}{l}{\textit{Panel: Linear}} \\
				\quad x1 & 2.000 & -0.0041 & 0.0005 & 0.0000 & 0.0006 & 0.0011 \\
				&  & (0.2284) & (0.1531) & (0.1129) & (0.0715) & (0.0495) \\
				\quad x2 & 3.000 & 0.0040 & 0.0087 & 0.0075 & 0.0012 & 0.0009 \\
				&  & (0.2194) & (0.1627) & (0.1135) & (0.0727) & (0.0514) \\
				\quad x3 & 0.000 & -0.0114 & -0.0088 & -0.0054 & 0.0010 & -0.0009 \\
				&  & (0.1388) & (0.1002) & (0.0725) & (0.0450) & (0.0329) \\
				\quad x4 & 0.000 & -0.0073 & 0.0004 & -0.0006 & 0.0013 & 0.0005 \\
				&  & (0.1443) & (0.1018) & (0.0690) & (0.0453) & (0.0321) \\
				\midrule
				\multicolumn{7}{l}{\textit{Panel: Random Forest}} \\
				\quad x1 & 2.000 & -0.0755 & -0.0310 & -0.0433 & -0.0724 & -0.0804 \\
				&  & (0.2566) & (0.1663) & (0.1169) & (0.0979) & (0.0925) \\
				\quad x2 & 3.000 & -0.0634 & -0.0690 & -0.0922 & -0.1112 & -0.1143 \\
				&  & (0.2806) & (0.1913) & (0.1497) & (0.1306) & (0.1244) \\
				\quad x3 & 0.000 & -0.0031 & -0.0004 & -0.0009 & 0.0003 & -0.0001 \\
				&  & (0.0612) & (0.0352) & (0.0182) & (0.0067) & (0.0029) \\
				\quad x4 & 0.000 & -0.0023 & 0.0004 & -0.0001 & -0.0002 & 0.0000 \\
				&  & (0.0665) & (0.0349) & (0.0184) & (0.0069) & (0.0030) \\
				\midrule
				\multicolumn{7}{l}{\textit{Panel: XGBoost}} \\
				\quad x1 & 2.000 & 0.3184 & 0.1016 & -0.0009 & -0.0406 & -0.0429 \\
				&  & (0.4539) & (0.2184) & (0.1163) & (0.0758) & (0.0602) \\
				\quad x2 & 3.000 & 0.3672 & 0.1135 & -0.0083 & -0.0572 & -0.0589 \\
				&  & (0.4920) & (0.2373) & (0.1227) & (0.0898) & (0.0746) \\
				\quad x3 & 0.000 & -0.0031 & 0.0011 & -0.0017 & 0.0005 & -0.0001 \\
				&  & (0.1040) & (0.0551) & (0.0319) & (0.0151) & (0.0078) \\
				\quad x4 & 0.000 & -0.0082 & 0.0003 & -0.0003 & -0.0004 & -0.0001 \\
				&  & (0.1028) & (0.0564) & (0.0314) & (0.0149) & (0.0080) \\
				\midrule
				\multicolumn{7}{l}{\textit{Panel: Neural Net}} \\
				\quad x1 & 2.000 & -0.0842 & -0.0083 & -0.0056 & 0.0003 & 0.0022 \\
				&  & (0.2380) & (0.1114) & (0.0779) & (0.0573) & (0.0536) \\
				\quad x2 & 3.000 & -0.1231 & -0.0068 & -0.0052 & 0.0004 & 0.0026 \\
				&  & (0.2467) & (0.1120) & (0.0783) & (0.0579) & (0.0533) \\
				\quad x3 & 0.000 & -0.0026 & -0.0016 & -0.0024 & 0.0004 & -0.0004 \\
				&  & (0.0986) & (0.0551) & (0.0383) & (0.0293) & (0.0316) \\
				\quad x4 & 0.000 & -0.0079 & 0.0015 & 0.0002 & -0.0005 & -0.0019 \\
				&  & (0.0981) & (0.0565) & (0.0387) & (0.0284) & (0.0331) \\
				\bottomrule
				\end{tabular}
				\begin{tablenotes}
				\footnotesize
				\item \textit{Notes:} Mean bias and RMSE (in parentheses) computed over 1,000 Monte Carlo iterations. True AMEs computed via Monte Carlo integration with $n=1{,}000{,}000$ observations. Noise features x3 and x4 have true AME of zero.
				\end{tablenotes}
		\end{threeparttable}
	\end{table}
	
	Table~\ref{tab:sim1_regression_interaction} reports results for the interaction regression DGP, $y = 2x_1 + 3x_2 + 2x_1 x_2 + \varepsilon$, with true AMEs of 2.0 and 3.0 for $x_1$ and $x_2$. Linear regression recovers these AMEs accurately despite omitting the interaction term, confirming the robustness result from the classification case: the interaction contribution to the marginal effect of $x_1$ is $\mathbb{E}[2x_2] = 0$ and similarly for $x_2$, so the misspecified linear model nevertheless recovers the correct AME. The RMSE for the noise features under linear regression is elevated relative to the linear DGP --- around 0.14 at $n=250$ --- because the interaction term inflates the residual variance and increases estimation uncertainty for all coefficients. Random forests again show persistent negative bias that grows rather than shrinks with sample size, reaching $-0.080$ and $-0.114$ for $x_1$ and $x_2$ respectively at $n=5{,}000$. XGBoost shows the same large positive bias at small samples followed by convergence, but the bias does not fully vanish at $n=5{,}000$ --- around $-0.043$ for $x_1$ and $-0.059$ for $x_2$ --- suggesting that the interaction structure requires larger samples for the gradient boosting to fully adapt. The neural network performs well, with bias below 0.003 for both signal features at $n=2{,}500$ and $n=5{,}000$, and RMSE declining steadily. Taken together the six bias and RMSE tables establish that the AME estimator is consistent for logistic regression, neural networks, and XGBoost across all three DGPs and both outcome types, with random forests showing persistent regularization bias in the regression setting that practitioners should be aware of when reporting AMEs from default-tuned forest models.
	
	\subsection{Simulation 2: The Refitting Bootstrap}
	
	\begin{table}[htbp]
		\centering
		\footnotesize
		\begin{threeparttable}
			\caption{Simulation 2: Bootstrap SE Calibration --- Linear DGP (Classification). Nominal coverage target: 95\%. Coverage rate and mean 95\% CI width (in parentheses) across 500 Monte Carlo iterations.}
			\label{tab:sim2_calibration_linear}
			\begin{tabular}{lrrrr}
				\toprule
				& & \multicolumn{3}{c}{Coverage (CI Width)} \\
				\cmidrule(lr){3-5}
				Variable & True AME & $n=250$ & $n=1,000$ & $n=5,000$ \\
				\midrule
				\multicolumn{5}{l}{\textit{Panel: Logistic}} \\
				\quad x1 & 0.199 & 0.881 & 0.878 & 0.886 \\
				&  & (0.0652) & (0.0334) & (0.0150) \\
				\quad x2 & 0.298 & 0.726 & 0.766 & 0.781 \\
				&  & (0.0505) & (0.0257) & (0.0115) \\
				\quad x3 & 0.000 & 0.944 & 0.939 & 0.946 \\
				&  & (0.0752) & (0.0386) & (0.0174) \\
				\quad x4 & 0.000 & 0.933 & 0.954 & 0.944 \\
				&  & (0.0754) & (0.0386) & (0.0174) \\
				\midrule
				\multicolumn{5}{l}{\textit{Panel: Random Forest}} \\
				\quad x1 & 0.199 & 0.968 & 0.933 & 0.186 \\
				&  & (0.1343) & (0.0553) & (0.0205) \\
				\quad x2 & 0.298 & 0.961 & 0.958 & 0.268 \\
				&  & (0.1500) & (0.0614) & (0.0219) \\
				\quad x3 & 0.000 & 0.962 & 0.942 & 0.949 \\
				&  & (0.0976) & (0.0386) & (0.0132) \\
				\quad x4 & 0.000 & 0.969 & 0.964 & 0.953 \\
				&  & (0.0978) & (0.0386) & (0.0132) \\
				\midrule
				\multicolumn{5}{l}{\textit{Panel: XGBoost}} \\
				\quad x1 & 0.199 & 0.926 & 0.947 & 0.877 \\
				&  & (0.1864) & (0.0702) & (0.0230) \\
				\quad x2 & 0.298 & 0.908 & 0.954 & 0.910 \\
				&  & (0.2136) & (0.0794) & (0.0239) \\
				\quad x3 & 0.000 & 0.976 & 0.961 & 0.960 \\
				&  & (0.1139) & (0.0436) & (0.0141) \\
				\quad x4 & 0.000 & 0.973 & 0.969 & 0.959 \\
				&  & (0.1149) & (0.0436) & (0.0142) \\
				\midrule
				\multicolumn{5}{l}{\textit{Panel: Neural Net}} \\
				\quad x1 & 0.199 & 0.870 & 0.870 & 0.958 \\
				&  & (0.0679) & (0.0351) & (0.0203) \\
				\quad x2 & 0.298 & 0.658 & 0.774 & 0.894 \\
				&  & (0.0586) & (0.0287) & (0.0154) \\
				\quad x3 & 0.000 & 0.936 & 0.936 & 0.986 \\
				&  & (0.0764) & (0.0391) & (0.0225) \\
				\quad x4 & 0.000 & 0.938 & 0.936 & 0.982 \\
				&  & (0.0772) & (0.0392) & (0.0225) \\
				\bottomrule
				\end{tabular}
				\begin{tablenotes}
				\footnotesize
				\item \textit{Notes:} Coverage rate and mean 95\% CI width (in parentheses). Intervals are bootstrap percentile intervals. Logistic, random forest and XGBoost use 1{,}000 Monte Carlo iterations and 1{,}000 resamples; the neural network uses 500 and 200, since every one of its resamples is a fresh training run. True AMEs computed via Monte Carlo integration with $n=1{,}000{,}000$ observations. Noise features x3 and x4 have true AME of zero.
				\end{tablenotes}
		\end{threeparttable}
	\end{table}
	
	Table~\ref{tab:sim2_calibration_linear} reports coverage of nominal 95\% intervals for the linear classification design. Two distinct failures are visible, and Section~\ref{M-sec:inference} predicts both.
	
	The first is the one the paper is built around. Under the random forest, coverage at $n=5{,}000$ falls to $0.186$ for $x_1$ and $0.268$ for $x_2$ while remaining near nominal for the noise features, whose true effect is zero and whose regularization bias is therefore also zero. This is the mechanism of Remark~\ref{M-rem:refit}: the refitting bootstrap is consistent for the dispersion of the plug-in average around its own expectation, but that expectation is displaced from $\theta_{j,h}$ by $\E[\alpha_h(\hat f - f)]$, and resampling cannot see the displacement because every replicate recentres at a similarly regularized fit. The interval shrinks at the usual rate; the bias does not shrink at all; coverage follows the ratio down. XGBoost, whose bias does vanish with $n$, keeps coverage between $0.877$ and $0.976$ throughout, so the failure is specific to learners whose regularization bias persists rather than general to flexible learners.
	
	The second failure is milder, and it affects even a correctly specified model. Logistic regression covers at $0.881$, $0.878$ and $0.886$ for $x_1$ and at $0.726$, $0.766$ and $0.781$ for $x_2$, with no tendency to improve as the sample grows, while covering at roughly the nominal rate for the two noise features. The pattern orders the features by coefficient size --- $\beta_2 = 3$, $\beta_1 = 2$, and zero for $x_3$ and $x_4$ --- and that ordering identifies the cause. As Section~\ref{M-sec:inference} notes, each replicate is evaluated at the observed sample rather than at the resample, so the bootstrap distribution omits the sampling variability of the average of the derivative over the covariates. That omitted component is identically zero when the fitted derivative does not vary with $x$ and grows with how much it does, and under a logistic link the derivative varies most for the feature with the largest coefficient. The estimated standard errors are correspondingly too small: averaged over cells the ratio of the bootstrap standard error to the Monte Carlo standard deviation is $0.86$ for logistic regression, $0.89$ for the network and $0.90$ for the forest. Neither failure is detectable from the bootstrap output alone, which is the practical point: both produce intervals that look unremarkable.
	
	\begin{table}[htbp]
		\centering
		\footnotesize
		\begin{threeparttable}
			\caption{Simulation 2: Bootstrap SE Calibration --- Linear DGP (Regression). Nominal coverage target: 95\%. Coverage rate and mean 95\% CI width (in parentheses) across 500 Monte Carlo iterations.}
			\label{tab:sim2_calibration_regression_linear}
			\begin{tabular}{lrrrr}
				\toprule
				& & \multicolumn{3}{c}{Coverage (CI Width)} \\
				\cmidrule(lr){3-5}
				Variable & True AME & $n=250$ & $n=1,000$ & $n=5,000$ \\
				\midrule
				\multicolumn{5}{l}{\textit{Panel: Linear}} \\
				\quad x1 & 2.000 & 0.928 & 0.945 & 0.954 \\
				&  & (0.2478) & (0.1235) & (0.0553) \\
				\quad x2 & 3.000 & 0.955 & 0.948 & 0.956 \\
				&  & (0.2489) & (0.1239) & (0.0552) \\
				\quad x3 & 0.000 & 0.951 & 0.947 & 0.948 \\
				&  & (0.2495) & (0.1234) & (0.0552) \\
				\quad x4 & 0.000 & 0.949 & 0.948 & 0.957 \\
				&  & (0.2489) & (0.1234) & (0.0552) \\
				\midrule
				\multicolumn{5}{l}{\textit{Panel: Random Forest}} \\
				\quad x1 & 2.000 & 0.870 & 0.886 & 0.092 \\
				&  & (0.7447) & (0.2990) & (0.1170) \\
				\quad x2 & 3.000 & 0.967 & 0.962 & 0.323 \\
				&  & (0.9327) & (0.3632) & (0.1441) \\
				\quad x3 & 0.000 & 0.996 & 0.996 & 0.999 \\
				&  & (0.2555) & (0.0707) & (0.0082) \\
				\quad x4 & 0.000 & 0.992 & 0.995 & 1.000 \\
				&  & (0.2566) & (0.0708) & (0.0083) \\
				\midrule
				\multicolumn{5}{l}{\textit{Panel: XGBoost}} \\
				\quad x1 & 2.000 & 0.849 & 0.982 & 0.521 \\
				&  & (0.8174) & (0.3009) & (0.1041) \\
				\quad x2 & 3.000 & 0.917 & 0.990 & 0.503 \\
				&  & (0.9709) & (0.3552) & (0.1184) \\
				\quad x3 & 0.000 & 0.988 & 0.969 & 0.971 \\
				&  & (0.3863) & (0.1217) & (0.0382) \\
				\quad x4 & 0.000 & 0.978 & 0.975 & 0.969 \\
				&  & (0.3778) & (0.1210) & (0.0382) \\
				\midrule
				\multicolumn{5}{l}{\textit{Panel: Neural Net}} \\
				\quad x1 & 2.000 & 0.950 & 0.978 & 1.000 \\
				&  & (0.3442) & (0.1497) & (0.1938) \\
				\quad x2 & 3.000 & 0.936 & 0.974 & 1.000 \\
				&  & (0.4000) & (0.1537) & (0.2004) \\
				\quad x3 & 0.000 & 0.972 & 0.970 & 1.000 \\
				&  & (0.3075) & (0.1471) & (0.1887) \\
				\quad x4 & 0.000 & 0.968 & 0.972 & 1.000 \\
				&  & (0.3116) & (0.1474) & (0.1869) \\
				\bottomrule
				\end{tabular}
				\begin{tablenotes}
				\footnotesize
				\item \textit{Notes:} Coverage rate and mean 95\% CI width (in parentheses). Intervals are bootstrap percentile intervals. Logistic, random forest and XGBoost use 1{,}000 Monte Carlo iterations and 1{,}000 resamples; the neural network uses 500 and 200, since every one of its resamples is a fresh training run. True AMEs computed via Monte Carlo integration with $n=1{,}000{,}000$ observations. Noise features x3 and x4 have true AME of zero.
				\end{tablenotes}
		\end{threeparttable}
	\end{table}
	
	Table~\ref{tab:sim2_calibration_regression_linear} reports the corresponding regression results, where the two mechanisms separate cleanly. Linear regression covers between $0.928$ and $0.957$ at every sample size and every feature, and its bootstrap standard error matches the Monte Carlo standard deviation to within one percent. That is the second mechanism switching off exactly where the theory says it must: the fitted derivative of a linear model is constant in $x$, so the variance component that the evaluation point omits is identically zero, and the understatement that afflicted every classification panel disappears. Comparing the two tables therefore isolates the cause more sharply than either does alone.
	
	The first mechanism is correspondingly starker. The forest covers at $0.092$ for $x_1$ and $0.323$ for $x_2$ at $n=5{,}000$, and XGBoost --- which was well behaved in classification --- falls to $0.521$ and $0.503$. Here the forest's bootstrap standard error is if anything conservative, averaging $1.24$ times the Monte Carlo standard deviation, so the coverage failure cannot be attributed to an understated standard error. It is bias alone. The neural network sits at the opposite extreme, covering at $1.000$ for every feature at $n=5{,}000$: its estimate is dominated by optimization noise rather than sampling noise at that sample size, and because each replicate re-draws the initialization independently of the point fit, the replicate distribution centres near the truth rather than near the estimate. Its intervals are therefore wide enough to be useless rather than wrong, which the bootstrap output again does not reveal.
	
	\subsection{Simulation 3: The Debiased Estimator}
	
	Simulation 3 repeats the design of Simulation 2 for the debiased cross-fitted estimator of Section~\ref{M-sec:inference}. The covariates are independent standard normals, so the Riesz representer of Lemma~\ref{M-lem:riesz} is available in closed form, $\alpha_h(u) = e^{-h^2/2}\sinh(h u_j)/h$, and the estimator requires no representer estimation at all; a second arm reported at the end of this subsection replaces it with a sieve estimate, to show what is lost when the covariate density is unknown. Cross-fitting uses $K=5$ folds. The data generating processes, learners, hyperparameters, sample sizes, window $h$ and Monte Carlo ground truth are all held at the values used in the two preceding studies, so any difference between Simulation 3 and its predecessors is attributable to the orthogonal correction and to nothing else. Trimming is inactive throughout because the covariate support is unbounded and the trimming weight of Section~2.1 is therefore identically one. For the same reason the design satisfies Assumption~\ref{M-ass:unbounded} rather than the density bounds of Assumption~\ref{M-ass:sampling}: a Gaussian density is not bounded below on $\R^d$, and the representer $\alpha_h(u) = e^{-h^2/2}\sinh(h u_j)/h$ is unbounded. It is square integrable, with
	\begin{equation}
		\E\left[\alpha_h(X)^2\right]
		= \frac{e^{-h^2}}{h^2}\,\E\left[\sinh^2(h X_j)\right]
		= \frac{e^{-h^2}\left(e^{2h^2} - 1\right)}{2h^2} < \infty ,
		\label{eq:gaussalpha}
	\end{equation}
	and all higher moments finite as well; the shift ratio $\exp(h u_j - h^2/2)$ has moments of every order, so both conditions of Assumption~\ref{M-ass:unbounded} hold for every $r$ and $q$. The theorems demonstrated here are therefore the ones proved under that assumption set, and the closed-form representer used below is the exact one for this design rather than a bounded approximation to it.
	
	\begin{table}[htbp]
		\centering
		\small
		\begin{threeparttable}
			\caption{Simulation 3: debiased estimator, linear regression DGP. Bias and RMSE (in parentheses) by sample size.}
			\label{tab:sim3_bias_regression_linear}
			\begin{tabular}{lrrrrrr}
				\toprule
				& & \multicolumn{5}{c}{Bias (RMSE)} \\
				\cmidrule(lr){3-7}
				Variable & True AME & $n=250$ & $n=500$ & $n=1,000$ & $n=2,500$ & $n=5,000$ \\
				\midrule
				\multicolumn{7}{l}{\textit{Panel: Linear}} \\
				\quad x1 & 2.000 & -0.002 & -0.000 & -0.000 & -0.001 & -0.000 \\
				&  & (0.067) & (0.044) & (0.032) & (0.021) & (0.014) \\
				\quad x2 & 3.000 & -0.004 & -0.002 & 0.001 & -0.000 & -0.000 \\
				&  & (0.064) & (0.045) & (0.032) & (0.020) & (0.014) \\
				\quad x3 & 0.000 & -0.000 & -0.001 & -0.000 & 0.001 & -0.000 \\
				&  & (0.065) & (0.046) & (0.031) & (0.020) & (0.014) \\
				\quad x4 & 0.000 & -0.003 & 0.001 & -0.000 & -0.000 & -0.000 \\
				&  & (0.065) & (0.045) & (0.031) & (0.020) & (0.014) \\
				\midrule
				\multicolumn{7}{l}{\textit{Panel: Random Forest}} \\
				\quad x1 & 2.000 & -0.002 & 0.002 & -0.002 & 0.000 & 0.000 \\
				&  & (0.171) & (0.106) & (0.074) & (0.044) & (0.031) \\
				\quad x2 & 3.000 & -0.001 & 0.010 & 0.003 & 0.001 & -0.000 \\
				&  & (0.228) & (0.147) & (0.094) & (0.056) & (0.038) \\
				\quad x3 & 0.000 & 0.002 & 0.001 & -0.001 & 0.001 & -0.000 \\
				&  & (0.087) & (0.054) & (0.035) & (0.022) & (0.015) \\
				\quad x4 & 0.000 & -0.002 & 0.000 & -0.001 & -0.001 & -0.000 \\
				&  & (0.089) & (0.056) & (0.036) & (0.021) & (0.015) \\
				\midrule
				\multicolumn{7}{l}{\textit{Panel: XGBoost}} \\
				\quad x1 & 2.000 & -0.002 & 0.001 & -0.001 & -0.002 & 0.000 \\
				&  & (0.230) & (0.142) & (0.088) & (0.048) & (0.030) \\
				\quad x2 & 3.000 & -0.002 & 0.007 & 0.001 & 0.001 & 0.001 \\
				&  & (0.280) & (0.173) & (0.101) & (0.054) & (0.034) \\
				\quad x3 & 0.000 & -0.001 & 0.000 & -0.000 & 0.001 & -0.000 \\
				&  & (0.095) & (0.058) & (0.036) & (0.022) & (0.015) \\
				\quad x4 & 0.000 & -0.004 & 0.001 & 0.000 & -0.000 & -0.000 \\
				&  & (0.099) & (0.061) & (0.036) & (0.021) & (0.014) \\
				\midrule
				\multicolumn{7}{l}{\textit{Panel: Neural Net}} \\
				\quad x1 & 2.000 & -0.002 & -0.001 & 0.000 & -0.001 & -0.000 \\
				&  & (0.084) & (0.051) & (0.034) & (0.022) & (0.014) \\
				\quad x2 & 3.000 & -0.004 & -0.002 & 0.001 & -0.001 & -0.000 \\
				&  & (0.090) & (0.055) & (0.036) & (0.021) & (0.015) \\
				\quad x3 & 0.000 & 0.000 & -0.001 & -0.000 & 0.001 & -0.000 \\
				&  & (0.079) & (0.049) & (0.033) & (0.021) & (0.014) \\
				\quad x4 & 0.000 & -0.003 & 0.001 & -0.001 & 0.000 & -0.000 \\
				&  & (0.080) & (0.051) & (0.033) & (0.020) & (0.014) \\
				\bottomrule
				\end{tabular}
				\begin{tablenotes}
				\footnotesize
				\item \textit{Notes:} Debiased cross-fitted estimator with $K=5$ folds and the closed-form Riesz representer $\alpha_h(u) = e^{-h^2/2}\sinh(h u_j)/h$, exact for independent standard normal covariates. Trimming is inactive because the support is unbounded. Bias is the mean of $\hat{\theta}_{j,h} - \theta_{j,h}$ and RMSE its root mean square, over the Monte Carlo iterations. True values computed by Monte Carlo integration at $n=1{,}000{,}000$. Noise features $x_3$ and $x_4$ have a true effect of zero.
				\end{tablenotes}
		\end{threeparttable}
	\end{table}
	
	Table~\ref{tab:sim3_bias_regression_linear} reports bias and RMSE for the linear regression DGP. The result to compare against is the random forest panel of Table~\ref{tab:sim1_regression_linear}, where the plug-in's bias was $-0.097$ at $n=250$ and, rather than shrinking, grew to $-0.091$ for $x_1$ and $-0.095$ for $x_2$ at $n=5{,}000$. Under the correction that bias is $-0.002$ at $n=250$ and $0.000$ at $n=5{,}000$, and the RMSE now falls monotonically from $0.171$ to $0.031$, at the $O(n^{-1/2})$ rate the plug-in never achieved. The correction does not merely shrink the bias at a given sample size; it restores the rate, which is what inference requires. The other learners were already close to unbiased on this design and are essentially unchanged, as they should be: the orthogonal score corrects a first-order term that is near zero for them.
	
	\begin{table}[htbp]
		\centering
		\small
		\begin{threeparttable}
			\caption{Coverage of nominal 95\% confidence intervals, linear regression DGP: refitting bootstrap against the influence function.}
			\label{tab:sim3_coverage_regression_linear}
			\begin{tabular}{lrrrrrrr}
				\toprule
				& & \multicolumn{3}{c}{Refitting bootstrap} & \multicolumn{3}{c}{Influence function} \\
				\cmidrule(lr){3-5} \cmidrule(lr){6-8}
				Variable & True AME & $n=250$ & $n=1,000$ & $n=5,000$ & $n=250$ & $n=1,000$ & $n=5,000$ \\
				\midrule
				\multicolumn{8}{l}{\textit{Panel: Linear}} \\
				\quad x1 & 2.000 & 0.928 & 0.945 & 0.954 & 0.930 & 0.951 & 0.957 \\
				\quad x2 & 3.000 & 0.955 & 0.948 & 0.956 & 0.955 & 0.946 & 0.960 \\
				\quad x3 & 0.000 & 0.951 & 0.947 & 0.948 & 0.944 & 0.949 & 0.950 \\
				\quad x4 & 0.000 & 0.949 & 0.948 & 0.957 & 0.941 & 0.949 & 0.956 \\
				\midrule
				\multicolumn{8}{l}{\textit{Panel: Random Forest}} \\
				\quad x1 & 2.000 & 0.870 & 0.886 & 0.092 & 0.936 & 0.939 & 0.946 \\
				\quad x2 & 3.000 & 0.967 & 0.962 & 0.323 & 0.936 & 0.952 & 0.952 \\
				\quad x3 & 0.000 & 0.996 & 0.996 & 0.999 & 0.953 & 0.963 & 0.950 \\
				\quad x4 & 0.000 & 0.992 & 0.995 & 1.000 & 0.940 & 0.950 & 0.964 \\
				\midrule
				\multicolumn{8}{l}{\textit{Panel: XGBoost}} \\
				\quad x1 & 2.000 & 0.849 & 0.982 & 0.521 & 0.932 & 0.942 & 0.938 \\
				\quad x2 & 3.000 & 0.917 & 0.990 & 0.503 & 0.948 & 0.946 & 0.941 \\
				\quad x3 & 0.000 & 0.988 & 0.969 & 0.971 & 0.957 & 0.952 & 0.952 \\
				\quad x4 & 0.000 & 0.978 & 0.975 & 0.969 & 0.940 & 0.957 & 0.962 \\
				\midrule
				\multicolumn{8}{l}{\textit{Panel: Neural Net}} \\
				\quad x1 & 2.000 & 0.950 & 0.978 & 1.000 & 0.940 & 0.948 & 0.956 \\
				\quad x2 & 3.000 & 0.936 & 0.974 & 1.000 & 0.950 & 0.939 & 0.959 \\
				\quad x3 & 0.000 & 0.972 & 0.970 & 1.000 & 0.947 & 0.951 & 0.948 \\
				\quad x4 & 0.000 & 0.968 & 0.972 & 1.000 & 0.949 & 0.952 & 0.946 \\
				\bottomrule
				\end{tabular}
				\begin{tablenotes}
				\footnotesize
				\item \textit{Notes:} Empirical coverage of nominal 95\% intervals. The bootstrap columns resample and refit the learner, centring on the plug-in average; the influence-function columns use the debiased cross-fitted estimator with $K=5$ folds and the closed-form representer. Both are evaluated against the same window estimand $\theta_{j,h}$ and the same Monte Carlo ground truth. Noise features $x_3$ and $x_4$ have a true effect of zero.
				\end{tablenotes}
		\end{threeparttable}
	\end{table}
	
	Table~\ref{tab:sim3_coverage_regression_linear} places the two inference procedures side by side on the same designs and the same ground truth. The bootstrap columns range from $0.092$ to $1.000$; the influence-function columns range from $0.930$ to $0.964$. Across all 300 cells of the closed-form arm --- three data generating processes, two outcome types, five learners, five sample sizes and four features --- empirical coverage has mean $0.948$ and lies between $0.920$ and $0.970$, with $98.7\%$ of cells inside $[0.93, 0.97]$. No learner is an exception: mean coverage by model runs from $0.947$ to $0.949$. The forest cell that motivates the paper moves from $0.092$ to $0.946$.
	
	The forest panel carries a second reading. These forests are capped at depth six, and a depth-capped forest does not converge to the regression function; it converges to the best depth-limited approximation of it. The estimator is therefore consistent for neither nuisance in the usual sense, and yet its intervals cover at the nominal rate. This is Proposition~\ref{M-prop:robust} operating exactly as stated: with the representer known, $\sqrt{n}$-normality survives an outcome model that converges to the wrong function, and the learner's approximation error is paid in the width of the interval rather than in its location. The forest's RMSE at $n=5{,}000$ is $0.031$ against the linear model's $0.014$ --- wider, and correctly centred. That is the practical content of double robustness here, and it is why the choice of learner becomes a question of efficiency rather than of validity.
	
	\begin{table}[htbp]
		\centering
		\small
		\begin{threeparttable}
			\caption{Simulation 3: debiased estimator, linear classification DGP. Bias and RMSE (in parentheses) by sample size.}
			\label{tab:sim3_bias_classification_linear}
			\begin{tabular}{lrrrrrr}
				\toprule
				& & \multicolumn{5}{c}{Bias (RMSE)} \\
				\cmidrule(lr){3-7}
				Variable & True AME & $n=250$ & $n=500$ & $n=1,000$ & $n=2,500$ & $n=5,000$ \\
				\midrule
				\multicolumn{7}{l}{\textit{Panel: Logistic}} \\
				\quad x1 & 0.1986 & 0.0019 & 0.0005 & 0.0006 & 0.0000 & -0.0002 \\
				&  & (0.0199) & (0.0144) & (0.0105) & (0.0064) & (0.0048) \\
				\quad x2 & 0.2978 & 0.0004 & -0.0001 & 0.0003 & 0.0002 & 0.0002 \\
				&  & (0.0206) & (0.0142) & (0.0102) & (0.0065) & (0.0046) \\
				\quad x3 & 0.0000 & 0.0013 & 0.0007 & 0.0007 & 0.0003 & 0.0001 \\
				&  & (0.0208) & (0.0139) & (0.0100) & (0.0063) & (0.0045) \\
				\quad x4 & 0.0000 & 0.0008 & -0.0002 & -0.0008 & -0.0001 & 0.0003 \\
				&  & (0.0206) & (0.0139) & (0.0101) & (0.0062) & (0.0045) \\
				\midrule
				\multicolumn{7}{l}{\textit{Panel: Random Forest}} \\
				\quad x1 & 0.1986 & 0.0034 & -0.0009 & 0.0008 & 0.0003 & -0.0003 \\
				&  & (0.0339) & (0.0223) & (0.0152) & (0.0091) & (0.0060) \\
				\quad x2 & 0.2978 & -0.0003 & 0.0024 & 0.0004 & 0.0006 & 0.0004 \\
				&  & (0.0437) & (0.0269) & (0.0176) & (0.0101) & (0.0063) \\
				\quad x3 & 0.0000 & 0.0003 & 0.0009 & 0.0003 & 0.0003 & 0.0003 \\
				&  & (0.0246) & (0.0162) & (0.0111) & (0.0068) & (0.0047) \\
				\quad x4 & 0.0000 & -0.0001 & -0.0010 & -0.0010 & -0.0002 & 0.0002 \\
				&  & (0.0249) & (0.0168) & (0.0113) & (0.0065) & (0.0046) \\
				\midrule
				\multicolumn{7}{l}{\textit{Panel: XGBoost}} \\
				\quad x1 & 0.1986 & 0.0050 & -0.0025 & 0.0013 & 0.0001 & -0.0003 \\
				&  & (0.0495) & (0.0304) & (0.0193) & (0.0107) & (0.0069) \\
				\quad x2 & 0.2978 & -0.0003 & 0.0035 & 0.0005 & 0.0010 & 0.0003 \\
				&  & (0.0630) & (0.0395) & (0.0220) & (0.0119) & (0.0069) \\
				\quad x3 & 0.0000 & 0.0011 & 0.0007 & 0.0004 & 0.0002 & 0.0002 \\
				&  & (0.0260) & (0.0171) & (0.0112) & (0.0068) & (0.0047) \\
				\quad x4 & 0.0000 & 0.0004 & -0.0003 & -0.0004 & -0.0001 & 0.0003 \\
				&  & (0.0268) & (0.0174) & (0.0117) & (0.0065) & (0.0046) \\
				\midrule
				\multicolumn{7}{l}{\textit{Panel: Neural Net}} \\
				\quad x1 & 0.1986 & 0.0019 & 0.0005 & 0.0004 & 0.0001 & -0.0001 \\
				&  & (0.0217) & (0.0158) & (0.0115) & (0.0068) & (0.0049) \\
				\quad x2 & 0.2978 & 0.0004 & -0.0003 & 0.0000 & 0.0004 & 0.0002 \\
				&  & (0.0229) & (0.0171) & (0.0117) & (0.0072) & (0.0048) \\
				\quad x3 & 0.0000 & 0.0012 & 0.0010 & 0.0005 & 0.0004 & 0.0002 \\
				&  & (0.0216) & (0.0144) & (0.0101) & (0.0063) & (0.0045) \\
				\quad x4 & 0.0000 & 0.0008 & -0.0005 & -0.0008 & -0.0001 & 0.0003 \\
				&  & (0.0216) & (0.0141) & (0.0105) & (0.0063) & (0.0045) \\
				\bottomrule
				\end{tabular}
				\begin{tablenotes}
				\footnotesize
				\item \textit{Notes:} Debiased cross-fitted estimator with $K=5$ folds and the closed-form Riesz representer $\alpha_h(u) = e^{-h^2/2}\sinh(h u_j)/h$, exact for independent standard normal covariates. Trimming is inactive because the support is unbounded. Bias is the mean of $\hat{\theta}_{j,h} - \theta_{j,h}$ and RMSE its root mean square, over the Monte Carlo iterations. True values computed by Monte Carlo integration at $n=1{,}000{,}000$. Noise features $x_3$ and $x_4$ have a true effect of zero.
				\end{tablenotes}
		\end{threeparttable}
	\end{table}
	
	\begin{table}[htbp]
		\centering
		\small
		\begin{threeparttable}
			\caption{Coverage of nominal 95\% confidence intervals, linear classification DGP: refitting bootstrap against the influence function.}
			\label{tab:sim3_coverage_classification_linear}
			\begin{tabular}{lrrrrrrr}
				\toprule
				& & \multicolumn{3}{c}{Refitting bootstrap} & \multicolumn{3}{c}{Influence function} \\
				\cmidrule(lr){3-5} \cmidrule(lr){6-8}
				Variable & True AME & $n=250$ & $n=1,000$ & $n=5,000$ & $n=250$ & $n=1,000$ & $n=5,000$ \\
				\midrule
				\multicolumn{8}{l}{\textit{Panel: Logistic}} \\
				\quad x1 & 0.199 & 0.881 & 0.878 & 0.886 & 0.944 & 0.934 & 0.940 \\
				\quad x2 & 0.298 & 0.726 & 0.766 & 0.781 & 0.924 & 0.956 & 0.954 \\
				\quad x3 & 0.000 & 0.944 & 0.939 & 0.946 & 0.934 & 0.950 & 0.952 \\
				\quad x4 & 0.000 & 0.933 & 0.954 & 0.944 & 0.948 & 0.950 & 0.944 \\
				\midrule
				\multicolumn{8}{l}{\textit{Panel: Random Forest}} \\
				\quad x1 & 0.199 & 0.968 & 0.933 & 0.186 & 0.942 & 0.956 & 0.944 \\
				\quad x2 & 0.298 & 0.961 & 0.958 & 0.268 & 0.932 & 0.952 & 0.946 \\
				\quad x3 & 0.000 & 0.962 & 0.942 & 0.949 & 0.944 & 0.960 & 0.948 \\
				\quad x4 & 0.000 & 0.969 & 0.964 & 0.953 & 0.944 & 0.944 & 0.944 \\
				\midrule
				\multicolumn{8}{l}{\textit{Panel: XGBoost}} \\
				\quad x1 & 0.199 & 0.926 & 0.947 & 0.877 & 0.932 & 0.952 & 0.950 \\
				\quad x2 & 0.298 & 0.908 & 0.954 & 0.910 & 0.934 & 0.968 & 0.972 \\
				\quad x3 & 0.000 & 0.976 & 0.961 & 0.960 & 0.964 & 0.948 & 0.946 \\
				\quad x4 & 0.000 & 0.973 & 0.969 & 0.959 & 0.954 & 0.938 & 0.946 \\
				\midrule
				\multicolumn{8}{l}{\textit{Panel: Neural Net}} \\
				\quad x1 & 0.199 & 0.870 & 0.870 & 0.958 & 0.940 & 0.930 & 0.942 \\
				\quad x2 & 0.298 & 0.658 & 0.774 & 0.894 & 0.934 & 0.950 & 0.958 \\
				\quad x3 & 0.000 & 0.936 & 0.936 & 0.986 & 0.932 & 0.946 & 0.954 \\
				\quad x4 & 0.000 & 0.938 & 0.936 & 0.982 & 0.946 & 0.932 & 0.944 \\
				\bottomrule
				\end{tabular}
				\begin{tablenotes}
				\footnotesize
				\item \textit{Notes:} Empirical coverage of nominal 95\% intervals. The bootstrap columns resample and refit the learner, centring on the plug-in average; the influence-function columns use the debiased cross-fitted estimator with $K=5$ folds and the closed-form representer. Both are evaluated against the same window estimand $\theta_{j,h}$ and the same Monte Carlo ground truth. Noise features $x_3$ and $x_4$ have a true effect of zero.
				\end{tablenotes}
		\end{threeparttable}
	\end{table}
	
	Tables~\ref{tab:sim3_bias_classification_linear} and~\ref{tab:sim3_coverage_classification_linear} report the corresponding classification results, where the same pattern holds on a compressed scale: the forest's coverage at $n=5{,}000$ moves from $0.186$ to $0.944$ for $x_1$ and from $0.268$ to $0.946$ for $x_2$, and the debiased intervals are centred for every learner at every sample size. The nonlinear and interaction designs behave identically and are reported in the supplement.
	
	The closed-form representer used above is a luxury of the simulation design, and the natural question is what the estimator costs when $\alpha_h$ must itself be estimated. The second arm answers it by replacing the closed form with a sieve estimate, holding everything else fixed. Coverage survives, with a mild conservatism: over the same cells the sieve arm has mean coverage $0.958$ against the closed form's $0.948$. The cost appears where theory says it should, in the product bias of Assumption~\ref{M-ass:nuisance}(b) and in efficiency, and it vanishes with sample size. Mean absolute bias is $3.3$ times the closed form's at $n=250$, then $1.6$, $1.2$, $1.1$, and $0.9$ at $n=5{,}000$ --- indistinguishable by the largest sample size. Standard errors are $15\%$ wider at $n=250$ and $1\%$ wider at $n=5{,}000$. Estimating the representer costs something real at small samples and asymptotically nothing, which is the behaviour the product structure of the bias predicts: neither nuisance need converge quickly on its own, so long as the product of their errors is $o(n^{-1/2})$.
	
	One design remains to be accounted for. The arms reported so far draw independent standard normal covariates, and a Gaussian density is not bounded below on $\R^d$, so those designs satisfy Assumption~\ref{M-ass:unbounded} rather than the density bounds of Assumption~\ref{M-ass:sampling}. That is a real gap between what the theorems assume and what the simulation exhibits, and the natural way to close it is a design that meets the density-bounded hypotheses literally. The final arm does so: the same linear predictor, with the covariates drawn independently and uniformly on $[-\sqrt{3}, \sqrt{3}]$ so that each coordinate again has unit variance and the two designs are directly comparable. There the density is constant, and the representer follows from Lemma~\ref{M-lem:riesz} with the density cancelling entirely,
	
	\begin{equation}
		\alpha_h(u) = \frac{w(u - h e_j) - w(u + h e_j)}{2h}
		= \frac{\mathbb{1}\{u_j \geq a + 2h\} - \mathbb{1}\{u_j \leq b - 2h\}}{2h}
		\label{eq:unifalpha}
	\end{equation}
	
	on $[a, b] = [-\sqrt{3}, \sqrt{3}]$: the step function taking $-1/(2h)$ on the lower trimming shell, $+1/(2h)$ on the upper one, and zero between them. It is bounded by $1/(2h)$ and integrates to zero, so Assumption~\ref{M-ass:sampling} and Assumption~\ref{M-ass:nuisance}(c) hold as originally stated. The trimming weight is also active here rather than vacuous, which the Gaussian arms never exercise: the true effect of $x_1$ is $1.943$ rather than $2$, being exactly $2\,(1 - 2h/(b-a))$, so the boundary shell that the weight excises is visible in the target itself.
	
	\begin{table}[htbp]
		\centering
		\small
		\begin{threeparttable}
			\caption{Simulation 3: debiased estimator, bounded-support linear regression DGP. Bias and RMSE (in parentheses) by sample size.}
			\label{tab:sim3_bias_regression_linear_uniform}
			\begin{tabular}{lrrrrrr}
				\toprule
				& & \multicolumn{5}{c}{Bias (RMSE)} \\
				\cmidrule(lr){3-7}
				Variable & True AME & $n=250$ & $n=500$ & $n=1,000$ & $n=2,500$ & $n=5,000$ \\
				\midrule
				\multicolumn{7}{l}{\textit{Panel: Linear}} \\
				\quad x1 & 1.943 & 0.005 & 0.002 & -0.004 & -0.001 & -0.001 \\
				&  & (0.151) & (0.111) & (0.075) & (0.049) & (0.034) \\
				\quad x2 & 2.913 & -0.001 & -0.002 & -0.002 & -0.001 & -0.001 \\
				&  & (0.161) & (0.113) & (0.077) & (0.049) & (0.034) \\
				\quad x3 & 0.000 & 0.004 & 0.001 & 0.003 & 0.002 & 0.001 \\
				&  & (0.159) & (0.107) & (0.076) & (0.048) & (0.033) \\
				\quad x4 & 0.000 & -0.002 & 0.001 & -0.002 & 0.000 & 0.001 \\
				&  & (0.156) & (0.105) & (0.079) & (0.049) & (0.034) \\
				\midrule
				\multicolumn{7}{l}{\textit{Panel: Random Forest}} \\
				\quad x1 & 1.943 & 0.003 & 0.002 & -0.005 & 0.001 & -0.000 \\
				&  & (0.252) & (0.162) & (0.107) & (0.064) & (0.044) \\
				\quad x2 & 2.913 & -0.008 & -0.002 & 0.001 & -0.001 & -0.000 \\
				&  & (0.288) & (0.188) & (0.124) & (0.070) & (0.048) \\
				\quad x3 & 0.000 & 0.004 & -0.001 & 0.002 & 0.002 & 0.001 \\
				&  & (0.184) & (0.123) & (0.082) & (0.051) & (0.034) \\
				\quad x4 & 0.000 & -0.006 & 0.002 & -0.003 & -0.000 & 0.001 \\
				&  & (0.180) & (0.120) & (0.086) & (0.051) & (0.036) \\
				\midrule
				\multicolumn{7}{l}{\textit{Panel: XGBoost}} \\
				\quad x1 & 1.943 & -0.002 & -0.002 & -0.003 & -0.001 & -0.000 \\
				&  & (0.302) & (0.175) & (0.111) & (0.065) & (0.042) \\
				\quad x2 & 2.913 & -0.004 & 0.002 & 0.001 & -0.002 & 0.000 \\
				&  & (0.334) & (0.200) & (0.127) & (0.068) & (0.043) \\
				\quad x3 & 0.000 & 0.006 & -0.002 & 0.000 & 0.001 & 0.001 \\
				&  & (0.196) & (0.127) & (0.085) & (0.051) & (0.035) \\
				\quad x4 & 0.000 & -0.004 & 0.002 & -0.003 & 0.000 & 0.001 \\
				&  & (0.193) & (0.125) & (0.085) & (0.051) & (0.036) \\
				\midrule
				\multicolumn{7}{l}{\textit{Panel: Neural Net}} \\
				\quad x1 & 1.943 & 0.008 & 0.003 & -0.004 & -0.001 & -0.002 \\
				&  & (0.192) & (0.120) & (0.079) & (0.050) & (0.035) \\
				\quad x2 & 2.913 & -0.002 & -0.003 & -0.001 & -0.000 & -0.001 \\
				&  & (0.202) & (0.119) & (0.080) & (0.050) & (0.034) \\
				\quad x3 & 0.000 & 0.003 & 0.000 & 0.004 & 0.002 & 0.000 \\
				&  & (0.184) & (0.112) & (0.078) & (0.050) & (0.034) \\
				\quad x4 & 0.000 & -0.002 & 0.003 & -0.003 & 0.000 & 0.001 \\
				&  & (0.182) & (0.114) & (0.081) & (0.050) & (0.035) \\
				\bottomrule
				\end{tabular}
				\begin{tablenotes}
				\footnotesize
				\item \textit{Notes:} Debiased cross-fitted estimator with $K=5$ folds and the closed-form Riesz representer $\alpha_h(u) = e^{-h^2/2}\sinh(h u_j)/h$, exact for independent standard normal covariates. Trimming is inactive because the support is unbounded. Bias is the mean of $\hat{\theta}_{j,h} - \theta_{j,h}$ and RMSE its root mean square, over the Monte Carlo iterations. True values computed by Monte Carlo integration at $n=1{,}000{,}000$. Noise features $x_3$ and $x_4$ have a true effect of zero.
				\end{tablenotes}
		\end{threeparttable}
	\end{table}
	
	Table~\ref{tab:sim3_bias_regression_linear_uniform} reports bias and RMSE. The estimator behaves as it does on the unbounded designs: bias is at most $0.002$ on the signal features at any sample size and for any learner, and RMSE declines at the $O(n^{-1/2})$ rate. Coverage is the quantity of interest, and across the $160$ combinations of learner, sample size and feature in this arm it has mean $0.944$ and lies between $0.890$ and $0.970$, against a mean of $0.948$ over the matched Gaussian cells. The two assumption sets therefore deliver the same behaviour to within four thousandths. Since the two designs share a predictor and a coordinate variance and differ only in the covariate law, that agreement isolates boundedness as making no material difference: what the moment conditions buy is coverage of the Gaussian case, not a weaker guarantee within it. The influence-function standard errors are honest on both, with a mean ratio to the Monte Carlo standard deviation of $0.99$ here.
	
	The bounded design is nonetheless the slightly harder of the two at small samples, and the pattern is worth recording rather than averaging away. Fourteen of the $160$ cells fall below $0.93$; thirteen are classification and eleven have $n \leq 500$. The weakest is the logistic coefficient on $x_2$, whose coverage runs $0.912$, $0.890$, $0.908$, $0.948$ and $0.958$ across the five sample sizes while the ratio of the estimated standard error to the Monte Carlo standard deviation runs $0.87$, $0.91$, $0.96$, $1.00$ and $1.07$. This is finite-sample understatement of the standard error converging to correctness, not a failure of the assumption set: the mean ratio over all cells at $n=250$ is $0.958$ and at $n=5{,}000$ is $1.015$. The trimming weight is the plausible source, since it discards a shell of probability $O(h)$ and so costs effective sample size exactly where there is least of it to spare.
	

	
	% ============================================================
	\section{Proof of Lemma~\ref{M-lem:riesz} (Riesz Representation)}
	\label{app:lem_riesz}
	% ============================================================
	
	\begin{proof}
		Recall $q = wp$ extended by zero outside $\Omega$, so that $q$ is defined on
		all of $\R^d$, $\norm{q}_\infty \leq p_{\max}$, and $q(x) \neq 0$ implies
		$x \in \operatorname{supp}(w) \subseteq \Omega_{j,h}$. Fix $g \in L^1(P)$ and
		consider the positive-shift half of the finite difference:
		\begin{equation}
			\E\left[w(X)\, g(X + h e_j)\right]
			= \int_{\R^d} q(x)\, g(x + h e_j)\, dx
			= \int_{\R^d} q(u - h e_j)\, g(u)\, du
		\end{equation}
		by the substitution $u = x + h e_j$. The manipulation is legitimate because
		on the set where the integrand is nonzero we have $x \in \Omega_{j,h}$, hence
		$u = x + h e_j \in \Omega$, so $g$ is evaluated only on $\Omega$; and the
		integral is absolutely convergent since:
		\begin{equation}
			\int q(u - h e_j)\, \abs{g(u)}\, du
			\leq p_{\max} \int_\Omega \abs{g(u)}\, du
			\leq \frac{p_{\max}}{p_{\min}} \, \E\abs{g(X)} < \infty
		\end{equation}
		The negative-shift half is identical with $u = x - h e_j$. Subtracting and
		dividing by $2h$:
		\begin{equation}
			\E\left[w(X)\, \Dh g(X)\right]
			= \int_\Omega \frac{q(u - h e_j) - q(u + h e_j)}{2h}\, g(u)\, du
			= \E\left[\alpha_h(X)\, g(X)\right]
		\end{equation}
		where the restriction of the domain of integration to $\Omega$ in the second
		step uses that $g$ is supported there, and the final step multiplies and
		divides by $p(u) \geq p_{\min} > 0$. The uniform bound follows from
		$\abs{q(u - h e_j) - q(u + h e_j)} \leq 2 p_{\max}$ and $p(u) \geq p_{\min}$.
		Applying the representation to $g \equiv 1$ gives
		$\E[\alpha_h(X)] = \E[w(X) \Dh 1] = 0$ since the finite difference of a
		constant vanishes.
	\end{proof}
	
	% ============================================================
	\section{Proof of Proposition~\ref{M-prop:hlimit} (Recovery of the Average
		Derivative)}
	\label{app:prop_hlimit}
	% ============================================================
	
	\begin{proof}
		Let $f \in W^{1,1}(\Omega)$ and fix $h \leq h_0$, so that
		$\operatorname{supp}(w) \subseteq \Omega_{j,h_0} \subseteq \Omega_{j,h}$ and
		every segment $\{x + t e_j : \abs{t} \leq h\}$ with
		$x \in \operatorname{supp}(w)$ lies in $\Omega$. Sobolev functions are
		absolutely continuous on almost every line parallel to a coordinate axis, with
		the classical derivative along the line agreeing almost everywhere with the
		weak derivative \citep[Section~4.9]{evans2015measure}. Hence for almost every
		$x \in \operatorname{supp}(w)$:
		\begin{equation}
			f(x + h e_j) - f(x - h e_j)
			= \int_{-h}^{h} \partial_j f(x + t e_j)\, dt
		\end{equation}
		and therefore:
		\begin{equation}
			\theta_h(f) - \E\left[w(X)\, \partial_j f(X)\right]
			= \int_{\R^d} q(x)\, \frac{1}{2h} \int_{-h}^{h}
			\left[\partial_j f(x + t e_j) - \partial_j f(x)\right] dt\, dx
		\end{equation}
		Bounding $q$ by $p_{\max}$, applying Tonelli, and extending
		$\partial_j f$ by zero outside $\Omega$:
		\begin{equation}
			\abs{\theta_h(f) - \E\left[w\, \partial_j f\right]}
			\leq p_{\max} \sup_{\abs{t} \leq h}
			\int_{\R^d} \abs{\partial_j f(x + t e_j) - \partial_j f(x)}
			\, \mathbf{1}\{x \in \operatorname{supp}(w)\}\, dx
		\end{equation}
		where the extension introduces no error because $x + t e_j \in \Omega$ for
		$x \in \operatorname{supp}(w)$ and $\abs{t} \leq h \leq h_0$. Continuity of
		translation in $L^1(\R^d)$ applied to the extension of $\partial_j f$ sends
		the right side to zero as $h \downarrow 0$. For the second claim,
		$\abs{w_m \partial_j f} \leq \abs{\partial_j f}$, which is $P$-integrable
		because $\partial_j f \in L^1(\Omega)$ and $p \leq p_{\max}$; dominated
		convergence gives
		$\E[w_m \partial_j f] \to \E[\partial_j f] = \mathrm{AME}_j$.
	\end{proof}
	
	% ============================================================
	\section{Proof of Lemma~\ref{M-lem:orth} (Exact Product Bias)}
	\label{app:lem_orth}
	% ============================================================
	
	\begin{proof}
		Conditioning on $X$ and using $\E[Y \mid X] = f(X)$:
		\begin{equation}
			\E\left[\tilde{\alpha}(X)\,(Y - \tilde{f}(X))\right]
			= \E\left[\tilde{\alpha}(X)\,(f(X) - \tilde{f}(X))\right]
		\end{equation}
		By linearity of $\Dh$ and Lemma~\ref{M-lem:riesz} applied to
		$g = \tilde{f} - f \in L^1(P)$:
		\begin{equation}
			\E\left[w(X)\, \Dh \tilde{f}(X)\right] - \theta_h(f)
			= \E\left[w(X)\, \Dh (\tilde{f} - f)(X)\right]
			= \E\left[\alpha_h(X)\,(\tilde{f}(X) - f(X))\right]
		\end{equation}
		Summing the two displays gives:
		\begin{equation}
			\E\left[\psi(X, Y; \tilde{f}, \tilde{\alpha}, \theta_h(f))\right]
			= \E\left[(\alpha_h(X) - \tilde{\alpha}(X))(\tilde{f}(X) - f(X))\right]
		\end{equation}
		and the bound is Cauchy--Schwarz.
	\end{proof}
	
	% ============================================================
	\section{Proof of Theorem~\ref{M-thm:main} (Asymptotic Linearity and Normality)}
	\label{app:thm_main}
	% ============================================================
	
	\subsection{Step 0: The Trimmed Difference Operator is Bounded on $L^2(P)$}
	
	For any $g \in L^2(\Omega)$, using $(a - b)^2 \leq 2a^2 + 2b^2$ and the
	change of variables from the proof of Lemma~\ref{M-lem:riesz}:
	\begin{align}
		\E\left[\bigl(w(X)\, \Dh g(X)\bigr)^2\right]
		&\leq \frac{1}{2h^2}\left(
		\E\left[w(X)^2 g(X + h e_j)^2\right]
		+ \E\left[w(X)^2 g(X - h e_j)^2\right]\right) \notag \\
		&\leq \frac{1}{h^2} \cdot \frac{p_{\max}}{p_{\min}}\, \E\left[g(X)^2\right]
		\label{eq:opbound_proof}
	\end{align}
	which establishes \eqref{M-eq:opbound} with constant
	$\kappa = h^{-1}(p_{\max}/p_{\min})^{1/2}$. This is the only property of the
	finite difference used below; no derivative of $g$ appears anywhere.
	
	\subsection{Step 1: Decomposition by Folds}
	
	Write $\theta = \theta_h(f)$,
	$\psi_i = \psi(\mathbf{x}_i, y_i; f, \alpha_h, \theta)$, and for
	$i \in I_k$ let
	$\hat{\psi}_i = \psi(\mathbf{x}_i, y_i; \hat{f}^{(-k)},
	\hat{\alpha}^{(-k)}, \theta)$. Then:
	\begin{equation}
		\sqrt{n}\,(\hat{\theta}_h - \theta)
		= \frac{1}{\sqrt{n}} \sum_{i=1}^{n} \psi_i
		+ \sum_{k=1}^{K} \sqrt{\frac{n_k}{n}}\,\bigl(\sqrt{n_k}\, A_k
		+ \sqrt{n_k}\, B_k\bigr)
	\end{equation}
	where, with $\E_k$ denoting expectation over a fresh draw $(X, Y)$ holding the
	fold-$k$ nuisances fixed:
	\begin{equation}
		A_k = \E_k\left[\hat{\psi} - \psi\right],
		\qquad
		B_k = \frac{1}{n_k} \sum_{i \in I_k}
		\left(\hat{\psi}_i - \psi_i - A_k\right)
	\end{equation}
	Since $K$ is fixed and $n_k / n \to 1/K$, it suffices that
	$\sqrt{n_k} A_k = o_p(1)$ and $\sqrt{n_k} B_k = o_p(1)$ for each $k$.
	
	\subsection{Step 2: The Bias Term}
	
	By Lemma~\ref{M-lem:orth} with
	$(\tilde{f}, \tilde{\alpha}) = (\hat{f}^{(-k)}, \hat{\alpha}^{(-k)})$:
	\begin{equation}
		\abs{\sqrt{n_k}\, A_k}
		\leq \sqrt{n_k}\, \Ptwo{\hat{\alpha}^{(-k)} - \alpha_h}
		\cdot \Ptwo{\hat{f}^{(-k)} - f} = o_p(1)
	\end{equation}
	by Assumption~\ref{M-ass:nuisance}(b). This is the only place the product rate is
	used, and it is used without any expansion: the bound is exact.
	
	\subsection{Step 3: The Empirical Process Term}
	
	Conditional on the out-of-fold data, the summands of $B_k$ are i.i.d.\ and
	mean zero, so:
	\begin{equation}
		\E\left[(\sqrt{n_k}\, B_k)^2 \mid \text{out-of-fold}\right]
		\leq \E_k\left[(\hat{\psi} - \psi)^2\right]
	\end{equation}
	Decompose, adding and subtracting $\hat{\alpha}^{(-k)}(x)(y - f(x))$:
	\begin{equation}
		\hat{\psi} - \psi
		= w\, \Dh(\hat{f}^{(-k)} - f)
		+ (\hat{\alpha}^{(-k)} - \alpha_h)(y - f)
		- \hat{\alpha}^{(-k)}\,(\hat{f}^{(-k)} - f)
	\end{equation}
	and bound each piece in $L^2(P)$. The middle piece is bounded by
	$\bar{\sigma}\, \Ptwo{\hat{\alpha}^{(-k)} - \alpha_h}$ after conditioning on $X$
	under either assumption set, since it uses only the conditional variance bound of
	Assumption~\ref{M-ass:sampling}. The first and third pieces are where the two
	assumption sets diverge.

	Under the density bounds of Assumption~\ref{M-ass:sampling}, the first is bounded by
	$\kappa \Ptwo{\hat{f}^{(-k)} - f}$ via \eqref{eq:opbound_proof} and the third by
	$\bar{A}\, \Ptwo{\hat{f}^{(-k)} - f}$ using Assumption~\ref{M-ass:nuisance}(c). All
	three are then $o_p(1)$ by Assumption~\ref{M-ass:nuisance}(a).

	Under Assumption~\ref{M-ass:unbounded} neither bound is available: the representer
	is unbounded, and the finite-difference operator is genuinely unbounded on
	$L^2(P)$ rather than merely unbounded by this argument. For independent standard
	normal covariates, taking $g_a(x) = e^{a x_j}$ gives
	$\Ptwo{g_a}^2 = e^{2a^2}$ and $\Ptwo{g_a(\cdot + h e_j)}^2 = e^{2ah}e^{2a^2}$, so
	the operator norm is $\sup_a e^{ah} = \infty$. Both pieces are instead bounded by
	H\"{o}lder. Writing $\Delta = \hat{f}^{(-k)} - f$ and letting $r$ be as in
	\eqref{M-eq:shiftmoment} with conjugate $r'$, a change of variables gives
	\begin{equation}
		\E\left[\Delta(X \pm h e_j)^2\right]
		= \E\left[\Delta(X)^2\, \frac{p(X \mp h e_j)}{p(X)}\right]
		\leq \Pq{\Delta}{2r'}^2
		\left(\E\left[\left(\frac{p(X \mp h e_j)}{p(X)}\right)^{r}\right]\right)^{1/r}
		\label{eq:shiftbound}
	\end{equation}
	which is $o_p(1)$ by Assumption~\ref{M-ass:nuisance}(a) in its $L^{2+\delta}$ form,
	with $2r' = 2 + \delta$; the first piece follows from \eqref{eq:shiftbound} and
	$\abs{w} \leq 1$. For the third,
	$\E[(\hat{\alpha}^{(-k)})^2 \Delta^2]
	\leq \Pq{\hat{\alpha}^{(-k)}}{q}^2 \Pq{\Delta}{2s'}^2$ by H\"{o}lder with
	$s = q/2$, which is $o_p(1)$ by Assumption~\ref{M-ass:nuisance}(c) in its moment
	form together with (a).

	Either way $\E_k[(\hat{\psi} - \psi)^2] = o_p(1)$, and conditional Chebyshev gives
	$\sqrt{n_k}\, B_k = o_p(1)$. Note that only consistency of the nuisances is used
	here, never a rate; the rate enters at Step~2 alone.
	
	\subsection{Step 4: Central Limit Theorem and Variance Estimation}
	
	The $\psi_i$ are i.i.d., mean zero by Lemma~\ref{M-lem:orth} at
	$(\tilde{f}, \tilde{\alpha}) = (f, \alpha_h)$, with variance:
	\begin{equation}
		V_h = \E[\psi_i^2]
		\leq 2\,\E\left[(w \Dh f)^2\right] + 2\,\norm{\alpha_h}_\infty^2\,
		\E[\sigma^2(X)] < \infty
	\end{equation}
	by \eqref{eq:opbound_proof} and Lemma~\ref{M-lem:riesz} under the density bounds of
	Assumption~\ref{M-ass:sampling}. Under Assumption~\ref{M-ass:unbounded} the same
	finiteness follows from \eqref{eq:shiftbound} applied to $f$, which is bounded,
	together with $\alpha_h \in L^2(P)$ from
	Assumption~\ref{M-ass:unbounded}(c) and $\E[Y^2] < \infty$; for the Gaussian design
	the representer's second moment is \eqref{eq:gaussalpha}. Either way
	$\mathrm{Var}(\psi) < \infty$ and the classical CLT applies. For the variance estimator, write
	$\hat{\psi}_i(\hat{\theta}_h) = \hat{\psi}_i - (\hat{\theta}_h - \theta)$; then:
	\begin{equation}
		\hat{V}_h = \frac{1}{n} \sum_i \hat{\psi}_i^2
		- 2(\hat{\theta}_h - \theta)\,\frac{1}{n}\sum_i \hat{\psi}_i
		+ (\hat{\theta}_h - \theta)^2
	\end{equation}
	The last two terms vanish in probability since
	$\hat{\theta}_h - \theta = o_p(1)$ and
	$n^{-1}\sum_i \hat{\psi}_i = O_p(n^{-1/2})$ by Steps~1--3. For the first term:
	\begin{equation}
		\begin{split}
			\abs{\frac{1}{n}\sum_i \hat{\psi}_i^2 - \frac{1}{n}\sum_i \psi_i^2}
			&\leq \frac{1}{n}\sum_i (\hat{\psi}_i - \psi_i)^2 \\
			&\quad + 2\left(\frac{1}{n}\sum_i (\hat{\psi}_i - \psi_i)^2\right)^{1/2}
			\left(\frac{1}{n}\sum_i \psi_i^2\right)^{1/2}
		\end{split}
	\end{equation}
	by Cauchy--Schwarz; $n^{-1}\sum_i (\hat{\psi}_i - \psi_i)^2 = o_p(1)$ by the
	fold-wise conditional Markov inequality and the Step~3 bound, and
	$n^{-1}\sum_i \psi_i^2 \xrightarrowp V_h$ by the law of large numbers. Hence
	$\hat{V}_h \xrightarrowp V_h$.
	\qed
	
	% ============================================================
	\section{Proof of Proposition~\ref{M-prop:robust} (Double Robustness)}
	\label{app:prop_robust}
	% ============================================================
	
	The argument retraces the proof of Theorem~\ref{M-thm:main} with the pseudo-true
	influence function in place of the efficient one; we record the steps that change.
	
	\subsection{Step 1: The Pseudo-True Influence Function}
	
	Define $\psi_{g,i} = w(\mathbf{x}_i)\, \Dh g(\mathbf{x}_i)
	+ \alpha_h(\mathbf{x}_i)\,(y_i - g(\mathbf{x}_i)) - \theta_h(f)$. By
	Lemma~\ref{M-lem:orth} with $(\tilde{f}, \tilde{\alpha}) = (g, \alpha_h)$:
	\begin{equation}
		\E[\psi_{g,i}]
		= \E\bigl[(\alpha_h - \alpha_h)(g - f)\bigr] = 0
	\end{equation}
	for every $g \in L^2(\Omega)$, which also establishes
	$\E[m_g(X)] = \theta_h(f)$ after conditioning on $X$. Conditioning likewise gives
	$\E[\psi_{g,i} \mid X] = m_g(X) - \theta_h(f)$ and
	$\mathrm{Var}(\psi_{g,i} \mid X) = \alpha_h(X)^2 \sigma^2(X)$, so
	$\mathrm{Var}(\psi_{g,i}) = V_h(g)$ by the law of total variance; finiteness
	follows from \eqref{eq:opbound_proof}, $g \in L^2(\Omega)$, the bound on
	$\alpha_h$ from Lemma~\ref{M-lem:riesz}, and $\E[Y^2] < \infty$ under the density
	bounds; under Assumption~\ref{M-ass:unbounded} it follows from
	\eqref{eq:shiftbound}, $\alpha_h \in L^2(P)$, and $\E[Y^2] < \infty$.
	
	\subsection{Step 2: The Bias Term}
	
	With $A_k = \E_k[\hat{\psi} - \psi_g] = \E_k[\hat{\psi}]$, Lemma~\ref{M-lem:orth}
	gives $A_k = \E[(\alpha_h - \hat{\alpha}^{(-k)})(\hat{f}^{(-k)} - f)]$, so:
	\begin{equation}
		\begin{split}
			\abs{\sqrt{n_k}\, A_k}
			&\leq \sqrt{n_k}\, \Ptwo{\hat{\alpha}^{(-k)} - \alpha_h}
			\left( \Ptwo{\hat{f}^{(-k)} - g} + \Ptwo{g - f} \right) \\
			&= \sqrt{n_k}\, o_p(n^{-1/2}) \cdot O_p(1) = o_p(1)
		\end{split}
	\end{equation}
	by conditions (a) and (b); the second factor is $O_p(1)$ rather than $o_p(1)$
	because $g$ need not equal $f$, and it is exactly here that condition (b) must be
	strengthened to $o_p(n^{-1/2})$ relative to Assumption~\ref{M-ass:nuisance}(b).
	When $\hat{\alpha}^{(-k)} = \alpha_h$, $A_k = 0$ identically at every $n$;
	averaging the fold-wise identity
	$\E[\hat{\psi}_i + \theta_h(f) \mid \text{out-of-fold data}] = \theta_h(f)$,
	which holds because the observations in fold $k$ are independent of
	$\hat{f}^{(-k)}$, gives the exact unbiasedness claim.
	
	\subsection{Step 3: The Empirical Process Term}
	
	Decompose, for $i \in I_k$:
	\begin{equation}
		\hat{\psi}_i - \psi_{g,i}
		= w\, \Dh(\hat{f}^{(-k)} - g)
		+ (\hat{\alpha}^{(-k)} - \alpha_h)(y_i - \hat{f}^{(-k)})
		+ \alpha_h\,(g - \hat{f}^{(-k)})
	\end{equation}
	The first and third terms have conditional $L^2(P)$ norms bounded by
	$\kappa\, \Ptwo{\hat{f}^{(-k)} - g}$ and
	$\norm{\alpha_h}_\infty \Ptwo{\hat{f}^{(-k)} - g}$ respectively, both $o_p(1)$
	by condition (a). For the middle term, write
	$y - \hat{f}^{(-k)} = (y - f) + (f - g) + (g - \hat{f}^{(-k)})$ and bound the
	three resulting pieces: the first by
	$\bar{\sigma}^2 \Ptwo{\hat{\alpha}^{(-k)} - \alpha_h}^2 = o_p(1)$ after
	conditioning on $X$; the third by
	$(\bar{A} + \norm{\alpha_h}_\infty)^2 \Ptwo{\hat{f}^{(-k)} - g}^2 = o_p(1)$; and
	the second piece, $\E[(\hat{\alpha}^{(-k)} - \alpha_h)^2 (f - g)^2]$, tends to
	zero. Under the density bounds this is dominated convergence along subsequences:
	the factor $(\hat{\alpha}^{(-k)} - \alpha_h)^2$ is bounded by
	$(\bar{A} + \norm{\alpha_h}_\infty)^2$ and converges to zero in $P$-measure by
	condition (b), while $(f - g)^2 \in L^1(P)$ dominates. Under
	Assumption~\ref{M-ass:unbounded} the domination is unavailable and H\"{o}lder is
	used instead:
	$\E[(\hat{\alpha}^{(-k)} - \alpha_h)^2 (f - g)^2]
	\leq \Pq{\hat{\alpha}^{(-k)} - \alpha_h}{q}^2\, \Pq{f - g}{2t'}^2$ with
	$t = q/2$, which requires the fixed approximation error $f - g$ to lie in
	$L^{2t'}(P)$ --- a mild condition, and one that holds whenever $g$ and $f$ are
	bounded. Conditional Chebyshev then
	gives $\sqrt{n_k}\, B_k = o_p(1)$ exactly as in Step~3 of the proof of
	Theorem~\ref{M-thm:main}.
	
	\subsection{Step 4: Conclusion}
	
	The $\psi_{g,i}$ are i.i.d., mean zero, with finite positive variance $V_h(g)$,
	so the classical CLT applies to their normalized sum, and Steps~2--3 identify
	$\sqrt{n}(\hat{\theta}_h - \theta_h(f))$ with that sum up to $o_p(1)$. The
	consistency of $\hat{V}_h$ for $V_h(g)$ repeats Step~4 of the proof of
	Theorem~\ref{M-thm:main} verbatim with $\psi_{g}$ in place of $\psi$, using the
	Step~3 bound just established. Finally, when $g \neq f$ the influence function
	$\psi_g$ differs from the efficient influence function $\psi$ of
	Theorem~\ref{M-thm:efficiency}; since in the nonparametric model every regular
	asymptotically linear estimator has influence function equal to the efficient
	one, the estimator is not regular in this case, and no conflict with the
	convolution theorem arises.
	\qed
	
	% ============================================================
	\section{Proof of Theorem~\ref{M-thm:efficiency} (Semiparametric Efficiency)}
	\label{app:thm_eff}
	% ============================================================
	
	\begin{proof}
		Consider one-dimensional submodels
		$dP_t = (1 + t s)\, dP$ indexed by bounded functions
		$s \in L^2_0(P) = \{s : \E[s(X, Y)] = 0\}$, which are well defined for
		$\abs{t} < 1/\norm{s}_\infty$ and differentiable in quadratic mean at $t = 0$
		with score $s$; bounded mean-zero functions are dense in $L^2_0(P)$, so the
		tangent space of the nonparametric model is all of $L^2_0(P)$
		\citep[Chapter~25]{vaart1998}. Write
		$\bar{s}_X(x) = \E[s(X, Y) \mid X = x]$ for the induced score of the
		covariate marginal and
		$f_t(x) = \E_t[Y \mid X = x]$ for the regression function along the path.
		Differentiating the Bayes formula for $f_t$ at $t = 0$ gives, for
		$P$-almost every $x$:
		\begin{equation}
			\dot{f}(x) := \frac{\partial}{\partial t} f_t(x)\Big|_{t=0}
			= \E\left[(Y - f(X))\, s(X, Y) \mid X = x\right]
		\end{equation}
		where differentiation under the integral is justified by
		$\E[Y^2] < \infty$ and the boundedness of $s$. The functional along the path
		is $\theta(t) = \E_t[w(X)\, \Dh f_t(X)]$, and the product rule gives:
		\begin{equation}
			\frac{d\theta}{dt}\Big|_{t=0}
			= \E\left[w\, \Dh f(X)\cdot \bar{s}_X(X)\right]
			+ \E\left[w\, \Dh \dot{f}(X)\right]
		\end{equation}
		For the first term, note $\E[\bar{s}_X] = 0$, so the term equals
		$\E[(w \Dh f(X) - \theta_h(f))\, \bar{s}_X(X)]
		= \E[(w \Dh f(X) - \theta_h(f))\, s(X, Y)]$, the last equality because
		$w \Dh f(X) - \theta_h(f)$ is $X$-measurable. For the second term,
		$\dot{f} \in L^1(P)$ (indeed $L^2$) since $s$ is bounded, so
		Lemma~\ref{M-lem:riesz} applies:
		\begin{equation}
			\E\left[w\, \Dh \dot{f}(X)\right]
			= \E\left[\alpha_h(X)\, \dot{f}(X)\right]
			= \E\left[\alpha_h(X)\,(Y - f(X))\, s(X, Y)\right]
		\end{equation}
		by the tower property. Summing:
		\begin{equation}
			\frac{d\theta}{dt}\Big|_{t=0}
			= \E\left[\psi(X, Y; f, \alpha_h, \theta_h(f))\; s(X, Y)\right]
		\end{equation}
		for every score $s$, so $\theta_h$ is pathwise differentiable with influence
		function $\psi$; since $\psi \in L^2_0(P)$ lies in the tangent space, it is
		the efficient influence function and the bound is
		$\E[\psi^2]$ \citep[Theorem~25.20]{vaart1998}. The cross term in
		$\E[\psi^2]$ vanishes by conditioning on $X$, yielding
		\eqref{M-eq:effbound}. Attainment follows from Theorem~\ref{M-thm:main}, whose
		expansion exhibits $\hat{\theta}_h$ as asymptotically linear with influence
		function $\psi$, hence regular and efficient
		\citep[Lemma~25.23]{vaart1998}.
	\end{proof}
	
	% ============================================================
	\section{Proof of Lemma~\ref{lem:sobolev_existence}
		(Derivative Existence, Sobolev)}
	\label{app:lem_sobolev_existence}
	% ============================================================
	
	\subsection{What Needs to Be Shown}
	
	That $\fn \in W^{2,p}(\Omega)$ for $p > d$ implies $\nabla \fn \in
	W^{1,p}(\Omega)^d$ with the norm bound:
	\begin{equation}
		\norm{\nabla \fn}_{L^p}
		\leq \norm{\fn}_{W^{1,p}}
		\leq \norm{\fn}_{W^{2,p}}
	\end{equation}
	
	\begin{remark}
		The condition $p > d$ is required by the proof. It is the condition under
		which Theorem~6.5 of \citet{evans2015measure} guarantees pointwise
		differentiability of $W^{1,p}_{loc}$ functions. This is slightly stronger
		than the condition $p > 1$ stated in the main text and is an explicit
		assumption of the lemma.
	\end{remark}
	
	\subsection{Step 1: Existence of the Gradient as an $L^p$ Function}
	
	Since $\fn \in W^{2,p}(\Omega)$ and $p > d$, we have in particular that
	$\fn \in W^{1,p}(\Omega)$. Extending $\fn$ to a $W^{1,p}_{loc}(\R^d)$
	function by a standard extension operator \citep[Section~4.4]{evans2015measure},
	we may apply Theorem~6.5 of \citet{evans2015measure}:
	
	\begin{quote}
		\textit{Theorem~6.5 \citep{evans2015measure}.}
		Assume $f \in W^{1,p}_{loc}(\R^n)$ for some $n < p \leq \infty$. Then $f$
		is differentiable $\mathcal{L}^n$-a.e., and its classical derivative equals
		its weak derivative $\mathcal{L}^n$-a.e.
	\end{quote}
	
	Applying this theorem to $\fn$ with $n = d$ and $p > d$, we conclude that
	$\fn$ is classically differentiable at $\mathcal{L}^d$-almost every $x \in
	\Omega$, and:
	\begin{equation}
		\nabla \fn(x) = D\fn(x) \quad \mathcal{L}^d\text{-a.e.}
	\end{equation}
	where $D\fn$ denotes the weak gradient of $\fn$. Since $\fn \in W^{1,p}(\Omega)$,
	the weak gradient $D\fn \in L^p(\Omega)^d$ by definition of the Sobolev space.
	Therefore $\nabla \fn \in L^p(\Omega)^d$.
	
	Applying the same argument to each first-order weak derivative $\partial_j \fn
	\in W^{1,p}(\Omega)$ — which belongs to $W^{1,p}(\Omega)$ because $\fn \in
	W^{2,p}(\Omega)$ — we conclude that each $\partial_j \fn$ is classically
	differentiable $\mathcal{L}^d$-a.e.\ with classical second derivatives in
	$L^p(\Omega)$. Therefore $\nabla \fn \in W^{1,p}(\Omega)^d$.
	
	\subsection{Step 2: The Norm Bound}
	
	We establish the chain of inequalities:
	\begin{equation}
		\norm{\nabla \fn}_{L^p}
		\leq \norm{\fn}_{W^{1,p}}
		\leq \norm{\fn}_{W^{2,p}}
	\end{equation}
	
	\subsubsection*{First Inequality}
	
	By definition of the $W^{1,p}(\Omega)$ norm:
	\begin{equation}
		\norm{\fn}_{W^{1,p}}
		= \norm{\fn}_{L^p} + \sum_{j=1}^d \norm{\partial_j \fn}_{L^p}
		= \norm{\fn}_{L^p} + \norm{\nabla \fn}_{L^p}
	\end{equation}
	Since $\norm{\fn}_{L^p} \geq 0$:
	\begin{equation}
		\norm{\fn}_{W^{1,p}}
		= \norm{\fn}_{L^p} + \norm{\nabla \fn}_{L^p}
		\geq \norm{\nabla \fn}_{L^p}
	\end{equation}
	
	\subsubsection*{Second Inequality}
	
	By definition of the $W^{2,p}(\Omega)$ norm using multi-index notation:
	\begin{equation}
		\norm{\fn}_{W^{2,p}}
		= \sum_{|\alpha| \leq 2} \norm{\partial^\alpha \fn}_{L^p}
	\end{equation}
	The multi-indices with $|\alpha| \leq 1$ form a strict subset of those with
	$|\alpha| \leq 2$, and all terms are nonneg­ative, so:
	\begin{equation}
		\norm{\fn}_{W^{2,p}}
		= \sum_{|\alpha| \leq 2} \norm{\partial^\alpha \fn}_{L^p}
		\geq \sum_{|\alpha| \leq 1} \norm{\partial^\alpha \fn}_{L^p}
		= \norm{\fn}_{W^{1,p}}
	\end{equation}
	
	\subsection{Conclusion}
	
	For $p > d$, membership of $\fn$ in $W^{2,p}(\Omega)$ implies via
	Theorem~6.5 of \citet{evans2015measure} that $\nabla \fn$ exists
	$\mathcal{L}^d$-almost everywhere as a classical gradient, coincides with
	the weak gradient, and belongs to $W^{1,p}(\Omega)^d \subset L^p(\Omega)^d$.
	The gradient operator $\nabla: W^{2,p}(\Omega) \to W^{1,p}(\Omega)^d$ is
	bounded with:
	\begin{equation}
		\norm{\nabla \fn}_{L^p}
		\leq \norm{\fn}_{W^{1,p}}
		\leq \norm{\fn}_{W^{2,p}}
	\end{equation}
	\qed
	
	% ============================================================
	\section{Proof of Lemma~\ref{lem:sobolev_consistency}
		(Derivative Consistency, Sobolev)}
	\label{app:lem_sobolev_consistency}
	% ============================================================
	
	\subsection{What Needs to Be Shown}
	
	That well-behavedness in the Sobolev sense implies
	$\norm{\nabla \fn - \nabla f}_{L^p} \xrightarrowp 0$.
	
	\subsection{Proof Strategy}

	Compactness is a statement about deterministic sequences and $\fn$ is
	random, so the two are kept apart. We first prove a deterministic lemma: on
	a ball of $\Wtwop$, closeness to $f$ in $L^p$ forces closeness of the
	gradients in $L^p$, with a modulus depending only on the radius of the ball.
	The probabilistic conclusion then follows in a few lines, the tightness
	condition \eqref{eq:tight_sobolev} supplying the radius.

	The order matters. Applying Rellich--Kondrachov directly to $\{\fn\}$
	extracts a subsequence that depends on the realization, and an almost sure
	statement along a realization-dependent subsequence does not combine with a
	probability bound along a deterministic one. Separating the compactness from
	the randomness avoids the difficulty rather than managing it, and has the
	further advantage of needing only tightness of the $\Wtwop$ norms, which is
	what Definition~\ref{def:wellbehaved_sobolev} provides and all that the
	estimators of Section~\ref{sec:learners} deliver.

	\subsection{Step 1: A Deterministic Gradient Modulus}

	\begin{lemma}[Gradient modulus on a Sobolev ball]
		\label{lem:sobolev_modulus}
		Let $\Omega$ be bounded and Lipschitz, let $p > 1$, and let
		$f \in \Wonep$. For every $M < \infty$ and $\delta > 0$ there is an
		$\epsilon = \epsilon(M,\delta) > 0$ such that every $g \in \Wtwop$ with
		\begin{equation}
			\norm{g}_{W^{2,p}} \leq M
			\qquad\text{and}\qquad
			\norm{g - f}_{L^p} \leq \epsilon
		\end{equation}
		satisfies $\norm{\nabla g - \nabla f}_{L^p} \leq \delta$.
	\end{lemma}

	\begin{proof}
		Suppose not. Then for some $M$ and some $\delta$ there is a sequence
		$\{g_k\} \subset \Wtwop$ with $\norm{g_k}_{W^{2,p}} \leq M$ and
		$\norm{g_k - f}_{L^p} \to 0$, but
		$\norm{\nabla g_k - \nabla f}_{L^p} > \delta$ for every $k$: negating the
		statement for a fixed $M$ and $\delta$ produces, for each
		$\epsilon = 1/k$, a witness $g_k$.

		This is a deterministic bounded sequence in $\Wtwop$. By the
		Rellich--Kondrachov theorem the embedding
		$W^{2,p}(\Omega) \hookrightarrow\hookrightarrow W^{1,p}(\Omega)$ is
		compact \citep[Theorem~5,~Section~5.7]{evans2010pde}, so some
		subsequence $\{g_{k_j}\}$ converges strongly in $\Wonep$ to a limit
		$h \in \Wonep$. Strong $\Wonep$ convergence implies strong $L^p$
		convergence, and $g_k \to f$ in $L^p$ by construction, so $h = f$
		$\mathcal{L}^d$-almost everywhere by uniqueness of limits in
		$L^p(\Omega)$. Hence $\norm{g_{k_j} - f}_{W^{1,p}} \to 0$, and since
		$g_{k_j} - f \in \Wonep$ with
		$\norm{\nabla g - \nabla f}_{L^p} \leq \norm{g - f}_{W^{1,p}}$, this
		gives $\norm{\nabla g_{k_j} - \nabla f}_{L^p} \to 0$, contradicting
		$\norm{\nabla g_{k_j} - \nabla f}_{L^p} > \delta$.
	\end{proof}

	Being a compactness statement, the lemma carries no rate: $\epsilon(M,\delta)$
	is not quantified and cannot be without further assumptions on $f$. That is
	sufficient here, because what is being established is consistency and not a
	rate of convergence. It is also where the extra degree of differentiability
	is spent --- the hypothesis is a bound in $\Wtwop$ and the conclusion is
	about gradients in $L^p$, a gap of one derivative that the compact embedding
	closes and the continuous embedding of Lemma~\ref{lem:sobolev_existence}
	does not.

	\subsection{Step 2: The Probabilistic Conclusion}

	Fix $\delta > 0$ and $\eta > 0$. By the tightness condition
	\eqref{eq:tight_sobolev} there is an $M < \infty$ with
	$\sup_n \Prob(\norm{\fn}_{W^{2,p}} > M) \leq \eta/2$; let
	$\epsilon = \epsilon(M,\delta) > 0$ be the radius supplied by
	Lemma~\ref{lem:sobolev_modulus} for this $M$ and $\delta$. On the event
	\begin{equation}
		A_n = \left\{\norm{\fn}_{W^{2,p}} \leq M\right\}
		\cap \left\{\norm{\fn - f}_{L^p} \leq \epsilon\right\}
	\end{equation}
	the lemma applies pointwise in $\omega$ --- for each such $\omega$ the
	realization $\fn(\omega)$ is a fixed element of $\Wtwop$ satisfying both
	hypotheses --- and yields
	$\norm{\nabla \fn - \nabla f}_{L^p} \leq \delta$ on $A_n$. Taking
	complements,
	\begin{equation}
		\Prob\!\left(\norm{\nabla \fn - \nabla f}_{L^p} > \delta\right)
		\leq \Prob(A_n^c)
		\leq \frac{\eta}{2}
		+ \Prob\!\left(\norm{\fn - f}_{L^p} > \epsilon\right) .
	\end{equation}
	The remaining term tends to zero by the $L^p$ consistency condition of
	Definition~\ref{def:wellbehaved_sobolev}, with $\epsilon$ held fixed, so the
	left-hand side is at most $\eta$ for all $n$ sufficiently large. As $\delta$
	and $\eta$ were arbitrary,
	\begin{equation}
		\norm{\nabla \fn - \nabla f}_{L^p} \xrightarrowp 0
	\end{equation}
	which is the conclusion of Lemma~\ref{lem:sobolev_consistency}. \qed
	
	\begin{remark}
		The key step is the application of Rellich--Kondrachov in Step~1. The
		theorem requires $\Omega$ to be bounded and Lipschitz --- both of which are
		assumed throughout --- and $p > 1$, which is assumed in the lemma statement.
		The compact embedding $W^{2,p} \hookrightarrow\hookrightarrow W^{1,p}$ is
		strictly stronger than the continuous embedding $W^{2,p} \hookrightarrow
		W^{1,p}$: the latter follows from the norm bound in
		Lemma~\ref{lem:sobolev_existence}, while the former provides the
		compactness needed to extract convergent subsequences. It is this
		compactness, not merely the norm bound, that allows the interchange of
		limit and gradient in the proof.
	\end{remark}
	
	% ============================================================
	\section{Proof of Lemma~\ref{lem:bv_existence}
		(Derivative Existence, BV)}
	\label{app:lem_bv_existence}
	% ============================================================
	
	\subsection{What Needs to Be Shown}
	
	Two things: that $BV$ membership implies $D\fn$ exists as a Radon measure
	with finite total variation, and the explicit formula for piecewise constant
	functions.
	
	\subsection{Step 1: Existence of $D\fn$ as a Radon Measure}
	
	We establish existence in three steps, using the Riesz representation
	theorem and the Radon-Nikodym theorem.
	
	\subsubsection*{Substep 1a: The Distributional Gradient as a Linear Functional}
	
	For any $\fn \in \BV$, the distributional gradient is defined as the linear
	functional $\Lambda: C_c^1(\Omega)^d \to \R$ given by:
	\begin{equation}
		\Lambda(\phi) = -\int_\Omega \fn \operatorname{div} \phi \, dx
		\quad \forall \phi \in C_c^1(\Omega)^d
	\end{equation}
	By the definition of $\BV$, this functional has finite total variation:
	\begin{equation}
		|\Lambda|(\Omega)
		= \sup\left\{
		\Lambda(\phi) : \phi \in C_c^1(\Omega)^d,\, \norm{\phi}_{L^\infty} \leq 1
		\right\}
		< \infty
	\end{equation}
	This is precisely the condition that $\fn \in \BV$ imposes --- the
	distributional gradient is a bounded linear functional on $C_c^1(\Omega)^d$
	with respect to the supremum norm.
	
	\subsubsection*{Substep 1b: Riesz Representation}
	
	Since $\Lambda$ is a bounded linear functional on $C_c^1(\Omega)^d$ with
	finite total variation, the Riesz representation theorem
	\citep[Theorem~1.38]{evans2015measure} guarantees the existence of a
	vector-valued Radon measure $D\fn \in \mathcal{M}(\Omega)^d$ such that:
	\begin{equation}
		\Lambda(\phi) = \int_\Omega \phi \cdot d(D\fn)
		\quad \forall \phi \in C_c^1(\Omega)^d
	\end{equation}
	That is:
	\begin{equation}
		-\int_\Omega \fn \operatorname{div} \phi \, dx
		= \int_\Omega \phi \cdot d(D\fn)
		\quad \forall \phi \in C_c^1(\Omega)^d
	\end{equation}
	This is the distributional gradient equation, and $D\fn$ is the
	distributional gradient of $\fn$ represented as a Radon measure.
	
	\subsubsection*{Substep 1c: Finiteness of Total Variation}
	
	By the Radon-Nikodym theorem for Radon measures
	\citep[Theorem~1.29]{evans2015measure}, the derivative $D_{\mathcal{L}^d}(D\fn)$
	exists and is finite $\mathcal{L}^d$-almost everywhere on $\Omega$. Since
	$\Omega$ is bounded and the total variation of $\Lambda$ is finite by the
	$BV$ assumption, we conclude:
	\begin{equation}
		|D\fn|(\Omega) = |\Lambda|(\Omega) < \infty
	\end{equation}
	Therefore $D\fn \in \mathcal{M}(\Omega)^d$ with $|D\fn|(\Omega) < \infty$,
	as claimed.
	
	\subsection{Step 2: Explicit Formula for Piecewise Constant Functions}
	
	Let $\fn = \sum_j \bar{y}_j \mathbf{1}_{R_j}$ where $\{R_j\}$ is a finite
	partition of $\Omega$ into sets with Lipschitz boundaries, and $\bar{y}_j$
	denotes the response mean on leaf $R_j$. We derive the explicit formula for
	$D\fn$ by applying the Gauss--Green theorem to each leaf.
	
	\subsubsection*{Substep 2a: Applying the Definition}
	
	For any test function $\phi \in C_c^1(\Omega)^d$, the distributional
	gradient equation from Substep~1b gives:
	\begin{equation}
		\int_\Omega \phi \cdot d(D\fn)
		= -\int_\Omega \fn \operatorname{div} \phi \, dx
		= -\sum_j \bar{y}_j \int_{R_j} \operatorname{div} \phi \, dx
	\end{equation}
	where the second equality uses the piecewise constant structure of $\fn$ and
	linearity of the integral.
	
	\subsubsection*{Substep 2b: Gauss--Green on Each Leaf}
	
	Since each $R_j$ has Lipschitz boundary, the Gauss--Green theorem for $BV$
	functions \citep[Theorem~5.16]{evans2015measure} applies to each leaf:
	\begin{equation}
		\int_{R_j} \operatorname{div} \phi \, dx
		= \int_{\partial R_j} \phi \cdot \nu_j \, d\mathcal{H}^{d-1}
	\end{equation}
	where $\nu_j$ denotes the outward unit normal to $\partial R_j$ and
	$\mathcal{H}^{d-1}$ is the $(d-1)$-dimensional Hausdorff measure on the
	boundary.
	
	\subsubsection*{Substep 2c: Cancellation on Interior Boundaries}
	
	Substituting the Gauss--Green formula into the expression from Substep~2a:
	\begin{equation}
		\int_\Omega \phi \cdot d(D\fn)
		= -\sum_j \bar{y}_j \int_{\partial R_j} \phi \cdot \nu_j \, d\mathcal{H}^{d-1}
	\end{equation}
	Every interior boundary $\partial R_j \cap \partial R_k$ between adjacent
	leaves $j$ and $k$ appears twice in this sum --- once with outward normal
	$\nu_j$ for leaf $R_j$ and once with outward normal $\nu_k = -\nu_j$ for
	leaf $R_k$. Collecting these paired contributions and writing $\nu_{jk} =
	\nu_j$ for the unit normal pointing from $R_j$ into $R_k$:
	\begin{align}
		-\sum_j \bar{y}_j \int_{\partial R_j} \phi \cdot \nu_j \, d\mathcal{H}^{d-1}
		&= -\sum_{j \sim k} \bar{y}_j \int_{\partial R_j \cap \partial R_k}
		\phi \cdot \nu_{jk} \, d\mathcal{H}^{d-1} \notag \\
		&\quad + \sum_{j \sim k} \bar{y}_k \int_{\partial R_j \cap \partial R_k}
		\phi \cdot \nu_{jk} \, d\mathcal{H}^{d-1} \notag \\
		&= \sum_{j \sim k} (\bar{y}_k - \bar{y}_j)
		\int_{\partial R_j \cap \partial R_k} \phi \cdot \nu_{jk}
		\, d\mathcal{H}^{d-1}
	\end{align}
	Boundary terms on $\partial\Omega$ vanish because $\phi \in C_c^1(\Omega)^d$
	has compact support in $\Omega$.
	
	\subsubsection*{Substep 2d: Identifying the Measure}
	
	Since the equation:
	\begin{equation}
		\int_\Omega \phi \cdot d(D\fn)
		= \sum_{j \sim k} (\bar{y}_k - \bar{y}_j)
		\int_{\partial R_j \cap \partial R_k} \phi \cdot \nu_{jk}
		\, d\mathcal{H}^{d-1}
	\end{equation}
	holds for all $\phi \in C_c^1(\Omega)^d$, the uniqueness part of the Riesz
	representation theorem \citep[Theorem~1.38]{evans2015measure} identifies
	$D\fn$ uniquely as the measure:
	\begin{equation}
		D\fn = \sum_{j \sim k} (\bar{y}_k - \bar{y}_j) \cdot \nu_{jk} \cdot
		\mathcal{H}^{d-1}\lfloor_{\partial R_j \cap \partial R_k}
	\end{equation}
	where $\mathcal{H}^{d-1}\lfloor_{\partial R_j \cap \partial R_k}$ denotes
	the $(d-1)$-dimensional Hausdorff measure restricted to the shared boundary
	between leaves $j$ and $k$. This is the formula stated in the lemma.
	\qed
	
	% ============================================================
	\section{Proof of Lemma~\ref{lem:bv_consistency}
		(Derivative Consistency, BV)}
	\label{app:lem_bv_consistency}
	% ============================================================
	
	\subsection{What Needs to Be Shown}
	
	That well-behavedness in the BV sense implies $D\fn \xrightarrow{w*} Df$
	in $\cM$, and the strict convergence upgrade under the additional condition
	$\E[|D\fn|(\Omega)] \to |Df|(\Omega)$.
	
	\subsection{Proof Strategy}
	
	The proof mirrors the subsequence argument of
	Lemma~\ref{lem:sobolev_consistency}, replacing the Rellich--Kondrachov
	compact embedding with the BV compactness theorem and replacing strong $L^p$
	convergence with weak-* convergence of measures. We show that every
	subsequence of $\{D\fn\}$ has a further subsequence converging weak-* to
	$Df$. By the metric space characterization of convergence in probability,
	this implies $D\fn \xrightarrow{w*} Df$ in $\cM$ in probability.
	
	Suppose for contradiction that $D\fn$ does not converge weak-* to $Df$ in
	probability. Then there exist $\epsilon > 0$, $\delta > 0$, a test function
	$\phi \in C_0(\Omega)^d$ with $\norm{\phi}_\infty \leq 1$, and a subsequence
	$\{n_k\}$ such that:
	\begin{equation}
		\mathbb{P}\!\left(\left|
		\int_\Omega \phi \cdot d(D\hat{f}_{n_k})
		- \int_\Omega \phi \cdot d(Df)
		\right| > \epsilon\right) > \delta
		\quad \forall k
	\end{equation}
	We derive a contradiction by extracting a further subsequence along which
	weak-* convergence holds almost surely.
	
	\subsection{Part I: Weak-* Convergence}
	
	\subsubsection*{Step 1: Tightness of the $BV$ Norms}

	The well-behaved condition in the $BV$ sense controls $|D\fn|(\Omega)$ in
	expectation, and an expectation bound does not make the sequence bounded
	almost surely: the total variation may be arbitrarily large on an event of
	vanishing but positive probability for every $n$. What the condition does
	deliver is uniform tightness, and tightness is what the subsequence
	architecture needs. Fix $\epsilon$, $\delta$, $\phi$ and $\{n_k\}$ as in the
	contradiction hypothesis and write
	\begin{equation}
		A_k = \left\{\left|
		\int_\Omega \phi \cdot d(D\hat{f}_{n_k})
		- \int_\Omega \phi \cdot d(Df)\right| > \epsilon\right\},
		\qquad \Prob(A_k) > \delta .
	\end{equation}

	Since $\norm{\fn - f}_{L^1} \xrightarrowp 0$, pass once to a subsequence of
	$\{n_k\}$ along which $\norm{\hat{f}_{n_k} - f}_{L^1} \to 0$ almost surely.
	The contradiction hypothesis is inherited by every subsequence, so nothing
	is lost and we continue to write $\{n_k\}$. Let $\Omega_0$ be the
	probability-one event on which that convergence holds. Securing it here,
	rather than by a second extraction later, is what keeps the argument free of
	a subsequence that depends on the sample point.

	For the total variation term, $C_1 = \sup_n \E[|D\fn|(\Omega)] < \infty$,
	so Markov's inequality gives $\Prob(|D\fn|(\Omega) > M) \leq C_1/M$
	uniformly in $n$; fix $M$ with $C_1/M < \delta/8$. For the $L^1$ term,
	$\norm{\fn}_{L^1} \leq \norm{\fn - f}_{L^1} + \norm{f}_{L^1}$ with
	$\norm{f}_{L^1} < \infty$ because $f \in \Wonep \subset \Lone$, so there are
	$C_2 < \infty$ and $K_0$ with
	$\Prob(\norm{\hat{f}_{n_k}}_{L^1} > C_2) < \delta/8$ for $k \geq K_0$.
	Setting
	\begin{equation}
		E_k = \left\{|D\hat{f}_{n_k}|(\Omega) \leq M\right\} \cap
		\left\{\norm{\hat{f}_{n_k}}_{L^1} \leq C_2\right\} ,
	\end{equation}
	we have $\Prob(E_k) \geq 1 - \delta/4$ for $k \geq K_0$, and on $E_k$ the
	realized $BV$ norm satisfies
	$\norm{\hat{f}_{n_k}}_{BV} \leq R := C_2 + M$.

	\subsubsection*{Step 2: A Positive-Probability Event of Bounded Realizations}

	For $k \geq K_0$ we have $\Prob(A_k \cap E_k) \geq \delta - \delta/4 =
	3\delta/4$. Because $\Prob(\limsup_k B_k) \geq \limsup_k \Prob(B_k)$ for any
	sequence of events $B_k$ in a probability space, the event
	\begin{equation}
		G = \left(\limsup_{k \to \infty} \, A_k \cap E_k\right) \cap \Omega_0
	\end{equation}
	satisfies $\Prob(G) \geq 3\delta/4 > 0$. Fix $\omega \in G$. By
	construction $\omega$ lies in $A_k \cap E_k$ for infinitely many $k$; let
	$\{k_j\}$ enumerate those indices. Along them the realized functions
	$\hat{f}_{n_{k_j}}(\omega)$ are a deterministic sequence bounded in $BV$ by
	$R$, so the $BV$ compactness theorem
	\citep[Theorem~5.5]{evans2015measure} applies to it as a statement about
	functions, with no measurable selection required: there are a further
	subsequence $\{k_{j_i}\}$ and $g \in \BV$ with
	\begin{enumerate}
		\item $\hat{f}_{n_{k_{j_i}}}(\omega) \to g$ strongly in $\Lone$,
		\item $D\hat{f}_{n_{k_{j_i}}}(\omega) \xrightarrow{w*} Dg$ in $\cM$.
	\end{enumerate}
	This is the $BV$ analogue of Rellich--Kondrachov: one extraction delivers
	both the $L^1$ convergence of the functions and the weak-* convergence of
	their gradient measures.

	\subsubsection*{Step 3: Identify the Limit}

	Because $\omega \in \Omega_0$, $\norm{\hat{f}_{n_k}(\omega) - f}_{L^1} \to 0$
	along the whole index set and hence along $\{k_{j_i}\}$. Limits in $\Lone$
	are unique, so $g = f$ $\mathcal{L}^d$-almost everywhere on $\Omega$,
	$Dg = Df$ as Radon measures, and item~(2) of Step~2 gives
	\begin{equation}
		\int_\Omega \phi \cdot d\!\left(D\hat{f}_{n_{k_{j_i}}}(\omega)\right)
		\longrightarrow \int_\Omega \phi \cdot d(Df) .
	\end{equation}

	\subsubsection*{Step 4: Derive the Contradiction}

	But $\omega \in A_{k_j}$ for every $j$, so
	\begin{equation}
		\left|\int_\Omega \phi \cdot d\!\left(D\hat{f}_{n_{k_j}}(\omega)\right)
		- \int_\Omega \phi \cdot d(Df)\right| > \epsilon
	\end{equation}
	for every $j$, in particular along $\{k_{j_i}\}$, contradicting the
	convergence just established. As $\omega \in G$ was arbitrary and
	$\Prob(G) > 0$, the contradiction hypothesis is untenable.

	\subsubsection*{Conclusion of Part I}

	The test function $\phi \in C_0(\Omega)^d$ in the contradiction hypothesis
	was arbitrary, so for every such $\phi$ the integrals
	$\int_\Omega \phi \cdot d(D\fn)$ converge in probability to
	$\int_\Omega \phi \cdot d(Df)$, which is to say
	\begin{equation}
		D\fn \xrightarrow{w*} Df \quad \text{in } \cM
		\quad \text{in probability.}
	\end{equation}

	\subsection{Part II: Strict Convergence Upgrade}
	
	We now prove the additional claim: if $\E[|D\fn|(\Omega)] \to |Df|(\Omega)$
	then $D\fn \to Df$ strictly in $\cM$, meaning weak-* convergence holds and
	$|D\fn|(\Omega) \to |Df|(\Omega)$.
	
	\subsubsection*{Step 1: Lower Semicontinuity Gives One Direction}
	
	The total variation functional $\mu \mapsto |\mu|(\Omega)$ is lower
	semicontinuous with respect to weak-* convergence of Radon measures
	\citep[Theorem~5.2]{evans2015measure}. Since $D\fn \xrightarrow{w*} Df$
	in $\cM$ by Part~I:
	\begin{equation}
		|Df|(\Omega) \leq \liminf_{n \to \infty} |D\fn|(\Omega)
	\end{equation}
	This gives the lower bound on the total variation of the limit.
	
	\subsubsection*{Step 2: The Additional Condition Gives the Other Direction}

	The additional assumption constrains a sequence of expectations and does not
	by itself bound $\limsup_n |D\fn|(\Omega)$ pathwise; a sequence of random
	variables can have convergent means while fluctuating without limit. The two
	directions combine through Fatou's lemma instead. Write
	$Z_n = |D\fn|(\Omega) - |Df|(\Omega)$ and take any subsequence. By Part~I it
	has a further subsequence along which $D\fn \xrightarrow{w*} Df$ almost
	surely, and along that further subsequence lower semicontinuity gives
	$\liminf_n Z_n \geq 0$ almost surely. Hence the negative part satisfies
	$Z_n^- \to 0$ almost surely, and since $0 \leq Z_n^- \leq |Df|(\Omega)
	< \infty$ the bounded convergence theorem gives $\E[Z_n^-] \to 0$. The
	assumption gives $\E[Z_n] \to 0$, so
	$\E[Z_n^+] = \E[Z_n] + \E[Z_n^-] \to 0$, whence $Z_n^+ \to 0$ in $L^1$ and
	therefore in probability. Both parts vanish, so $Z_n \to 0$ in probability
	along that further subsequence; as the initial subsequence was arbitrary,
	\begin{equation}
		|D\fn|(\Omega) \longrightarrow |Df|(\Omega) \quad \text{in probability.}
	\end{equation}

	\subsubsection*{Step 3: Conclude Strict Convergence}

	Part~I gives $D\fn \xrightarrow{w*} Df$ in $\cM$ in probability and Step~2
	gives convergence of the total masses in probability. Together these are
	precisely strict convergence in $\cM$
	\citep[Definition~3.13]{ambrosio2000bv}. \qed

	\begin{remark}
		The parallel between Lemma~\ref{lem:sobolev_consistency} and
		Lemma~\ref{lem:bv_consistency} is close but not exact. In both cases
		uniform boundedness in a strong space gives compactness via a compact
		embedding or compactness theorem, the consistency condition identifies
		the limit, and the subsequence argument promotes deterministic
		convergence to convergence in probability. The first difference is in the
		function spaces and the topology of convergence --- strong $L^p$ in the
		Sobolev case, weak-* of measures in the BV case --- reflecting the
		fundamental difference in the inductive biases of smooth and piecewise
		constant estimators. The second is in the mode of the boundedness
		hypothesis, and it is a difference of form rather than of substance.
		Definition~\ref{def:wellbehaved_sobolev} asks for tightness of
		$\norm{\fn}_{W^{2,p}}$, and Definition~\ref{def:wellbehaved_bv} asks for
		a bound on $\E[|D\fn|(\Omega)]$, which gives tightness by Markov's
		inequality. Neither makes the sequence almost surely bounded by a
		constant, so in neither case may a compactness theorem be applied to
		$\{\fn\}$ directly: a subsequence extracted from a random sequence
		depends on the realization, and an almost sure statement along such a
		subsequence will not combine with a probability bound along a
		deterministic one. Both proofs therefore work on an event where the
		realized norms are bounded and apply compactness there --- the Sobolev
		argument through the deterministic modulus of
		Lemma~\ref{lem:sobolev_modulus}, the BV argument through an event of
		positive probability.
	\end{remark}
	
	\begin{remark}
		The strict convergence upgrade requires the additional assumption that
		$\E[|D\fn|(\Omega)]$ converge to $|Df|(\Omega)$. This is not implied by
		weak-* convergence alone. This is because weak-* convergence of measures does not
		in general imply convergence of total masses --- mass can escape to the
		boundary or concentrate in ways invisible to continuous test functions. The
		lower semicontinuity of total variation \citep[Theorem~5.2]{evans2015measure}
		gives the liminf inequality for free; pinning the limsup requires the
		additional assumption. In applications where only weak-* convergence is
		needed --- as in the bootstrap corollaries --- the additional assumption
		is not required.
	\end{remark}
	
	% ============================================================
	\section{Proof of Theorem~\ref{M-thm:pairing} (Pairing Identity)}
	\label{app:thm_pairing}
	% ============================================================
	
	\begin{proof}
		Write points as $x = (x_{-j}, x_j)$ and for fixed $x_{-j}$ let
		$\Omega^{x_{-j}} = \{ \tau \in \R : (x_{-j}, \tau) \in \Omega \}$ denote the
		one-dimensional section of the domain and
		$g^{x_{-j}}(\tau) = g(x_{-j}, \tau)$ the section of $g$. By the
		one-dimensional section theory of BV functions
		\citep[Section~3.11]{ambrosio2000bv}, for $\mathcal{H}^{d-1}$-almost every
		$x_{-j}$ the section $g^{x_{-j}}$ of the precise representative is a good
		representative of a function of bounded variation on $\Omega^{x_{-j}}$, its
		distributional derivative $D g^{x_{-j}}$ is the corresponding section of
		$D_j g$, and the total variations disintegrate:
		\begin{equation}
			\int_\Omega \phi\; d D_j g
			= \int \left( \int_{\Omega^{x_{-j}}} \phi(x_{-j}, \tau)\;
			d D g^{x_{-j}}(\tau) \right) d\mathcal{H}^{d-1}(x_{-j})
			\label{eq:disintegration}
		\end{equation}
		for every bounded Borel $\phi$. For a good representative of a BV function of
		one variable, the fundamental theorem of calculus holds in the form
		$g^{x_{-j}}(b) - g^{x_{-j}}(a) = D g^{x_{-j}}\bigl((a, b]\bigr)$ for all
		$a < b$ in $\Omega^{x_{-j}}$ that are not atoms of $D g^{x_{-j}}$; since
		each section has at most countably many atoms, for
		Lebesgue-almost every $x_j$ the pair $(x_j - h, x_j + h)$ avoids them. Hence
		for almost every $x \in \operatorname{supp}(w)$ (where the whole segment lies
		in $\Omega$ by the trimming condition):
		\begin{equation}
			g(x + h e_j) - g(x - h e_j)
			= D g^{x_{-j}}\bigl((x_j - h,\; x_j + h]\bigr)
		\end{equation}
		Multiplying by $q(x)/2h$, integrating, and applying
		Tonelli on each section (justified since
		$\abs{D g^{x_{-j}}}$ is finite for a.e.\ section and $q$ is bounded):
		\begin{align}
			& \E\left[w(X)\, \Dh g(X)\right] \notag \\
			&= \int d\mathcal{H}^{d-1}(x_{-j}) \int dx_j\;
			\frac{q(x_{-j}, x_j)}{2h}
			\int \mathbf{1}\{\tau \in (x_j - h, x_j + h]\}\;
			d D g^{x_{-j}}(\tau) \notag \\
			&= \int d\mathcal{H}^{d-1}(x_{-j}) \int d D g^{x_{-j}}(\tau)\;
			\frac{1}{2h} \int \mathbf{1}\{x_j \in [\tau - h, \tau + h)\}\,
			q(x_{-j}, x_j)\; dx_j \notag \\
			&= \int d\mathcal{H}^{d-1}(x_{-j}) \int (K_h \star q)(x_{-j}, \tau)\;
			d D g^{x_{-j}}(\tau) \notag \\
			&= \int_\Omega (K_h \star q)\; d D_j g
		\end{align}
		using $\tau \in (x_j - h, x_j + h] \Leftrightarrow
		x_j \in [\tau - h, \tau + h)$ in the second line, the definition of
		$K_h \star q$ (half-open versus closed windows agree up to Lebesgue-null
		sets) in the third, and \eqref{eq:disintegration} in the last step. The
		piecewise constant case follows by substituting the explicit jump formula of
		Lemma~\ref{lem:bv_existence} for $D_j g$.
		
		For the corollary: if $q$ extends to a uniformly continuous function of
		compact support on $\R^d$, then $K_h \star q \to q$ uniformly as
		$h \downarrow 0$, so:
		\begin{equation}
			\abs{\int (K_h \star q)\; d D_j g - \int q\; d D_j g}
			\leq \norm{K_h \star q - q}_\infty \cdot \abs{D_j g}(\Omega) \to 0
		\end{equation}
		since $\abs{D_j g}(\Omega) < \infty$. When $g \in W^{1,1}(\Omega)$, the
		measure $D_j g$ is absolutely continuous with density $\partial_j g$, so
		$\int q\, d D_j g = \int w\, \partial_j g\, p\, dx
		= \E[w\, \partial_j g(X)]$.
	\end{proof}
	
	% ============================================================
	\section{Proof of Theorem~\ref{thm:neural_nets} (Neural Networks)}
	\label{app:neural_nets}
	% ============================================================
	
	\subsection{Setup and Notation}
	
	Let $\fn: \Omega \to \R$ be a feedforward neural network with $L$ layers,
	weight matrices $W_\ell^{(n)} \in \R^{d_\ell \times d_{\ell-1}}$, bias vectors
	$b_\ell^{(n)} \in \R^{d_\ell}$, and activation function $\sigma \in C^2(\R)$.
	The network computes:
	\begin{equation}
		\fn(x) = W_L^{(n)} z_{L-1} + b_L^{(n)}
	\end{equation}
	where the intermediate representations are defined recursively by:
	\begin{equation}
		z_\ell = \sigma\!\left(W_\ell^{(n)} z_{\ell-1} + b_\ell^{(n)}\right),
		\quad z_0 = x
	\end{equation}
	with $\sigma$ applied elementwise. Denote the pre-activation at layer $\ell$
	by $a_\ell = W_\ell^{(n)} z_{\ell-1} + b_\ell^{(n)}$ and the diagonal
	matrix of activation derivatives by $\Sigma'_\ell = \operatorname{diag}
	(\sigma'(a_\ell^{(1)}), \ldots, \sigma'(a_\ell^{(d_\ell)}))$ and similarly
	$\Sigma''_\ell = \operatorname{diag}(\sigma''(a_\ell^{(1)}), \ldots,
	\sigma''(a_\ell^{(d_\ell)}))$.
	
	Define the product of weight matrices from layer $k$ to layer $\ell$ as:
	\begin{equation}
		M_{\ell:k} = W_\ell^{(n)} \Sigma'_{\ell-1} W_{\ell-1}^{(n)}
		\Sigma'_{\ell-2} \cdots W_{k+1}^{(n)} \Sigma'_k
		\quad \text{for } \ell > k
	\end{equation}
	with the convention $M_{\ell:\ell} = \Sigma'_\ell$. This notation captures
	the Jacobian of the composition of layers from $k$ to $\ell$.
	
	The proof proceeds in two steps. Step~1 establishes uniform $\Wtwop$
	boundedness. Step~2 applies Lemma~\ref{lem:sobolev_consistency} to conclude.
	
	\subsection{Step 1: $W^{2,p}$ Tightness}
	
	\textbf{Claim:} $\norm{\fn}_{\Wtwop} = O_p(1)$.
	
	\subsubsection*{Substep 1a: Bound on $\fn$ in $L^p$}
	
	Since $\norm{\fn - f}_{L^p} \xrightarrowp 0$ and $f \in \Wonep \subset \Lp$,
	the triangle inequality gives:
	\begin{equation}
		\norm{\fn}_{L^p} \leq \norm{\fn - f}_{L^p} + \norm{f}_{L^p}
	\end{equation}
	The first term converges to zero in probability and the second is a finite
	constant since $f \in \Lp$. Hence $\norm{\fn}_{L^p} = O_p(1)$, which is the
	form of boundedness the claim asserts; no appeal to a subsequence or to an
	auxiliary condition on the network output is needed, because tightness is
	preserved under addition and a term converging to zero in probability is
	tight.
	
	\subsubsection*{Substep 1b: First Derivative Bound}
	
	Differentiating $\fn$ with respect to the input $x$ via the chain rule
	through all $L$ layers, the gradient at $x \in \Omega$ is:
	\begin{equation}
		\nabla \fn(x) = \left(W_L^{(n)} M_{L-1:1}\right)^\top
	\end{equation}
	where $M_{L-1:1} = W_{L-1}^{(n)} \Sigma'_{L-2} W_{L-2}^{(n)} \cdots
	\Sigma'_1 W_1^{(n)}$ is the product of weight-activation Jacobians from
	layer 1 to layer $L-1$. More explicitly, for the $j$-th input coordinate:
	\begin{equation}
		\frac{\partial \fn}{\partial x_j}(x)
		= \left(W_L^{(n)}\right)_{:} \cdot
		\prod_{\ell=L-1}^{1} \left(W_\ell^{(n)} \Sigma'_{\ell-1}\right) \cdot e_j
	\end{equation}
	where $e_j$ is the $j$-th standard basis vector.
	
	Taking the operator norm of this product and applying submultiplicativity:
	\begin{equation}
		\abs{\frac{\partial \fn}{\partial x_j}(x)}
		\leq \norm{W_L^{(n)}} \cdot \prod_{\ell=1}^{L-1}
		\norm{W_\ell^{(n)}} \cdot \norm{\Sigma'_\ell}
		\leq \prod_{\ell=1}^{L} \norm{W_\ell^{(n)}} \cdot \norm{\sigma'}_\infty^{L-1}
	\end{equation}
	using that $\norm{\Sigma'_\ell} = \max_i |\sigma'(a_\ell^{(i)})| \leq
	\norm{\sigma'}_\infty < \infty$ since $\sigma \in C^2$. This bound holds
	pointwise for every $x \in \Omega$, so integrating over $\Omega$:
	\begin{equation}
		\norm{\nabla \fn}_{L^p}
		\leq \norm{\nabla \fn}_{L^\infty} \cdot |\Omega|^{1/p}
		\leq \prod_\ell \norm{W_\ell^{(n)}} \cdot \norm{\sigma'}_\infty^{L-1}
		\cdot |\Omega|^{1/p}
	\end{equation}
	
	\subsubsection*{Substep 1c: Second Derivative Bound}
	
	Differentiating the gradient formula once more with respect to $x_i$ using
	the product rule and chain rule, the $(i,j)$ entry of the Hessian
	$\nabla^2 \fn(x)$ is:
	\begin{equation}
		\frac{\partial^2 \fn}{\partial x_i \partial x_j}(x)
		= \sum_{\ell=1}^{L-1}
		\left(W_L^{(n)} M_{L-1:\ell+1}\right)_{:}
		\cdot \Sigma''_\ell \cdot
		\left(M_{\ell:1} e_j\right) \cdot
		\left(M_{\ell:1} e_i\right)^\top \cdot
		\left(W_L^{(n)} M_{L-1:\ell+1}\right)_{:}^\top
	\end{equation}
	where the sum runs over all layers $\ell$ at which the second derivative of
	$\sigma$ is applied. The structure is: the forward pass from the input to
	layer $\ell$ contributes $M_{\ell:1} e_j$ for the $j$-th coordinate
	direction, the activation second derivative at layer $\ell$ contributes
	$\Sigma''_\ell$, and the backward pass from layer $\ell$ to the output
	contributes $W_L^{(n)} M_{L-1:\ell+1}$.
	
	Taking absolute values and applying submultiplicativity of operator norms:
	\begin{align}
		\abs{\frac{\partial^2 \fn}{\partial x_i \partial x_j}(x)}
		&\leq \sum_{\ell=1}^{L-1}
		\norm{W_L^{(n)} M_{L-1:\ell+1}}^2 \cdot
		\norm{\Sigma''_\ell} \cdot
		\norm{M_{\ell:1}}^2 \notag \\
		&\leq \sum_{\ell=1}^{L-1}
		\left(\prod_{k=\ell+1}^{L} \norm{W_k^{(n)}}\right)^2 \cdot
		\norm{\sigma''}_\infty \cdot
		\left(\prod_{k=1}^{\ell} \norm{W_k^{(n)}}\right)^2 \cdot
		\norm{\sigma'}_\infty^{2(L-2)} \notag \\
		&\leq (L-1) \cdot \left(\prod_{\ell=1}^{L} \norm{W_\ell^{(n)}}\right)^2
		\cdot \norm{\sigma''}_\infty \cdot \norm{\sigma'}_\infty^{2(L-2)}
	\end{align}
	using $\norm{\Sigma''_\ell} \leq \norm{\sigma''}_\infty < \infty$ since
	$\sigma \in C^2$, and collecting all weight norm factors into a single
	product. This bound holds pointwise for every $x \in \Omega$, so:
	\begin{equation}
		\norm{\nabla^2 \fn}_{L^p}
		\leq \norm{\nabla^2 \fn}_{L^\infty} \cdot |\Omega|^{1/p}
		\leq (L-1) \cdot \left(\prod_\ell \norm{W_\ell^{(n)}}\right)^2
		\cdot \norm{\sigma''}_\infty \cdot \norm{\sigma'}_\infty^{2(L-2)}
		\cdot |\Omega|^{1/p}
	\end{equation}
	
	\subsubsection*{Combining Substeps 1a--1c}
	
	Summing the three bounds:
	\begin{equation}
		\norm{\fn}_{W^{2,p}}
		= \norm{\fn}_{L^p} + \norm{\nabla \fn}_{L^p} + \norm{\nabla^2 \fn}_{L^p}
		\leq C\!\left(\prod_\ell \norm{W_\ell^{(n)}},\,
		\norm{\sigma'}_\infty,\, \norm{\sigma''}_\infty,\, L,\, |\Omega|\right)
	\end{equation}
	where $C$ is continuous and increasing in its first argument and depends on
	the remaining quantities but not on $n$ or $x$. The weight product
	$\prod_\ell \norm{W_\ell^{(n)}}$ is a function of the fitted weights and so
	is random; by assumption it is $O_p(1)$ --- which formalizes the effect of
	$L^2$ regularization during training \citep{bartlett2017spectrally} --- and
	$\norm{\sigma'}_\infty$, $\norm{\sigma''}_\infty$ are finite constants since
	$\sigma \in C^2$. A continuous increasing function of a tight sequence is
	tight, so
	\begin{equation}
		\norm{\fn}_{W^{2,p}} = O_p(1) ,
	\end{equation}
	which is the first condition of Definition~\ref{def:wellbehaved_sobolev}.
	This completes Step~1.
	
	\subsection{Step 2: $L^p$ Consistency}
	
	The $L^p$ consistency condition $\norm{\fn - f}_{L^p} \xrightarrowp 0$ is
	the standard universal approximation result for smooth neural networks.
	\citet{hornik1990universal} establish that multilayer feedforward networks
	with smooth activation functions are dense in $W^{1,p}(\Omega)$ with respect
	to the $W^{1,p}$ norm, implying in particular density in $L^p(\Omega)$.
	Combined with standard empirical risk minimization results establishing that
	the training procedure converges to the best approximation in the function
	class, this gives $\norm{\fn - f}_{L^p} \xrightarrowp 0$.
	
	\subsection{Step 3: Applying Lemma~\ref{lem:sobolev_consistency}}
	
	By Step~1, $\norm{\fn}_{W^{2,p}} = O_p(1)$, the bound being a fixed multiple
	of the random quantity $\prod_\ell \norm{W_\ell^{(n)}}$, which is $O_p(1)$
	by hypothesis. By Step~2,
	$\norm{\fn - f}_{L^p} \xrightarrowp 0$. These are precisely the two
	conditions of the well-behaved definition in the Sobolev sense
	(Definition~\ref{def:wellbehaved_sobolev}). Lemma~\ref{lem:sobolev_consistency}
	then gives:
	\begin{equation}
		\norm{\nabla \fn - \nabla f}_{L^p} \xrightarrowp 0
	\end{equation}
	which is the conclusion of Theorem~\ref{thm:neural_nets}. 
	\qed
	
	\begin{remark}
		The weight norm assumption $\sup_n \prod_\ell \norm{W_\ell^{(n)}} < \infty$
		is the key regularity condition on the estimator. It formalizes the effect
		of $L^2$ weight decay regularization during training, which penalizes large
		weights and keeps the product of operator norms bounded. Under standard
		conditions on the regularization parameter and optimization procedure,
		this bound holds with high probability over the training randomness
		\citep{bartlett2017spectrally}. The condition is stated as an assumption
		of the theorem rather than derived from first principles because its
		verification depends on details of the training procedure that are outside
		the scope of the present analysis.
	\end{remark}
	
	\begin{remark}
		Theorem~\ref{thm:neural_nets} is stated for smooth activations $\sigma \in
		C^2$. For ReLU networks $\sigma \notin C^2$ and the Hessian calculation in
		Substep~1c does not apply directly. The first derivative exists almost
		everywhere with $\sigma' = \mathbf{1}_{(0,\infty)}$ bounded, so Substep~1b
		goes through unchanged. The distributional second derivative is $\sigma'' =
		\delta_0$, a Dirac mass at zero, placing $\fn$ in a mixed $W^{1,p} \cap BV$
		space rather than $\Wtwop$. A modified version of
		Lemma~\ref{lem:sobolev_consistency} applicable in this mixed space, or a
		restriction to input distributions placing zero mass on the kink set
		$\{x : a_\ell(x) = 0\}$, would be required to handle this case. We leave
		this extension for future work.
	\end{remark}
	
	% ============================================================
	\section{Proof of Theorem~\ref{thm:boosting} (Gradient Boosting)}
	\label{app:boosting}
	% ============================================================
	
	\subsection{Setup and Notation}
	
	By Proposition~1 of \citet{buhlmann2003boosting}, the $L_2$Boost estimator after
	$m$ iterations can be represented as:
	\begin{equation}
		\hat{F}_m = (I - (I-S)^{m+1})Y
	\end{equation}
	where $S: \mathbb{R}^n \to \mathbb{R}^n$ is the base learner operator. The residual
	at step $m$ is:
	\begin{equation}
		u_m = Y - \hat{F}_m = (I-S)^m Y
	\end{equation}
	so that each base learner fit is $\hat{f}_m = Su_m$ and the ensemble updates as
	$\hat{F}_m = \hat{F}_{m-1} + \hat{f}_m$.

	Throughout this section $\norm{\cdot}_{L^2}$ on $\R^n$ denotes the
	\emph{normalized} empirical norm $\norm{v}_{L^2}^2 = n^{-1}\sum_{i=1}^n v_i^2$,
	so that it is the $L^2(P_n)$ norm of the function $i \mapsto v_i$ and is
	directly comparable across sample sizes. This matters because the constant
	in \eqref{eq:geometric_decay} is $\norm{Y}_{L^2}$: under the normalized norm
	$\norm{Y}_{L^2}^2 = n^{-1}\sum_i y_i^2 \to \E[Y^2] < \infty$ almost surely
	by Assumption~\ref{M-ass:sampling} and the law of large numbers, so
	$\norm{Y}_{L^2} = O_p(1)$ uniformly in $n$. Under the unnormalized Euclidean
	norm it would instead grow like $n^{1/2}$ and the uniform bounds below
	would be vacuous.

	By Proposition~2 of \citet{buhlmann2003boosting}, for a symmetric linear learner $S$
	with eigenvalues $\{\lambda_k; k = 1,\ldots,n\}$, the eigenvalues of the boosting
	operator $\mathcal{B}_m = I - (I-S)^{m+1}$ are $\{1-(1-\lambda_k)^{m+1}; k =
	1,\ldots,n\}$.
	
	Under the eigenvalue condition $0 < \lambda_k \leq 1$ for all $k$
	\citep[Theorem~1]{buhlmann2003boosting}, let:
	\begin{equation}
		\rho = \max_k (1-\lambda_k) \in [0,1)
	\end{equation}
	Since $\lambda_k > 0$ for all $k$, we have $\rho < 1$. The bias decays
	exponentially fast: by Theorem~1(a) of \citet{buhlmann2003boosting}:
	\begin{equation}
		\norm{u_m}_{L^2}^2 = \norm{(I-S)^m Y}_{L^2}^2
		\leq \rho^{2m} \norm{Y}_{L^2}^2
	\end{equation}
	In particular:
	\begin{equation}
		\norm{u_m}_{L^2} \leq C\rho^m, \quad C = \norm{Y}_{L^2} = O_p(1)
		\label{eq:geometric_decay}
	\end{equation}
	This geometric decay is the engine of the argument below. Note that $S$ is a
	contraction in the normalized norm exactly when it is one in the
	unnormalized norm, the normalization being a positive scalar multiple, so
	the eigenvalue condition and $\rho$ are unaffected.
	
	
	\subsection{$W^{2,p}$ Tightness and the Conclusion}
	
	\subsubsection*{Step 1: $W^{2,p}$ Bound for a Single Base Learner}
	
	Each $\hat{f}_m = Su_m \in \Wtwop$ by assumption on the base learner class. For
	shallow neural net base learners this follows from Section~\ref{app:neural_nets}.
	For spline base learners it follows from standard spline approximation theory
	\citep{wahba1990splines}. We require a uniform bound:
	\begin{equation}
		\sup_m \norm{\hat{f}_m}_{W^{2,p}} \leq C' \norm{u_m}_{L^2} < \infty
	\end{equation}
	where the bound on $\norm{\hat{f}_m}_{W^{2,p}}$ in terms of $\norm{u_m}_{L^2}$
	follows from the boundedness of $S$ as an operator from $L^2$ to $W^{2,p}$.
	
	\subsubsection*{Step 2: Ensemble Bound via Geometric Decay}
	
	By linearity of the $W^{2,p}$ norm and the triangle inequality:
	\begin{equation}
		\norm{\hat{F}_m}_{W^{2,p}}
		\leq \sum_{j=1}^m \norm{\hat{f}_j}_{W^{2,p}}
		\leq C' \sum_{j=1}^m \norm{u_j}_{L^2}
		\leq C'C \sum_{j=1}^m \rho^j
		\leq \frac{C'C\rho}{1-\rho}
		< \infty
	\end{equation}
	using the geometric decay \eqref{eq:geometric_decay} and the fact that
	$\sum_{j=1}^\infty \rho^j = \rho/(1-\rho) < \infty$ for $\rho < 1$. The bound
	is uniform in $m$, and uniform in $n$ as well once the two suprema of
	\eqref{eq:boost_uniform} are used: $C' \leq \sup_n \norm{S_n}_{L^2 \to W^{2,p}}$
	and $\rho/(1-\rho) \leq \bar{\rho}/(1-\bar{\rho})$.

	What the bound is not is deterministic. Its remaining factor is
	$C = \norm{Y}_{L^2}$, a random quantity, so what Step~2 establishes is
	\begin{equation}
		\norm{\hat{F}_m}_{W^{2,p}}
		\leq \frac{\bar{\rho}}{1-\bar{\rho}}
		\left(\sup_n \norm{S_n}_{L^2 \to W^{2,p}}\right) \norm{Y}_{L^2}
		= O_p(1),
		\label{eq:boost_tight}
	\end{equation}
	since $\norm{Y}_{L^2} = O_p(1)$ as noted after \eqref{eq:geometric_decay}.
	This is tightness in $\Wtwop$, which is what
	Definition~\ref{def:wellbehaved_sobolev} requires; a bound by a single
	deterministic constant would be stronger than the geometric-decay argument
	can deliver, because nothing caps $\norm{Y}_{L^2}$ almost surely under
	Assumption~\ref{M-ass:sampling}.

	\subsubsection*{Step 3: Apply Lemma~\ref{lem:sobolev_consistency}}

	$L^p$ consistency is assumed --- in $L^p$ and not $L^1$, matching
	Definition~\ref{def:wellbehaved_sobolev}. Tightness in $\Wtwop$ is
	\eqref{eq:boost_tight}. Lemma~\ref{lem:sobolev_consistency} gives
	$\norm{\nabla \hat{F}_m - \nabla f}_{L^p} \xrightarrowp 0$. \qed
	
	\begin{remark}[Why tree base learners are excluded]
		\label{rem:boosting_trees}
		An earlier form of this theorem carried a second case for tree stump base
		learners. It is withdrawn, because the hypotheses cannot be met by a tree.
		The eigenvalue condition $0 < \lambda_k \leq 1$ requires $S$ to have full
		rank, but a tree is the orthogonal projection onto the piecewise constant
		functions on its partition, so its spectrum is $\{0,1\}$ with
		$n - (\#\text{leaves})$ zeros. Then $\rho = \max_k(1-\lambda_k) = 1$ and
		the geometric series driving both steps above diverges --- the degenerate
		situation already noted in Remark~\ref{rem:degenerate}. Shrinkage does not
		help, since $\nu S$ has spectrum $\{0,\nu\}$, and nor does an $L^2$ leaf
		penalty; the zero eigenvalues are unavoidable for any finite-rank base
		learner. A tree whose split is instead chosen from the current residual,
		as in gradient boosting as it is implemented, is not a fixed linear
		operator at all, so the representation
		$\hat{F}_m = (I - (I-S)^{m+1})Y$ does not hold for it and neither does the
		decay that the representation supplies.

		Gradient boosting with greedily grown trees is therefore outside this
		theorem. It remains inside Theorem~\ref{M-thm:main}, which consumes an
		$L^2(P)$ rate and nothing else, and inside the pairing identity of
		Theorem~\ref{M-thm:pairing}, a finite sum of piecewise constant trees
		being of bounded variation. What is not established for it is gradient
		consistency of the \emph{uncorrected} plug-in, and there is reason to
		think it does not hold: at the number of rounds minimizing its own risk,
		such an ensemble can carry substantially more total variation than the
		function it estimates, with the gap widening in $n$ on targets poorly
		matched to axis-aligned splits. That is the regularization bias of
		Section~\ref{M-sec:inference} appearing in the interpretation theory, and
		it is an argument for the orthogonal correction rather than an omission
		from it.
	\end{remark}

\begin{remark}
		The key step is the geometric decay $\norm{u_m}_{L^2} \leq C\rho^m$,
		which follows from the eigenvalue condition $0 < \lambda_k \leq 1$ via Proposition~1
		and Theorem~1 of \citet{buhlmann2003boosting}. This decay is much faster than any
		polynomial rate and implies that the geometric series $\sum_j \rho^j$ converges
		without any condition on step sizes. The eigenvalue condition is the operative
		regularity assumption of the theorem, replacing the step size conditions common
		in the boosting literature. It is also what confines the result to full-rank
		smoothers; see Remarks~\ref{rem:degenerate} and~\ref{rem:boosting_trees}.
	\end{remark}
	
	\begin{remark}
		\label{rem:uniformity}
		Condition \eqref{eq:boost_uniform} is a genuine restriction, not a
		technicality. The operator $S_n$ acts on $\R^n$, so its spectrum moves
		with the sample size, and a gap for each $n$ separately does not bound
		$\sup_n (1-\rho_n)^{-1}$, which is what the geometric series in Step~2
		produces. The same applies to the constant $C'$ of Step~1, which is the
		$L^2 \to W^{2,p}$ norm of $S_n$. Both enter the uniform $\Wtwop$ bound
		that Lemma~\ref{lem:sobolev_consistency} consumes, and that bound is a
		statement about the sequence in $n$.

		The condition fails whenever the smoother's spectrum accumulates at
		zero as the sample grows, which includes a smoothing spline at a fixed
		smoothing parameter: there $\lambda_k^{(n)} = (1+\gamma d_k)^{-1}$ with
		the penalty eigenvalues $d_k$ unbounded in $n$, so $\rho_n \to 1$. It
		holds for smoothers that are bounded away from singularity by
		construction, for instance $S_n = \epsilon I + (1-\epsilon)S_n^0$ with
		$\epsilon \in (0,1)$ fixed and $S_n^0$ symmetric with spectrum in
		$[0,1]$, for which $\rho_n \leq 1 - \epsilon$ uniformly. Verifying
		\eqref{eq:boost_uniform} is therefore part of applying the theorem.
	\end{remark}

	\begin{remark}
		\label{rem:degenerate}
		When $\lambda_k \in \{0,1\}$ for all $k$ --- as occurs for linear projection
		operators --- Corollary~1 of \citet{buhlmann2003boosting} shows that
		$\mathcal{B}_m \equiv S$ for all $m$, meaning boosting does not improve on the
		base learner. In this degenerate case $\rho = \max_k(1-\lambda_k)$ takes the value
		1 for any $k$ with $\lambda_k = 0$, and the geometric series does not converge.
		The strict positivity condition $\lambda_k > 0$ excludes this case and corresponds
		to the requirement that the base learner has positive but incomplete explanatory
		power --- the standard ``weak learner'' assumption in the boosting literature.
	\end{remark}
	
	% ============================================================
	\section{Proof of Theorem~\ref{thm:forests} (Mondrian Forests)}
	\label{app:forests}
	% ============================================================
	
	\subsection{Setup}
	
	Let $\fn = \frac{1}{B}\sum_{b=1}^B T_b$ where $T_b$ are individual Mondrian
	trees. Each $T_b$ is piecewise constant on a random partition $\mathcal{P}_b$
	of $\Omega = [0,1]^d$ generated by the Mondrian process $MP(\lambda_n, [0,1]^d)$
	independently of the data, with budget $\lambda_n \asymp n^{1/(d+2)}$
	\citep{mourtada2017mondrian}. The partition $\mathcal{P}_b$ consists of
	$K_{\lambda_n} + 1$ hyperrectangular cells, where $K_{\lambda_n}$ is the
	number of splits \citep{roy2009mondrian}.
	
	\subsection{Step 1: Each Tree is in BV}
	
	Each $T_b$ is piecewise constant on a finite partition and therefore in
	$\BV$ by Lemma~\ref{lem:bv_existence}. The forest average $\fn =
	\frac{1}{B}\sum_b T_b$ is a finite convex combination of $BV$ functions,
	hence $\fn \in \BV$ with:
	\begin{equation}
		|D\fn|(\Omega) \leq \frac{1}{B} \sum_{b=1}^B |DT_b|(\Omega)
	\end{equation}
	by convexity of the total variation functional
	\citep[Proposition~3.6]{ambrosio2000bv}.
	
	\subsection{Step 2: Reduction to a Single Tree}

	Step~1 is already the whole of the forest argument. By convexity it suffices
	to bound one tree:
	\begin{equation}
		\E\!\left[|D\fn|(\Omega)\right]
		\leq \frac{1}{B}\sum_{b=1}^B \E\!\left[|DT_b|(\Omega)\right]
		= \E\!\left[|DT|(\Omega)\right]
		\label{eq:forest_convexity}
	\end{equation}
	for a single Mondrian tree $T$, the trees being identically distributed. No
	cancellation between trees is needed, and consequently no condition on $B$:
	the bound below holds for a forest of any fixed size, including $B = 1$.

	Steps~3 to~6 establish $\sup_n \E[|DT|(\Omega)] < \infty$. Write
	$\bar{y}_l$ for the mean of the responses falling in leaf $R_l$, with
	$\bar{y}_l = 0$ when $R_l$ is empty, and $N_l$ for that count.

	\subsection{Step 3: The Jump Set is Carried by the Split Faces}

	Let $k$ index the internal nodes of the tree, let $C_k$ be the cell split at
	node $k$ along coordinate $j_k$ at threshold $t_k$, and write
	\begin{equation}
		F_k = \overline{C}_k \cap \{x : x_{j_k} = t_k\}
	\end{equation}
	for the split face. Lemma~\ref{lem:bv_existence} gives the distributional
	gradient of a piecewise constant function as a sum over adjacent leaf pairs;
	the following re-indexes that sum by internal nodes, which is the form the
	later steps use.

	\begin{lemma}[Node-face representation]
		\label{lem:nodeface}
		Let $T$ be piecewise constant on the leaves of a hierarchical
		axis-aligned partition whose cuts are strictly interior to their cells.
		Then
		\begin{equation}
			|DT|(\Omega)
			= \sum_k \int_{F_k} \abs{\Delta_k(u)}\, d\mathcal{H}^{d-1}(u),
			\label{eq:nodeface}
		\end{equation}
		where $\Delta_k(u)$ is the difference of the two leaf values meeting at
		$u$.
	\end{lemma}

	\begin{proof}
		Two facts suffice. \emph{First}, each leaf-pair interface lies in the
		face of that pair's lowest common ancestor. Let $R_l$ and $R_m$ be
		adjacent leaves with lowest common ancestor $k$. They lie in opposite
		children of $k$, so after relabelling every $x \in \overline{R}_l$ has
		$x_{j_k} \leq t_k$ and every $x \in \overline{R}_m$ has
		$x_{j_k} \geq t_k$. A point of $\partial R_l \cap \partial R_m$ lies in
		both closures, hence satisfies $x_{j_k} = t_k$, and lies in
		$\overline{C}_k$; so $\partial R_l \cap \partial R_m \subseteq F_k$.

		\emph{Second}, distinct faces meet in an $\mathcal{H}^{d-1}$-null set.
		Take internal nodes $k \neq k'$. If neither is an ancestor of the other
		they lie in opposite subtrees of their own lowest common ancestor, so
		$C_k \cap C_{k'} = \emptyset$. Otherwise let $k$ be a strict ancestor of
		$k'$, so that $C_{k'}$ lies in one child of $k$. If $j_{k'} \neq j_k$
		the two hyperplanes are transverse and meet in dimension $d-2$, which is
		$\mathcal{H}^{d-1}$-null. If $j_{k'} = j_k$ they are parallel, and
		distinct: $t_{k'}$ is strictly interior to the extent of $C_{k'}$ in
		that coordinate while $C_{k'}$ lies entirely on one side of $t_k$.

		The map sending an adjacent leaf pair to its lowest common ancestor
		therefore sorts the interfaces into the faces without repetition, and
		conversely each point in the relative interior of $F_k$ lies between the
		leaf of $C_k \cap \{x_{j_k} \le t_k\}$ and the leaf of
		$C_k \cap \{x_{j_k} > t_k\}$ containing it, whose lowest common ancestor
		is $k$. Regrouping the sum of Lemma~\ref{lem:bv_existence} gives
		\eqref{eq:nodeface}.
	\end{proof}

	The argument is deterministic: it uses no probability, no property of the
	Mondrian process, and holds in every dimension. The hypothesis that cuts are
	strictly interior holds almost surely here, cut positions being uniform on an
	open interval, and is what the second fact requires --- a cut on a cell
	boundary would produce an empty child and a face coinciding with an
	ancestor's.

	\subsection{Step 4: Three Facts About the Mondrian Partition}

	\begin{lemma}[Expected interface area]
		\label{lem:interface}
		For $\mathcal{P} \sim MP(\lambda, C)$ on a box $C$,
		\begin{equation}
			\E\!\left[\sum_k \mathcal{H}^{d-1}(F_k)\right] = d\,\lambda\,\mathrm{vol}(C).
		\end{equation}
	\end{lemma}

	\begin{proof}
		Write $A(C,\lambda)$ for the left side, $s$ for the side lengths of $C$,
		$a = \sum_j s_j$ and $V = \mathrm{vol}(C)$. A cut of $C$ in coordinate
		$j$ has face area $\prod_{m \neq j} s_m = V/s_j$ whatever the cut
		position, and the two children have volumes summing to $V$. Conditioning
		on the first cut,
		\begin{equation}
			A(C,\lambda) = \int_0^\lambda a e^{-at}
			\left[ \frac{dV}{a} + d(\lambda - t)V \right] dt,
		\end{equation}
		where the induction hypothesis $A(\cdot,\lambda') = d\lambda'\,\mathrm{vol}(\cdot)$
		has been used on the children and is independent of which coordinate is
		cut and where. Evaluating,
		\begin{equation}
			A(C,\lambda)
			= dV\,\frac{1-e^{-a\lambda}}{a}
			+ dVa\left[\frac{\lambda}{a} - \frac{1-e^{-a\lambda}}{a^2}\right]
			= d\lambda V,
		\end{equation}
		the two boundary terms cancelling exactly.
	\end{proof}

	The cancellation is why the answer is free of the cell's shape. The second
	fact is the cell law of \citet{mourtada2017mondrian}: for $MP(\lambda)$ on
	$\R^d$ the cell containing a fixed point is
	$\prod_j [-E_j^-, E_j^+]$ with the $2d$ variables i.i.d.\ $\mathrm{Exp}(\lambda)$,
	so each side is $\mathrm{Gamma}(2,\lambda)$ and the sides are independent
	across coordinates.

	Both of the remaining steps need inverse moments of a single side of the
	cell, and both need them for the \emph{truncated} side, since the cells of
	$MP(\lambda,[0,1]^d)$ are the whole-space cells intersected with the cube.
	The following records them once.

	\begin{lemma}[Inverse moments of a truncated side]
		\label{lem:invmom}
		Let $s = \min(E^-, U) + \min(E^+, 1-U)$ with $E^\pm$ i.i.d.\
		$\mathrm{Exp}(\lambda)$ and $U \sim \mathrm{Unif}(0,1)$, which by
		projectivity is the law of a side of the cell of $MP(\lambda,[0,1]^d)$
		containing a uniform point. Then for $0 < p < 2$,
		\begin{equation}
			\E\!\left[s^{-p}\right]
			\leq \lambda^p\,\Gamma(2-p)\left(1 + \frac{p}{\lambda}\right)
			+ e^{-\lambda} .
			\label{eq:invmom}
		\end{equation}
	\end{lemma}

	\begin{proof}
		Write $s^{-p} = \Gamma(p)^{-1}\int_0^\infty r^{p-1}e^{-rs}\,dr$ and use
		$\E[e^{-r\min(E,u)}] = [\lambda + re^{-(\lambda+r)u}]/(\lambda+r)$. Given
		$U$ the two sides are independent, so integrating $U \sim \mathrm{Unif}(0,1)$
		out before $r$,
		\begin{equation}
			\E\!\left[e^{-rs}\right]
			= \frac{\lambda^2
			+ 2\lambda r\,(1-e^{-(\lambda+r)})/(\lambda+r)
			+ r^2 e^{-(\lambda+r)}}{(\lambda+r)^2} .
		\end{equation}
		Three elementary integrals follow, using
		$\int_0^\infty v^{a-1}(1+v)^{-a-b}dv = B(a,b)$ after $r = \lambda v$.
		The first gives
		$\Gamma(p)^{-1}\lambda^2\int r^{p-1}(\lambda+r)^{-2}dr
		= \lambda^p\Gamma(2-p)$. The second, after discarding
		$1-e^{-(\lambda+r)} \leq 1$, gives
		$\Gamma(p)^{-1} 2\lambda\int r^{p}(\lambda+r)^{-3}dr
		= p\,\lambda^{p-1}\Gamma(2-p)$. The third is at most $e^{-\lambda}$
		because $r^2/(\lambda+r)^2 \leq 1$ and
		$\Gamma(p)^{-1}\int r^{p-1}e^{-r}dr = 1$.
	\end{proof}

	The untruncated value is $\lambda^p\Gamma(2-p)$, so truncation costs the
	factor $1 + p/\lambda$ together with an exponentially small additive term:
	an $O(1/\lambda)$-thick boundary layer, as one would expect. At $p=1$ this
	is the bound $\lambda + 1 + e^{-\lambda}$ used in Lemma~\ref{lem:aspect}
	below; at $p = 1/2$ and $p = 3/2$ it is what Lemma~\ref{lem:noise} needs. It
	is convenient to write \eqref{eq:invmom} as
	\begin{equation}
		\E\!\left[s^{-p}\right] \leq c_p(\lambda)\,\lambda^p,
		\qquad
		c_p(\lambda) = \Gamma(2-p)\left(1 + \frac{p}{\lambda}\right)
		+ \lambda^{-p}e^{-\lambda},
		\label{eq:cp}
	\end{equation}
	so that $c_p(\lambda)$ decreases to the untruncated constant $\Gamma(2-p)$
	and is bounded by $c_p(1)$ for $\lambda \geq 1$. Only boundedness in
	$\lambda$ is used below; the rates come from the powers of $\lambda$.

	\begin{lemma}[Aspect ratio]
		\label{lem:aspect}
		Let $l(X)$ be the cell of $MP(\lambda, [0,1]^d)$ containing a uniform
		point $X$, with side lengths $s_j$. Then
		\begin{equation}
			\E\!\left[\frac{\mathrm{diam}(l(X))}{s_1(l(X))}\right]
			\leq 2d - 1 + \frac{4(d-1)}{\lambda}.
		\end{equation}
	\end{lemma}

	\begin{proof}
		The Mondrian is projective, so $MP(\lambda,[0,1]^d)$ is $MP(\lambda,\R^d)$
		restricted to the cube and restriction intersects every cell with the
		sub-box; hence
		$s_j = \min(E_j^-, X_j) + \min(E_j^+, 1 - X_j)$. The coordinates of $X$
		are independent and the $E_j^\pm$ are independent across $j$, so $s_m$
		and $s_1$ are independent and
		$\E[s_m/s_1] = \E[s_m]\,\E[1/s_1]$. Truncation only shrinks, so
		$\E[s_m] \leq 2/\lambda$.

		For the inverse moment apply Lemma~\ref{lem:invmom} with $p = 1$:
		$\E[1/s_1] \leq \lambda + 1 + e^{-\lambda}$, against the value $\lambda$
		on $\R^d$, so
		$\E[s_m/s_1] \leq 2 + 2(1+e^{-\lambda})/\lambda \leq 2 + 4/\lambda$ for
		$\lambda \geq 1$. Summing over $m$ with
		$\mathrm{diam} \leq \sum_m s_m$ and the $m = 1$ term equal to one gives
		the claim.
	\end{proof}

	The boundary therefore costs an additive constant rather than a rate: an
	$O(1/\lambda)$-thick layer contributing $O(1)$ to an integral of size
	$\lambda$. Since $\lambda_n \to \infty$ the constant converges to $2d-1$.

	\subsection{Step 5: The Signal Term}

	Write $\bar{y}_l = \bar{f}_l + \varepsilon_l$, where
	$\bar{f}_l = N_l^{-1}\sum_{i \in R_l} f(\mathbf{x}_i)$ and
	$\varepsilon_l = N_l^{-1}\sum_{i \in R_l}(y_i - f(\mathbf{x}_i))$, both zero
	on empty leaves, and set
	$T^{\mathrm{s}} = \sum_{N_l \geq 1} \bar{f}_l \mathbf{1}_{R_l}$,
	$T^{\mathrm{n}} = \sum_{N_l \geq 1} \varepsilon_l \mathbf{1}_{R_l}$, so that
	$T = T^{\mathrm{s}} + T^{\mathrm{n}}$. By \eqref{eq:nodeface} and the
	triangle inequality,
	\begin{equation}
		|DT|(\Omega) \leq S_1 + S_0 + N ,
	\end{equation}
	where $N$ collects the jumps of $T^{\mathrm{n}}$, $S_1$ the jumps of
	$T^{\mathrm{s}}$ across faces with both sides nonempty, and $S_0$ those
	across faces with an empty side.

	Fix the coordinate~$1$; the argument is identical for each. Along the line
	$\ell(x_{-1}) = \{(t, x_{-1}) : t \in [0,1]\}$ the tree is a step function
	whose jumps occur at the $1$-perpendicular cuts. The following records the
	only property of such a sum that is used, and in particular how a bound on a
	jump \emph{between two leaves} becomes a sum over single leaves.

	\begin{lemma}[Per-leaf charging]
		\label{lem:charging}
		Suppose each jump between consecutive leaves is bounded by
		$\delta(l') + \delta(l'')$ for the two leaves $l', l''$ meeting there and
		some nonnegative $\delta$ defined on leaves. Then
		\begin{equation}
			\sum_k \int_{F_k} \abs{\Delta_k}\, d\mathcal{H}^{d-1}
			\leq 2\sum_l \delta(l)\,\frac{\mathrm{vol}(l)}{s_1(l)}
			= 2\,\E_X\!\left[\frac{\delta(l(X))}{s_1(l(X))}\right],
			\label{eq:sections}
		\end{equation}
		the expectation being over a uniform $X$ with $l(X)$ its leaf.
	\end{lemma}

	\begin{proof}
		Let $l_0, \dots, l_{m}$ be the leaves $\ell$ crosses, in order, with
		values $v_0, \dots, v_m$. By hypothesis
		\begin{equation}
			\sum_{i=1}^{m} \abs{v_i - v_{i-1}}
			\leq \sum_{i=1}^{m}\left(\delta(l_{i-1}) + \delta(l_i)\right)
			\leq 2\sum_{i=0}^{m} \delta(l_i),
		\end{equation}
		because each index appears in at most two summands --- once for the
		neighbour on its left and once for the neighbour on its right. The
		neighbouring leaf's contribution is therefore absorbed, and what remains
		is a sum over single leaves. Integrating over $x_{-1}$ and using that
		leaf $l$ is crossed by a set of lines of measure
		$\mathrm{vol}(l)/s_1(l)$ gives the first form; the second is
		$\sum_l g(l)\mathrm{vol}(l) = \E_X[g(l(X))]$.
	\end{proof}

	It is \eqref{eq:sections}, with a per-leaf $\delta$, that the next three
	steps apply: $\delta(l) = L\,\mathrm{diam}(R_l)$ for the signal,
	$\delta(l) = \norm{f}_\infty \mathbf{1}\{N_l = 0\}$ for the empty-leaf term,
	and $\delta(l) = \abs{\varepsilon_l}$ for the noise.
Let $u \in \partial R_l \cap \partial R_m$ be a point of a face at which both
	leaves are nonempty, and let $f$ be Lipschitz with constant $L$. Every
	$\mathbf{x}_i \in R_l$ lies within $\mathrm{diam}(R_l)$ of $u$, so
	$|\bar{f}_l - f(u)| \leq L\,\mathrm{diam}(R_l)$, and likewise for $m$;
	hence
	\begin{equation}
		\abs{\bar{f}_m - \bar{f}_l}
		\leq L\left(\mathrm{diam}(R_l) + \mathrm{diam}(R_m)\right).
		\label{eq:signaljump}
	\end{equation}
	The diameters here are those of the two \emph{leaves} meeting at $u$, not
	that of the ancestor cell $C_k$, which may be larger than either. The bound
	is pairwise, involving both leaves; it is Lemma~\ref{lem:charging}, applied
	with the per-leaf quantity $\delta(l) = L\,\mathrm{diam}(R_l)$, that turns
	it into a sum over single leaves, the neighbour's diameter being absorbed by
	the factor of two. Then by Lemma~\ref{lem:aspect},
	\begin{equation}
		\E[S_1] \leq 2L\,
		\E\!\left[\frac{\mathrm{diam}(l(X))}{s_1(l(X))}\right]
		\leq 2L\left(2d - 1 + \frac{4(d-1)}{\lambda}\right).
	\end{equation}
	No rate in $n$ or $\lambda$ enters: the discrete variation of a Lipschitz
	function over a refining partition is controlled by the partition's aspect
	ratio, not by how fine it is.

	The bound \eqref{eq:signaljump} requires both leaves to be nonempty, since
	$\bar{f}_l$ is set to zero on an empty leaf and then bears no relation to
	$f$ near $u$. Faces with an empty side are collected in $S_0$ and bounded
	separately in Step~5a; a jump there is at most $\norm{f}_\infty$, which is
	not small, so the term is controlled by the scarcity of empty leaves rather
	than by the size of the jump.

	\subsection{Step 5a: Empty Leaves}

	At a face with one side empty the jump of $T^{\mathrm{s}}$ is at most
	$\norm{f}_\infty$, and at a face with both sides empty it is zero. By
	\eqref{eq:sections},
	\begin{equation}
		S_0 \leq 2\norm{f}_\infty
		\sum_{l\,:\,N_l = 0} \frac{\mathrm{vol}(l)}{s_1(l)} .
	\end{equation}
	Conditionally on the partition $\Prob(N_l = 0) = (1-\mathrm{vol}(l))^n
	\leq e^{-n\,\mathrm{vol}(l)}$, so with $V = \mathrm{vol}(l(X))$ as in
	Lemma~\ref{lem:noise},
	\begin{equation}
		\E[S_0] \leq 2\norm{f}_\infty\,
		\E\!\left[\frac{e^{-nV}}{s_1}\right]
		\leq 2\norm{f}_\infty n^{-1/2}\,
		\E\!\left[s_1^{-3/2}\right]\prod_{j \geq 2}\E\!\left[s_j^{-1/2}\right],
	\end{equation}
	using $e^{-x} \leq x^{-1/2}$, valid for $x > 0$ since
	$\sup_{x>0} x^{1/2}e^{-x} = (2e)^{-1/2} < 1$. The moments are those of
	Lemma~\ref{lem:invmom}, so by \eqref{eq:cp}
	\begin{equation}
		\E[S_0] \leq 2\norm{f}_\infty\, c_{3/2}\,c_{1/2}^{\,d-1}\,
		\lambda^{1+d/2} n^{-1/2},
	\end{equation}
	the same order as the noise term, and so $O(1)$ at
	$\lambda_n \asymp n^{1/(d+2)}$. Empty leaves are rare at this budget because
	the typical cell holds $n\,\mathrm{vol}(l) \asymp n^{2/(d+2)}$ observations,
	and that is what the bound exploits.

	\subsection{Step 6: The Noise Term}

	\begin{lemma}[Noise]
		\label{lem:noise}
		Under Assumption~\ref{M-ass:sampling},
		\begin{equation}
			\E[N] \leq 2\bar{\sigma} C_d(\lambda)\, \lambda^{1 + d/2} n^{-1/2},
			\qquad
			C_d(\lambda) = \sqrt{2}\,c_{3/2}\,c_{1/2}^{\,d-1}
			+ c_{3/2}^{\,(2+d)/3},
		\end{equation}
		with $c_p = c_p(\lambda)$ as in \eqref{eq:cp}. For $\lambda \geq 1$,
		$C_d(\lambda) \leq C_d(1) < \infty$, and $C_d(\lambda)$ decreases to
		$\sqrt{2\pi}(\sqrt{\pi}/2)^{d-1} + \pi^{(2+d)/6}$.
	\end{lemma}

	\begin{proof}
		Given the partition, $N_l \sim \mathrm{Bin}(n, \mathrm{vol}(l))$ and
		$\E[\abs{\varepsilon_l} \mid N_l] \leq \bar{\sigma} N_l^{-1/2}$ on
		$\{N_l \geq 1\}$. With $\mu = n\,\mathrm{vol}(l)$ and
		$\psi = \E[N^{-1/2}\mathbf{1}\{N \geq 1\}]$ we have
		$\psi \leq \mu$, since $N^{-1/2} \leq 1$ there, and also
		$\psi \leq \sqrt{2}\,\mu^{-1/2}$: indeed
		$N^{-1/2} \leq \sqrt{2}(N+1)^{-1/2}$ on $\{N \geq 1\}$, Jensen gives
		$\E[(N+1)^{-1/2}] \leq (\E[(N+1)^{-1}])^{1/2}$, and
		$\E[(N+1)^{-1}] = (1-(1-p)^{n+1})/((n+1)p) \leq 1/\mu$. So
		$\psi \leq \min(\sqrt{2}\mu^{-1/2}, \mu)$.

		Write $V = \mathrm{vol}(l(X)) = \prod_j s_j$, so that $\mu = nV$. The
		sides $s_j$ of the cell containing a uniform point are independent
		across $j$ and each has the truncated law of Lemma~\ref{lem:invmom}; we
		use that law throughout rather than the whole-space $\mathrm{Gamma}(2,\lambda)$,
		since both truncation effects --- shorter sides and hence larger inverse
		moments --- go against the bound. By \eqref{eq:sections} and
		$\sum_l g(l)\,\mathrm{vol}(l) = \E_X[g(l(X))]$ it suffices to show
		$\E[\psi/s_1] \leq C_d(\lambda)\,\lambda^{1+d/2} n^{-1/2}$.

		On $\{\mu \geq 1\}$ use $\psi \leq \sqrt{2}\mu^{-1/2}
		= \sqrt{2}\,n^{-1/2}V^{-1/2}$ and independence across coordinates:
		\begin{equation}
			\E\!\left[\frac{\psi}{s_1}; \mu \geq 1\right]
			\leq \sqrt{2}\,n^{-1/2}\,
			\E\!\left[s_1^{-3/2}\right]\prod_{j \geq 2}\E\!\left[s_j^{-1/2}\right]
			\leq \sqrt{2}\,c_{3/2}\,c_{1/2}^{\,d-1}\,\lambda^{1+d/2} n^{-1/2},
		\end{equation}
		by \eqref{eq:cp} twice, the exponents combining as
		$3/2 + (d-1)/2 = 1 + d/2$. The moment $\E[s_1^{-3/2}]$ is finite
		because the law of a side has density vanishing linearly at the origin,
		so that $t^{-3/2}\cdot t$ is integrable; $p = 2$ is the threshold in
		Lemma~\ref{lem:invmom} for exactly this reason. It is the last inverse
		moment that exists, and $3/2$ is precisely what the factor $1/s_1$ from
		\eqref{eq:sections} and the factor $\mu^{-1/2}$ from the binomial
		jointly require.

		On $\{\mu < 1\}$ use $\psi \leq \mu \leq 1$, then H\"{o}lder with
		exponents $3/2$ and $3$, then the Chernoff bound
		$\Prob(V < w) \leq w^s\,\E[V^{-s}] = w^s\prod_j \E[s_j^{-s}]
		\leq w^s (c_s \lambda^s)^d$ at $s = 3/2$:
		\begin{equation}
			\E\!\left[\frac{\psi}{s_1}; \mu < 1\right]
			\leq \left(\E\!\left[s_1^{-3/2}\right]\right)^{2/3}
			\Prob(V < 1/n)^{1/3}
			\leq c_{3/2}^{\,(2+d)/3}\,\lambda^{1+d/2} n^{-1/2}.
		\end{equation}
		Adding the two gives the claim.
	\end{proof}

	The small-cell regime, where $\psi$ behaves like $\mu$ rather than
	$\mu^{-1/2}$, needs no separate treatment: there the bound $\psi \leq \mu$ is
	the stronger of the two, and all that is required is that $\{V < 1/n\}$ carry
	little mass.

	\subsection{Step 7: Uniform BV Boundedness}

	Summing Steps~5, 5a and~6 over the $d$ coordinates and inserting
	$\lambda_n \asymp n^{1/(d+2)}$, for which
	$\lambda^{1+d/2}n^{-1/2} = O(1)$ because $(1 + d/2)/(d+2) = 1/2$ --- the
	constants in those three steps being bounded in $\lambda$ by
	\eqref{eq:cp} ---
	\begin{equation}
		\sup_n \E\!\left[|DT|(\Omega)\right] < \infty ,
	\end{equation}
	and by \eqref{eq:forest_convexity} the same bound holds for $\fn$ for every
	$B \geq 1$.

	\subsection{Step 8: Apply Lemma~\ref{lem:bv_consistency}}
	
	The $L^1$ consistency condition follows from the $L^2$ consistency result of
	\citet{mourtada2017mondrian} under the budget $\lambda_n \asymp n^{1/(d+2)}$
	by the Cauchy--Schwarz inequality on the bounded domain $\Omega$:
	\begin{equation}
		\norm{\fn - f}_{L^1}
		\leq |\Omega|^{1/2} \norm{\fn - f}_{L^2}
		\xrightarrowp 0
	\end{equation}
	Uniform BV boundedness is Step~7, for every $B \geq 1$.
	Lemma~\ref{lem:bv_consistency} then gives
	$D\fn \xrightarrow{w*} Df$ in $\cM$. \qed
	
	\begin{remark}
		Theorem~\ref{thm:forests} is stated for Mondrian forests, whose
		partitions are generated independently of the data. Independence enters
		in exactly two places, both in Step~4: the interface identity of
		Lemma~\ref{lem:interface} is a statement about the Mondrian process
		alone, and the cell law behind Lemma~\ref{lem:aspect} likewise. Neither
		survives data-dependent splitting. For Breiman forests the partition
		geometry and the response variation are coupled --- splits are chosen to
		maximize response differences across boundaries, which is the worst case
		for a bound on total variation --- so both lemmas would need replacing,
		not merely adapting. That is left as future work.

		What the argument no longer needs is any cancellation between trees.
		An earlier version bounded the second moment of $|D\fn|(\Omega)$ and
		relied on cross terms vanishing, which required $B$ to grow with $n$.
		The per-tree bound of Steps~3 to~7 removes both the cross-term argument
		and the condition on $B$, so the theorem now covers forests of the fixed
		size that is used in practice.
	\end{remark}
	

\bibliographystyle{plainnat}
\bibliography{references}

\end{document}
